% Électrolyse et synthèse de l'eau — PCT 3ème — Cours signé OnBuch+
% Programme officiel MINESEC. Compiler avec : tectonic N05-electrolyse-synthese-eau.tex
\newcommand{\DOCMATIERE}{PCT}
\newcommand{\DOCNIVEAU}{3ème}
\newcommand{\DOCMODULE}{Module 1 : Chimie}
\newcommand{\DOCLECON}{N05}
\newcommand{\DOCTITRE}{Électrolyse et synthèse de l'eau}
% OnBuch+ — préambule « cours signés » ENRICHI (compilé avec tectonic / XeLaTeX)
% Système visuel « Pop » d'OnBuch+ : fond crème (jamais blanc), bordures encre
% épaisses, ombres « dures » décalées, titres Archivo Black en capitales.
% Version MATHÉMATIQUES : graphes (pgfplots/tikz), boîtes pédagogiques
% (définition, propriété, méthode, attention, le savais-tu, exercice résolu,
% exercice, TP, bilan), page de garde et signature OnBuch+ ancrée en bas.
%
% Macros attendues dans main.tex AVANT \input{preamble} :
%   \DOCMATIERE  \DOCNIVEAU  \DOCMODULE  \DOCLECON (numéro)  \DOCTITRE
\documentclass[11pt]{article}

\usepackage{fontspec}
\usepackage[french]{babel}
\usepackage[a4paper,margin=2cm,top=3.4cm,bottom=2.8cm,headheight=1.4cm,footskip=1.3cm]{geometry}
\usepackage{amsmath,amssymb,amsfonts}
\usepackage{mathtools} % \xmapsto, \coloneqq... souvent utilisés par les modèles
\usepackage{cancel} % simplifications barrées (bilans de forces)
% \abs{z} (module, valeur absolue) et \norm{u} (norme), utilisés par les modèles
\providecommand{\abs}{}\let\abs\relax\DeclarePairedDelimiter{\abs}{\lvert}{\rvert}
\providecommand{\norm}{}\let\norm\relax\DeclarePairedDelimiter{\norm}{\lVert}{\rVert}
\usepackage[table]{xcolor} % [table] : fournit aussi \rowcolors (alternance de lignes)
\usepackage{pagecolor}
\usepackage{tikz}
\usetikzlibrary{arrows.meta,positioning,calc,shapes.geometric,shapes.misc,shapes.arrows,decorations.pathmorphing,decorations.pathreplacing,decorations.markings,patterns,shadows,mindmap,trees,fit,backgrounds,matrix,intersections,angles,quotes}
\usepackage{pgfplots}
\usepackage[version=4]{mhchem} % équations nucléaires \ce{^{235}_{92}U}
\usepackage[european,siunitx]{circuitikz} % schémas électriques
\pgfplotsset{compat=1.18}
\usepgfplotslibrary{fillbetween,groupplots} % aires entre courbes (calcul intégral)
\usepackage{enumitem}
\usepackage{fancyhdr}
% [most] charge toutes les bibliothèques tcolorbox (listings, minted, raster,
% documentation...) et gonfle inutilement la mémoire TeX sur un document long
% (~300 boîtes) : on ne charge que ce qui est réellement utilisé.
\usepackage[skins,breakable]{tcolorbox}
\usepackage{graphicx}
\usepackage{colortbl}
\usepackage{booktabs}
\usepackage{tabularx}
\usepackage{multirow}
\usepackage{diagbox} % case d'en-tête diagonale des tableaux à double entrée
% Colonnes extensibles centrée / gauche / droite (C, L, R), souvent utilisées par les modèles
\newcolumntype{C}{>{\centering\arraybackslash}X}
\newcolumntype{L}{>{\raggedright\arraybackslash}X}
\newcolumntype{R}{>{\raggedleft\arraybackslash}X}
\usepackage{array}
\usepackage{multicol}
\usepackage{siunitx}
% parse-numbers=false : accepte aussi \SI{2,5 \times 10^{-3}}{...}
\sisetup{locale=FR,output-decimal-marker={,},group-minimum-digits=4,parse-numbers=false}
\DeclareSIUnit\ohm{\ensuremath{\symup{\Omega}}} % U+2126 absent de la police
\DeclareSIUnit\division{div} % réglages d'oscilloscope (ms/div, V/div)
\DeclareSIUnit\tr{tr} % tours (vitesse de rotation, ex. \SI{1500}{\tr\per\minute})
\DeclareSIUnit\molar{M} % concentration molaire (ex. \SI{3}{\molar}, \SI{1}{\micro\molar})
\DeclareSIUnit\an{an} % année (le modèle confond parfois avec \year, un compteur TeX) — ex. \SI{5}{\kilo\meter\per\an}
\DeclareSIUnit\FCFA{FCFA} % franc CFA (ex. \SI{500}{\FCFA\per\kilo\gram})
\DeclareSIUnit\degreeBrix{\textdegree Brix} % teneur en sucre d'un sirop/jus (réfractométrie)
\DeclareSIUnit\month{mois} % le modèle utilise parfois \month, un compteur TeX, comme unité de durée
\DeclareSIUnit\microliter{\micro\liter} % le modèle écrit parfois \microliter en un seul token
\DeclareSIUnit\million{millions} % dénombrement (ex. \SI{15}{\million\per\mL} spermatozoïdes)
\usepackage{qrcode}
% \degree : commande fréquemment utilisée par les modèles pour un angle ou une
% température en dehors de \SI{}{} (non fournie par les paquets chargés).
\providecommand{\degree}{^{\circ}}
% \qed (fin de démonstration), utilisé par les modèles sans amsthm
\providecommand{\qed}{\hfill$\square$}
% \barre : double barre des valeurs interdites dans un tableau de variations (tkz-tab)
\providecommand{\barre}{\ensuremath{\|}}
% environnement proof (amsthm non chargé) utilisé par les modèles
\ifdefined\proof\else\newenvironment{proof}[1][Démonstration]{\par\noindent\textit{#1.}\ }{\qed\par}\fi
\usepackage{lastpage}
\usepackage{hyperref}
\hypersetup{hidelinks}

% ============================================================
% Palette « Pop » OnBuch+ — mêmes hex que la classe P (plus_ui.dart)
% ============================================================
\definecolor{poporange}{HTML}{F59321}
\definecolor{popdark}{HTML}{E07A0C}
\definecolor{popcream}{HTML}{FFF6E8}
\definecolor{popink}{HTML}{1C1714}
\definecolor{poppurple}{HTML}{7A5AE0}
\definecolor{popgreen}{HTML}{2E9E6A}
\definecolor{poppink}{HTML}{E0457B}
\definecolor{popblue}{HTML}{2F7DE1}
\definecolor{popgold}{HTML}{C98A12}
\definecolor{popmuted}{HTML}{8A7F76}
\definecolor{popline}{HTML}{EADFD2}
\definecolor{popgray}{HTML}{8A7F76}
\colorlet{popgrey}{popgray}
% Alias tolérants (le modèle nomme parfois les couleurs différemment)
\colorlet{popwhite}{white}
\colorlet{popblack}{popink}
\colorlet{popred}{poppink}
\colorlet{popyellow}{popgold}
% Teintes claires (fonds de tableaux/encadrés) — le modèle nomme parfois
% les couleurs avec un suffixe L (ex. popblueL) sans les avoir définies.
\colorlet{poporangeL}{poporange!20}
\colorlet{popdarkL}{popdark!20}
\colorlet{poppurpleL}{poppurple!20}
\colorlet{popgreenL}{popgreen!20}
\colorlet{poppinkL}{poppink!20}
\colorlet{popblueL}{popblue!20}
\colorlet{popgoldL}{popgold!20}
\colorlet{popredL}{popred!20}
\colorlet{popgrayL}{popgray!20}
% Teintes claires pour fonds de tableaux / zones
\colorlet{popblueL}{popblue!10!white}
\colorlet{poporangeL}{poporange!14!white}
\colorlet{popgreenL}{popgreen!12!white}
\colorlet{poppurpleL}{poppurple!12!white}
\colorlet{poppinkL}{poppink!12!white}
\colorlet{popgoldL}{popgold!14!white}

\pagecolor{popcream}

% ============================================================
% Typographie
% ============================================================
\defaultfontfeatures{Ligatures=TeX}
\setmainfont{PlusJakartaSans-Medium.ttf}[Path=../../../fonts/,BoldFont=PlusJakartaSans-Bold.ttf]
% Maths en sans-serif (Fira Math) avec chiffres et lettres droites de Plus
% Jakarta, pour que les formules se fondent dans le texte.
\usepackage[math-style=ISO,bold-style=ISO]{unicode-math}
\setmathfont{FiraMath-Regular.otf}[Path=../../../fonts/]
\setmathfont{PlusJakartaSans-Medium.ttf}[Path=../../../fonts/,range={up/num,up/{Latin,latin},\mathup}]
\usepackage{icomma} % virgule décimale sans espace en mode math
% Caractères mathématiques Unicode absents des polices de texte (Plus Jakarta,
% Archivo Black) : rendus par leur équivalent mathématique, en texte comme en
% formule — sinon ils s'affichent en carré vide (ex. « Arithmétique dans ℤ »).
\usepackage{newunicodechar}
\newunicodechar{ℝ}{\ensuremath{\mathbb{R}}}
\newunicodechar{ℤ}{\ensuremath{\mathbb{Z}}}
\newunicodechar{ℂ}{\ensuremath{\mathbb{C}}}
\newunicodechar{ℕ}{\ensuremath{\mathbb{N}}}
\newunicodechar{ℚ}{\ensuremath{\mathbb{Q}}}
\newunicodechar{𝔻}{\ensuremath{\mathbb{D}}}
\newunicodechar{α}{\ensuremath{\alpha}}
\newunicodechar{β}{\ensuremath{\beta}}
\newunicodechar{γ}{\ensuremath{\gamma}}
\newunicodechar{δ}{\ensuremath{\delta}}
\newunicodechar{ε}{\ensuremath{\varepsilon}}
\newunicodechar{θ}{\ensuremath{\theta}}
\newunicodechar{λ}{\ensuremath{\lambda}}
\newunicodechar{μ}{\ensuremath{\mu}}
\newunicodechar{σ}{\ensuremath{\sigma}}
\newunicodechar{τ}{\ensuremath{\tau}}
\newunicodechar{φ}{\ensuremath{\varphi}}
\newunicodechar{ψ}{\ensuremath{\psi}}
\newunicodechar{ω}{\ensuremath{\omega}}
\newunicodechar{Σ}{\ensuremath{\Sigma}}
% majuscules produites par \MakeUppercase dans les titres (α-aminés -> Α)
\newunicodechar{Α}{\ensuremath{\alpha}}
\newunicodechar{Β}{\ensuremath{\beta}}
\newunicodechar{Γ}{\ensuremath{\Gamma}}
\newunicodechar{Δ}{\ensuremath{\Delta}}
\newunicodechar{Ε}{\ensuremath{\varepsilon}}
\newunicodechar{Θ}{\ensuremath{\Theta}}
\newunicodechar{Λ}{\ensuremath{\Lambda}}
\newunicodechar{Μ}{\ensuremath{\mu}}
\newunicodechar{Τ}{\ensuremath{\tau}}
\newunicodechar{Φ}{\ensuremath{\Phi}}
\newunicodechar{Ψ}{\ensuremath{\Psi}}
\newunicodechar{Ω}{\ensuremath{\Omega}}
\newunicodechar{↦}{\ensuremath{\mapsto}}
\newunicodechar{∈}{\ensuremath{\in}}
\newunicodechar{∉}{\ensuremath{\notin}}
\newunicodechar{∧}{\ensuremath{\wedge}}
\newunicodechar{⇔}{\ensuremath{\Leftrightarrow}}
\newunicodechar{⟺}{\ensuremath{\Longleftrightarrow}}
\newunicodechar{⇒}{\ensuremath{\Rightarrow}}
\newunicodechar{⟹}{\ensuremath{\Longrightarrow}}
\newunicodechar{∘}{\ensuremath{\circ}}
\newunicodechar{∪}{\ensuremath{\cup}}
\newunicodechar{∩}{\ensuremath{\cap}}
\newunicodechar{≡}{\ensuremath{\equiv}}
\newunicodechar{⊂}{\ensuremath{\subset}}
\newunicodechar{⊥}{\ensuremath{\perp}}
\newunicodechar{∃}{\ensuremath{\exists}}
\newunicodechar{∀}{\ensuremath{\forall}}
\newunicodechar{∛}{\ensuremath{\sqrt[3]{\,}}}
\newunicodechar{∜}{\ensuremath{\sqrt[4]{\,}}}
\newunicodechar{⃗}{\ensuremath{{}^{\to}}}
\newunicodechar{ⁿ}{\ifmmode^{n}\else\textsuperscript{\ensuremath{n}}\fi}
\newunicodechar{ˣ}{\ifmmode^{x}\else\textsuperscript{\ensuremath{x}}\fi}
\newunicodechar{ᵇ}{\ifmmode^{b}\else\textsuperscript{\ensuremath{b}}\fi}
\newunicodechar{ʳ}{\ifmmode^{r}\else\textsuperscript{\ensuremath{r}}\fi}
\newunicodechar{ᵘ}{\ifmmode^{u}\else\textsuperscript{\ensuremath{u}}\fi}
\newunicodechar{ʸ}{\ifmmode^{y}\else\textsuperscript{\ensuremath{y}}\fi}
\newunicodechar{ᵃ}{\ifmmode^{a}\else\textsuperscript{\ensuremath{a}}\fi}
\newunicodechar{ᵉ}{\ifmmode^{e}\else\textsuperscript{\ensuremath{e}}\fi}
\newunicodechar{ᵏ}{\ifmmode^{k}\else\textsuperscript{\ensuremath{k}}\fi}
\newunicodechar{ᵖ}{\ifmmode^{p}\else\textsuperscript{\ensuremath{p}}\fi}
\newunicodechar{ᴺ}{\ifmmode^{N}\else\textsuperscript{\ensuremath{N}}\fi}
\newunicodechar{ᶜ}{\ifmmode^{c}\else\textsuperscript{\ensuremath{c}}\fi}
\newunicodechar{ʰ}{\ifmmode^{h}\else\textsuperscript{\ensuremath{h}}\fi}
\newunicodechar{ᵠ}{\ifmmode^{q}\else\textsuperscript{\ensuremath{q}}\fi}
\newunicodechar{ₙ}{\ifmmode_{n}\else\textsubscript{\ensuremath{n}}\fi}
\newunicodechar{ₐ}{\ifmmode_{a}\else\textsubscript{\ensuremath{a}}\fi}
\newunicodechar{ᵢ}{\ifmmode_{i}\else\textsubscript{\ensuremath{i}}\fi}
\newunicodechar{ᵣ}{\ifmmode_{r}\else\textsubscript{\ensuremath{r}}\fi}
\newunicodechar{ₖ}{\ifmmode_{k}\else\textsubscript{\ensuremath{k}}\fi}
\newunicodechar{ᵦ}{\ifmmode_{\beta}\else\textsubscript{\ensuremath{\beta}}\fi}
\newunicodechar{ₓ}{\ifmmode_{x}\else\textsubscript{\ensuremath{x}}\fi}
\newfontfamily\popdisplay{ArchivoBlack-Regular.ttf}[Path=../../../fonts/]
\newfontfamily\poplogo{SpaceGrotesk-Bold.ttf}[Path=../../../fonts/]
\newfontfamily\popsemi{PlusJakartaSans-SemiBold.ttf}[Path=../../../fonts/]
\newfontfamily\popxbold{PlusJakartaSans-ExtraBold.ttf}[Path=../../../fonts/]
\newfontfamily\popxboldls{PlusJakartaSans-ExtraBold.ttf}[Path=../../../fonts/,LetterSpace=3.0]

\setlist[enumerate]{leftmargin=1.7em,itemsep=5pt,topsep=4pt}
\setlist[itemize]{leftmargin=1.7em,itemsep=4pt,topsep=4pt,label={\color{poporange}\textbullet}}
\setlength{\parindent}{0pt}
\setlength{\parskip}{5pt}
\renewcommand{\arraystretch}{1.25}

% pgfplots : style Pop par défaut
\pgfplotsset{
  every axis/.append style={axis line style={popink,line width=1pt},tick style={popink},
    grid style={popline,line width=0.5pt},label style={font=\small},tick label style={font=\footnotesize}},
  popcurve/.style={poporange,line width=1.8pt,no markers,smooth},
}
% Même style disponible dans un \draw TikZ simple (hors pgfplots)
\tikzset{popcurve/.style={poporange,line width=1.8pt}}
\tikzset{popgrid/.style={popline,very thin}}
\colorlet{popcurve}{poporange} % le nom du style est parfois utilisé comme couleur

% ============================================================
% Pill / squiggle
% ============================================================
\newcommand{\poppill}[3]{%
  \tikz[baseline=-0.55ex]\node[draw=popink,line width=1.2pt,rounded corners=2.6pt,%
    fill=#2,inner xsep=6.5pt,inner ysep=3pt]{%
    \popxboldls\fontsize{7}{7}\selectfont\color{#3}\MakeUppercase{#1}};%
}
\newcommand{\popsquiggle}[1]{%
  \tikz[baseline=-0.4ex]{
    \draw[#1,line width=2.6pt,line cap=round]
      plot [smooth] coordinates {(0,0.09) (0.34,0.01) (0.68,0.21) (1.02,0.01) (1.36,0.10)};
  }%
}
% Mot-clé surligné (marqueur orange sous le texte)
\newcommand{\cle}[1]{{\popxbold\color{popdark}#1}}

% ============================================================
% En-tête / pied
% ============================================================
\pagestyle{fancy}
\fancyhf{}
\renewcommand{\headrulewidth}{0pt}
\renewcommand{\footrulewidth}{0pt}
\lhead{\raisebox{-0.15ex}{\poplogo\fontsize{13}{13}\selectfont\color{popink}OnBuch\color{poporange}+}}
\rhead{\poppill{\DOCMATIERE\ \textbullet\ \DOCNIVEAU\ \textbullet\ Leçon \DOCLECON}{white}{popink}}
\lfoot{\popxbold\fontsize{7.6}{7.6}\selectfont\color{popmuted}\MakeUppercase{Cours signé OnBuch+}\ \textbullet\ {\popsemi\DOCTITRE}}
\rfoot{\tikz[baseline=-0.6ex]\node[rounded corners=8.5pt,fill=popink,minimum height=17pt,minimum width=17pt,inner xsep=5pt,inner sep=0]{\popxbold\fontsize{8}{8}\selectfont\color{popcream}\thepage\,/\,\pageref*{LastPage}};}

% ============================================================
% Titres
% ============================================================
\newcommand{\chaptitre}[2]{%
  \begin{tcolorbox}[enhanced,colback=poporange,colframe=popink,boxrule=2.6pt,arc=11pt,
      drop shadow=popink,left=15pt,right=15pt,top=12pt,bottom=11pt,before skip=3pt,after skip=18pt]
    \poppill{#2}{popink}{popcream}\par\vspace{9pt}
    {\raggedright\hyphenpenalty=10000\exhyphenpenalty=10000\popdisplay\fontsize{22}{24}\selectfont\color{popcream}\MakeUppercase{#1}\par}
  \end{tcolorbox}
}
% Section numérotée : pastille orange avec le numéro + titre Archivo
\newcounter{popsec}
\newcommand{\coursec}[1]{%
  \stepcounter{popsec}\setcounter{popsub}{0}%
  \par\vspace{16pt}\noindent
  \begin{minipage}{\linewidth}
  \tikz[baseline=-0.7ex]\node[draw=popink,line width=1.6pt,fill=poporange,rounded corners=4pt,minimum size=22pt,inner sep=0]{\popdisplay\fontsize{12}{12}\selectfont\color{popcream}\thepopsec};%
  \hspace{8pt}{\popdisplay\fontsize{15}{17}\selectfont\color{popink}\MakeUppercase{#1}}\par
  \vspace{4pt}\noindent{\color{poporange}\rule{36pt}{3.6pt}}
  \end{minipage}\par\nopagebreak\vspace{8pt}
}
\newcounter{popsub}
\newcommand{\courssub}[1]{%
  \stepcounter{popsub}%
  \par\vspace{9pt}\noindent{\popxbold\fontsize{11.5}{13}\selectfont\color{popink}{\color{poporange}\thepopsec.\thepopsub}\hspace{6pt}#1}\par\nopagebreak\vspace{4pt}
}

% ============================================================
% Boîtes pédagogiques — même recette : fond blanc, bordure épaisse colorée,
% ombre dure encre, badge plein. Titre optionnel [#1].
% ============================================================
\tcbset{popbase/.style={breakable,enhanced,colback=white,boxrule=2.3pt,arc=9pt,
  shadow={3pt}{-3pt}{0pt}{popink},left=14pt,right=14pt,top=11pt,bottom=11pt,before skip=11pt,after skip=15pt,
  fontupper=\popsemi\fontsize{10.6}{14.5}\selectfont\color{popink},
  fontlower=\popsemi\fontsize{10.6}{14.5}\selectfont\color{popink}}}
% Coche dessinée (le glyphe ✓ n'existe pas dans les polices du document)
\newcommand{\popcheck}[1]{\tikz[baseline=-0.5ex]\draw[#1,line width=1.6pt,line cap=round,line join=round](-0.13,0.0)--(-0.04,-0.09)--(0.13,0.1);}
\providecommand{\coche}{\popcheck{popgreen}}
\providecommand{\resultat}[1]{\par\smallskip\noindent\fcolorbox{popgreen}{popgreenL}{\parbox{\dimexpr\linewidth-2\fboxsep-2\fboxrule\relax}{\textbf{Résultat.} #1}}\par\smallskip}
\newcommand{\popboxhead}[3]{\poppill{#1}{#2}{white}\ifx\relax#3\relax\else\hspace{6pt}{\popxbold\fontsize{10.6}{12}\selectfont\color{popink}#3}\fi\par\vspace{7pt}}

\newtcolorbox{aretenir}[1][]{popbase,colframe=popgreen,before upper={\popboxhead{À retenir}{popgreen}{#1}}}
\newtcolorbox{exemplebox}[1][]{popbase,colframe=poporange,before upper={\popboxhead{Exemple}{poporange}{#1}}}
\newtcolorbox{definition}[1][]{popbase,colframe=popblue,before upper={\popboxhead{Définition}{popblue}{#1}}}
\newtcolorbox{propriete}[1][]{popbase,colframe=poppurple,before upper={\popboxhead{Propriété}{poppurple}{#1}}}
\newtcolorbox{methode}[1][]{popbase,colframe=popgold,before upper={\popboxhead{Méthode}{popgold}{#1}}}
\newtcolorbox{attention}[1][]{popbase,colframe=poppink,before upper={\popboxhead{Attention}{poppink}{#1}}}
\newtcolorbox{experience}[1][]{popbase,colframe=popblue,colback=popblueL,before upper={\popboxhead{Expérience}{popblue}{#1}}}
% « Le savais-tu ? » : carte violette pleine (contexte, culture, Cameroun)
\newtcolorbox{savaistu}[1][]{popbase,colback=poppurple,colframe=popink,
  fontupper=\popsemi\fontsize{10.4}{14.2}\selectfont\color{white},
  before upper={\poppill{Le savais-tu\,?}{white}{poppurple}\ifx\relax#1\relax\else\hspace{6pt}{\popxbold\color{white}#1}\fi\par\vspace{7pt}}}
% Exercice résolu : énoncé en haut, \tcblower puis la solution sur fond crème
\newcounter{popexo}\newcounter{popexr}
\newtcolorbox{exoresolu}[1][]{popbase,colframe=poporange,colbacklower=popcream,
  segmentation style={popink,line width=1.2pt,dashed},
  before upper={\stepcounter{popexr}\popboxhead{Exercice résolu \thepopexr}{poporange}{#1}},
  before lower={\poppill{Solution}{popink}{popcream}\par\vspace{6pt}}}
% Exercice (non résolu en place) — la correction est dans la partie Corrigés
\newtcolorbox{exercice}[1][]{popbase,colframe=popink,
  before upper={\stepcounter{popexo}\popboxhead{Exercice \thepopexo}{popink}{#1}}}
% Corrigé
\newtcolorbox{corrige}[1][]{popbase,colframe=popgreen,colback=popgreenL,
  before upper={\popboxhead{Corrigé}{popgreen}{#1}}}
% Objectifs / prérequis / bilan
\newtcolorbox{objectifs}{popbase,colframe=popink,colback=poporangeL,
  before upper={\popboxhead{Objectifs de la leçon}{popink}{}}}
\newtcolorbox{prerequis}{popbase,colframe=popink,colback=popblueL,
  before upper={\popboxhead{Prérequis}{popblue}{}}}
\newtcolorbox{bilan}{popbase,colframe=popink,colback=white,boxrule=2.8pt,
  before upper={\popboxhead{Fiche bilan}{poporange}{L'essentiel à connaître pour l'examen}}}
% Figure encadrée avec légende
\newtcolorbox{popfigure}[1][]{enhanced,breakable=false,colback=white,colframe=popline,boxrule=1.4pt,arc=8pt,
  left=10pt,right=10pt,top=10pt,bottom=8pt,before skip=10pt,after skip=12pt,halign=center,
  lower separated=false,halign lower=center,
  fontlower=\popsemi\fontsize{9}{11.5}\selectfont\color{popmuted},
  #1}
\newcommand{\legende}[1]{\par\vspace{6pt}{\popsemi\fontsize{9}{11.5}\selectfont\color{popmuted}#1}}

% ============================================================
% Signature OnBuch+
% ============================================================
\newcommand{\poplogoinline}[1]{{\poplogo\fontsize{#1}{#1}\selectfont\color{popink}OnBuch\color{poporange}+}}
\newcommand{\signatureonbuch}{%
  \begin{tcolorbox}[enhanced,colback=popink,colframe=popink,boxrule=2.2pt,arc=10pt,
      drop shadow=poporange,left=14pt,right=14pt,top=10pt,bottom=10pt,before skip=0pt,after skip=0pt]
    \tikz[baseline=-0.7ex]\node[circle,fill=popgreen,draw=popcream,line width=1.3pt,minimum size=22pt,inner sep=0]{\popcheck{white}};%
    \hspace{9pt}\begin{minipage}[c]{0.62\linewidth}
      {\popxbold\fontsize{12}{13}\selectfont\color{popcream}Signé\ \textbullet\ {\poplogo OnBuch\color{poporange}+}}\par\vspace{2pt}
      {\raggedright\popsemi\fontsize{8}{10.5}\selectfont\color{popline}\MakeUppercase{Équipe pédagogique OnBuch+}\\\MakeUppercase{Conforme au programme officiel MINESEC}\par}
    \end{minipage}\hfill
    {\poplogo\fontsize{22}{22}\selectfont\color{popcream}OnBuch\color{poporange}+}
  \end{tcolorbox}%
}

% ============================================================
% Page de garde
% #1 titre · #2 matière · #3 niveau · #4 module · #5 n° leçon · #6 durée ·
% #7 description / objectifs (texte ou itemize)
% ============================================================
\newcommand{\pagegarde}[7]{%
  \thispagestyle{empty}
  \newgeometry{a4paper,margin=1.7cm,top=1.6cm,bottom=1.4cm}%
  \begin{tikzpicture}[remember picture,overlay]
    \fill[poporange!18] ([xshift=-3.2cm,yshift=-3.6cm]current page.north east) circle (4.6cm);
    \fill[poppurple!14] ([xshift=2.4cm,yshift=7.6cm]current page.south west) circle (3.8cm);
  \end{tikzpicture}%
  {\poplogo\fontsize{30}{30}\selectfont\color{popink}OnBuch\color{poporange}+}\hfill
  \raisebox{0.6ex}{\poppill{Cours signé \textbullet\ Premium}{popink}{popcream}}\par
  \vspace{1pt}\popsquiggle{poppurple}\par
  \vspace{8pt}
  \begin{tcolorbox}[enhanced,colback=poporange,colframe=popink,boxrule=2.8pt,arc=13pt,
      drop shadow=popink,left=18pt,right=18pt,top=12pt,bottom=13pt,after skip=10pt]
    \poppill{#2 \textbullet\ #3}{popink}{popcream}\hspace{5pt}\poppill{#4}{white}{popink}\par\vspace{8pt}
    {\popdisplay\fontsize{14}{14}\selectfont\color{popink}LEÇON #5}\par\vspace{4pt}
    {\raggedright\hyphenpenalty=10000\exhyphenpenalty=10000\popdisplay\fontsize{25}{27}\selectfont\color{popcream}\MakeUppercase{#1}\par}
    \vspace{8pt}
    {\popxbold\fontsize{9.5}{9.5}\selectfont\color{popink}\MakeUppercase{Durée conseillée : #6}}
  \end{tcolorbox}
  \begin{tcolorbox}[enhanced,colback=white,colframe=popink,boxrule=2.2pt,arc=10pt,
      drop shadow=popink,left=16pt,right=16pt,top=10pt,bottom=10pt,after skip=8pt]
    \poppill{Dans cette leçon}{poppurple}{white}\par\vspace{5pt}
    \popsemi\fontsize{9.3}{12.6}\selectfont\color{popink}#7
  \end{tcolorbox}
  \noindent\begin{minipage}[t]{0.487\linewidth}
    \begin{tcolorbox}[enhanced,colback=popgreen,colframe=popink,boxrule=2.2pt,arc=10pt,
        drop shadow=popink,left=11pt,right=11pt,top=10pt,bottom=10pt,height=3.1cm,valign=center]
      \tcbox[colback=white,colframe=popink,boxrule=1.3pt,arc=5pt,boxsep=0pt,nobeforeafter,
        left=3pt,right=3pt,top=3pt,bottom=3pt,valign=center]{\qrcode[height=1.9cm]{https://play.google.com/store/apps/details?id=com.luvvix.onbuch&hl=fr}}%
      \hspace{8pt}\begin{minipage}[c]{0.5\linewidth}
        {\popdisplay\fontsize{11}{12.5}\selectfont\color{white}\MakeUppercase{Télécharge OnBuch}}\par\vspace{4pt}
        {\raggedright\popsemi\fontsize{8.4}{11}\selectfont\color{white}Gratuit \textbullet\ Google Play\\Tuteur IA, annales, concours, résultats\par}
      \end{minipage}
    \end{tcolorbox}
  \end{minipage}\hfill
  \begin{minipage}[t]{0.487\linewidth}
    \begin{tcolorbox}[enhanced,colback=poppurple,colframe=popink,boxrule=2.2pt,arc=10pt,
        drop shadow=popink,left=11pt,right=11pt,top=10pt,bottom=10pt,height=3.1cm,valign=center]
      \tikz[baseline=-0.7ex]\node[circle,fill=white,draw=popink,line width=1.3pt,minimum size=1.6cm,inner sep=0]{\popdisplay\fontsize{24}{24}\selectfont\color{poppurple}+};%
      \hspace{8pt}\begin{minipage}[c]{0.6\linewidth}
        {\popdisplay\fontsize{11}{12.5}\selectfont\color{white}\MakeUppercase{Passe à OnBuch+}}\par\vspace{4pt}
        {\raggedright\popsemi\fontsize{8.4}{11}\selectfont\color{white}Tous les cours signés, Tuteur IA illimité, fiches PDF — onglet OnBuch+ de l'app\par}
      \end{minipage}
    \end{tcolorbox}
  \end{minipage}
  \vspace{8pt}
  \popcontenu
  \vfill
  \signatureonbuch
  \restoregeometry
  \newpage
}

% Rangée « Ce cours contient » (page de garde)
\newcommand{\poptile}[3]{% #1 couleur #2 gros texte #3 légende
  \begin{tcolorbox}[enhanced,colback=#1,colframe=popink,boxrule=2pt,arc=9pt,shadow={2.5pt}{-2.5pt}{0pt}{popink},
      left=6pt,right=6pt,top=9pt,bottom=9pt,halign=center,valign=center,height=1.75cm,width=\linewidth,nobeforeafter]
    {\popdisplay\fontsize{12.5}{13}\selectfont\color{white}\MakeUppercase{#2}}\par\vspace{3pt}
    {\centering\popsemi\fontsize{7.8}{9.5}\selectfont\color{white}#3\par}
  \end{tcolorbox}}
\newcommand{\popcontenu}{%
  \par\noindent{\popxboldls\fontsize{7.5}{7.5}\selectfont\color{popmuted}CE COURS SIGNÉ CONTIENT}\par\vspace{7pt}
  \noindent\begin{minipage}[t]{0.235\linewidth}\poptile{popblue}{Cours}{complet, illustré, pas à pas}\end{minipage}\hfill
  \begin{minipage}[t]{0.235\linewidth}\poptile{poporange}{Méthodes}{et exercices résolus}\end{minipage}\hfill
  \begin{minipage}[t]{0.235\linewidth}\poptile{poppink}{Exercices}{type examen + corrigés}\end{minipage}\hfill
  \begin{minipage}[t]{0.235\linewidth}\poptile{popgreen}{Fiche bilan}{l'essentiel à retenir}\end{minipage}\par}

% Sommaire
\newtcolorbox{sommaire}{enhanced,colback=white,colframe=popink,boxrule=2.2pt,arc=10pt,
  drop shadow=popink,left=18pt,right=18pt,top=13pt,bottom=13pt,before skip=6pt,after skip=20pt,
  before upper={\poppill{Sommaire}{poppurple}{white}\par\vspace{9pt}\popsemi\fontsize{10.6}{14.5}\selectfont\color{popink}}}

% Signature de fin de cours — poussée tout en bas de la dernière page
\newcommand{\findecours}{%
  \par\vfill\nopagebreak
  \begin{center}\popsquiggle{poporange}\end{center}\vspace{4pt}
  \signatureonbuch
}
\AtBeginDocument{\def\llbracket{\mathopen{[\mkern-3.5mu[}}\def\rrbracket{\mathclose{]\mkern-3.5mu]}}} % intervalles d'entiers (glyphe ⟦ absent de la police)
% Glyphes absents de Fira Math : redéfinis à partir de glyphes disponibles
\AtBeginDocument{%
  \def\setminus{\mathbin{\backslash}}\let\smallsetminus\setminus
  \def\vdots{\mathord{\vbox{\baselineskip3.2pt\lineskiplimit0pt\kern4pt\hbox{.}\hbox{.}\hbox{.}}}}%
  \def\ddots{\mathinner{\mkern1mu\raise6.5pt\hbox{.}\mkern2mu\raise3.6pt\hbox{.}\mkern2mu\raise0.7pt\hbox{.}\mkern1mu}}%
  \def\triangle{\mathord{\scalebox{0.9}{$\symup{\Delta}$}}}%
  \def\nabla{\mathord{\raisebox{\depth}{\scalebox{1}[-1]{$\symup{\Delta}$}}}}%
  \def\checkmark{\popcheck{popdark}}%
  \def\star{\mathbin{\ast}}\def\bigstar{\mathord{\ast}}%
  \def\diamond{\mathbin{\scalebox{0.6}{$\Diamond$}}}%
  \def\bigcup{\mathop{\mathchoice{\raisebox{-0.3em}{\scalebox{1.7}{$\cup$}}}{\scalebox{1.25}{$\cup$}}{\cup}{\cup}}\displaylimits}%
  \def\bigcap{\mathop{\mathchoice{\raisebox{-0.3em}{\scalebox{1.7}{$\cap$}}}{\scalebox{1.25}{$\cap$}}{\cap}{\cap}}\displaylimits}%
  \def\lesseqqgtr{\mathrel{\lessgtr}}\def\bigoplus{\mathop{\oplus}\displaylimits}%
}
\providecommand{\ssi}{si et seulement si} % abréviation fréquente
\AtBeginDocument{\providecommand{\wideparen}{\overparen}} % arc de cercle

%% FIGFIX-BEGIN v4 — correctifs globaux des figures (docs/onbuch-generation/tools/figfix)
\usepackage{adjustbox}\usepackage{etoolbox}
% toute figure TikZ/pgfplots plus large que la ligne est réduite à la largeur disponible
\BeforeBeginEnvironment{tikzpicture}{\begin{adjustbox}{max width=\linewidth}}
\AfterEndEnvironment{tikzpicture}{\end{adjustbox}}
% légendes pgfplots lisibles par-dessus les courbes
\pgfplotsset{every axis/.append style={legend style={fill=white,fill opacity=0.92,text opacity=1,draw=black!15,inner sep=3pt}}}
% étiquettes d'arêtes (syntaxe edge["..."]) avec fond blanc
\tikzset{every edge quotes/.append style={fill=white,inner sep=1.2pt}}
%% FIGFIX-END
\begin{document}

\pagegarde{Électrolyse et synthèse de l'eau}{PCT}{3ème}{Module 1 : Chimie}{N05}{1 séance (1 h chacune)}{Cette leçon étudie la décomposition de l'eau par le courant électrique (électrolyse) et sa reconstitution par combinaison des gaz obtenus (synthèse). Elle permet d'établir expérimentalement la nature et la formule de l'eau, et d'identifier les gaz hydrogène et dioxygène par leurs tests caractéristiques.
\vspace{4pt}
\begin{multicols}{2}\small
\begin{itemize}
\item À la fin de cette leçon, je sais définir l'électrolyse de l'eau et décrire le montage expérimental correspondant
\item À la fin de cette leçon, je sais identifier les gaz recueillis aux électrodes (H2 et O2) à l'aide des tests caractéristiques
\item À la fin de cette leçon, je sais interpréter les volumes de gaz recueillis pour établir la proportion 2 volumes d'hydrogène pour 1 volume de dioxygène
\item À la fin de cette leçon, je sais écrire et équilibrer l'équation-bilan de l'électrolyse de l'eau
\item À la fin de cette leçon, je sais décrire l'expérience de synthèse de l'eau (mélange détonant, eudiomètre) et interpréter ses résultats
\item À la fin de cette leçon, je sais déduire des expériences d'analyse et de synthèse la formule H2O de la molécule d'eau
\end{itemize}\end{multicols}}

\begin{sommaire}
\begin{multicols}{2}
\begin{enumerate}
\item Mise en évidence expérimentale : l'électrolyse de l'eau
\item Identification des gaz : tests du dihydrogène et du dioxygène
\item Interprétation de l'électrolyse : analyse de l'eau
\item La synthèse de l'eau
\item Formule et modèle de la molécule d'eau
\item Applications et sécurité
\item Activité d'intégration
\item Méthodes et exercices résolus
\item Exercices
\item Corrigés des exercices
\item Fiche bilan
\end{enumerate}
\end{multicols}
\end{sommaire}

\begin{objectifs}
\begin{itemize}
\item À la fin de cette leçon, je sais définir l'électrolyse de l'eau et décrire le montage expérimental correspondant
\item À la fin de cette leçon, je sais identifier les gaz recueillis aux électrodes (H2 et O2) à l'aide des tests caractéristiques
\item À la fin de cette leçon, je sais interpréter les volumes de gaz recueillis pour établir la proportion 2 volumes d'hydrogène pour 1 volume de dioxygène
\item À la fin de cette leçon, je sais écrire et équilibrer l'équation-bilan de l'électrolyse de l'eau
\item À la fin de cette leçon, je sais décrire l'expérience de synthèse de l'eau (mélange détonant, eudiomètre) et interpréter ses résultats
\item À la fin de cette leçon, je sais déduire des expériences d'analyse et de synthèse la formule H2O de la molécule d'eau
\item À la fin de cette leçon, je sais appliquer les consignes de sécurité liées à la manipulation du dihydrogène et du courant électrique
\item À la fin de cette leçon, je sais citer des applications concrètes de l'électrolyse et du dihydrogène dans la vie courante et l'industrie
\end{itemize}
\end{objectifs}

\begin{prerequis}
\begin{itemize}
\item Notion de molécule et d'atome (leçon sur la constitution de la matière, début d'année)
\item Écriture et équilibrage d'une équation-bilan simple
\item Distinction entre corps pur simple et corps pur composé
\item Notions d'électricité : courant continu, générateur, circuit électrique (partie Physique)
\item Propriétés générales des gaz (invisibilité, expansibilité)
\end{itemize}
\end{prerequis}

\newpage

\coursec{Mise en évidence expérimentale : l'électrolyse de l'eau}

À Yaoundé comme à Maroua, l'eau est partout : on la boit, on l'utilise pour cuisiner le ndolé ou le foléré, on la voit s'évaporer sous le soleil de midi. Pourtant, lors d'une panne d'électricité, un technicien d'ENEO explique aux élèves d'un collège que l'on peut \og casser \fg{} l'eau grâce au courant électrique pour obtenir deux gaz invisibles, dont l'un brûle avec une détonation ! Un liquide si familier, si calme, cacherait donc deux gaz ? Et comment prouver que l'eau est bien formée de ces deux gaz ?

C'est exactement la question que nous allons traiter dans cette leçon. Notre démarche sera celle du scientifique : nous allons d'abord \textbf{réaliser une expérience} (l'électrolyse de l'eau), observer attentivement ce qui se passe, identifier les produits obtenus, puis en tirer une conclusion sur la composition de l'eau. Dans cette première partie, nous nous concentrons sur l'expérience elle-même : le montage, les observations et le vocabulaire.

\courssub{Observation préalable : l'eau pure conduit-elle le courant ?}

\textbf{Une question qui surprend.} Tu as peut-être déjà entendu dire qu'il ne faut jamais toucher un appareil électrique avec les mains mouillées. On pourrait donc croire que l'eau conduit très bien le courant électrique. Vérifions cela par une expérience simple.

\begin{experience}[L'eau pure est-elle conductrice ?]
\textbf{Matériel :} un générateur de courant continu (pile de \SI{4,5}{V}), une lampe témoin, deux fils de connexion, un bécher d'eau distillée (eau pure), un bécher d'eau contenant quelques gouttes d'acide sulfurique dilué (ou de soude diluée), deux électrodes en graphite.

\textbf{Protocole :}
\begin{enumerate}
\item Plonger les deux électrodes dans le bécher d'eau distillée et fermer le circuit.
\item Observer la lampe.
\item Recommencer avec le bécher d'eau acidulée.
\end{enumerate}

\textbf{Observations :}
\begin{itemize}
\item Avec l'eau distillée, la lampe ne brille pratiquement pas : le courant qui traverse le circuit est extrêmement faible.
\item Avec l'eau acidulée, la lampe brille nettement : le courant passe.
\end{itemize}

\textbf{Interprétation :} l'eau pure est un très mauvais conducteur du courant électrique. En revanche, l'eau contenant un acide (ou une base comme la soude) conduit bien le courant, car elle contient des particules chargées (des ions) qui peuvent se déplacer et transporter le courant.
\end{experience}

\begin{aretenir}[Ce qu'il faut retenir]
L'\cle{eau pure} conduit très mal le courant électrique. Pour réaliser l'électrolyse de l'eau, on ajoute quelques gouttes d'\cle{acide sulfurique dilué} (ou de soude diluée) : on obtient de l'\cle{eau acidulée}, qui est conductrice. L'acide ajouté ne participe pas à la décomposition : il sert uniquement à rendre la solution conductrice.
\end{aretenir}

\begin{attention}[Erreur fréquente au BEPC]
Beaucoup d'élèves écrivent : \og l'eau conduit bien le courant électrique \fg{}. C'est faux ! C'est l'\textbf{eau acidulée} qui conduit le courant et c'est \textbf{l'eau} qui est décomposée. L'acide sulfurique n'est ni consommé, ni transformé : il joue seulement le rôle de \og transporteur \fg{} du courant. Dans une copie, précise toujours : \og on électrolyse de l'eau rendue conductrice par addition de quelques gouttes d'acide \fg{}.
\end{attention}

\courssub{Le voltamètre à eau : description du montage}

Pour décomposer l'eau par le courant électrique et recueillir les gaz formés, on utilise un appareil appelé \cle{voltamètre à eau} (ou électrolyseur à eau). Décrivons-le précisément.

\begin{definition}[Le voltamètre à eau]
Le \cle{voltamètre à eau} est un appareil destiné à réaliser l'électrolyse de l'eau. Il comprend :
\begin{itemize}
\item une \textbf{cuve} (en verre ou en plastique transparent) contenant l'eau acidulée ;
\item deux \textbf{électrodes} (tiges en graphite ou en platine) plongées dans le liquide et reliées à un \textbf{générateur de courant continu} ;
\item deux \textbf{tubes gradués renversés} (remplis d'eau acidulée au départ), placés au-dessus des électrodes pour recueillir les gaz dégagés.
\end{itemize}
\end{definition}

\begin{definition}[Vocabulaire de l'électrolyse]
\begin{itemize}
\item L'\cle{électrolyte} : la solution conductrice traversée par le courant (ici, l'eau acidulée).
\item Les \cle{électrodes} : les deux conducteurs solides plongés dans l'électrolyte et reliés au générateur.
\item L'\cle{anode} : l'électrode reliée à la borne \textbf{positive} ($+$) du générateur.
\item La \cle{cathode} : l'électrode reliée à la borne \textbf{négative} ($-$) du générateur.
\item L'\cle{électrolyseur} : l'ensemble cuve + électrolyte + électrodes, c'est-à-dire le récepteur dans lequel se déroule l'électrolyse.
\end{itemize}
\textbf{Moyen mnémotechnique :} retiens l'ordre alphabétique : \textbf{A}node va avec le pôle \textbf{+} et \textbf{C}athode avec le pôle \textbf{$-$}. On peut aussi retenir : \og \textbf{A}node = \textbf{A}limentée par le plus \fg{}.
\end{definition}

\begin{attention}[Sens du courant]
Par convention, à l'extérieur du générateur, le courant électrique sort de la borne $+$, traverse le circuit et rentre par la borne $-$. Dans le voltamètre, le courant entre donc par l'\textbf{anode} (borne $+$) et ressort par la \textbf{cathode} (borne $-$). Ne confonds jamais anode et cathode dans un schéma : c'est une faute classique qui coûte des points !
\end{attention}

\begin{popfigure}
\begin{center}
\begin{tikzpicture}[scale=0.95, font=\small]
% Générateur
\draw[fill=popgoldL, draw=popdark, rounded corners] (-0.9,-3.4) rectangle (0.9,-2.4);
\node at (0,-2.9) {\textbf{Générateur}};
\node at (0,-3.15) {\footnotesize courant continu};
\node[popdark] at (-1.15,-2.55) {$+$};
\node[popdark] at (1.15,-2.55) {$-$};
% Fils
\draw[very thick, popdark] (-0.6,-2.4) -- (-2.2,-2.4) -- (-2.2,0.4);
\draw[very thick, popdark] (0.6,-2.4) -- (2.2,-2.4) -- (2.2,0.4);
% Cuve
\draw[fill=popblueL, draw=popblue, thick] (-3.2,0.4) -- (-3.2,3.2) -- (3.2,3.2) -- (3.2,0.4) -- cycle;
\node[popblue!60!black] at (0,1.0) {eau acidulée (électrolyte)};
% Électrodes
\draw[fill=popdark] (-2.35,0.4) rectangle (-2.05,2.2);
\draw[fill=popdark] (2.05,0.4) rectangle (2.35,2.2);
\node[popdark, align=center] at (-2.2,-0.15) {\footnotesize anode (graphite)};
\node[popdark, align=center] at (2.2,-0.15) {\footnotesize cathode (graphite)};
% Tubes gradués renversés
\draw[draw=popink, thick, fill=white] (-2.75,1.3) rectangle (-1.65,4.6);
\draw[draw=popink, thick, fill=white] (1.65,1.3) rectangle (2.75,4.6);
% Gaz dans les tubes (anode: petit volume, cathode: double)
\draw[fill=poporangeL, draw=none] (-2.73,3.55) rectangle (-1.67,4.58);
\draw[fill=poporangeL, draw=none] (1.67,2.5) rectangle (2.73,4.58);
% Graduations
\foreach \y in {1.6,2.1,2.6,3.1,3.6,4.1} {
  \draw[popmuted] (-2.75,\y) -- (-2.6,\y);
  \draw[popmuted] (2.75,\y) -- (2.6,\y);
}
\node[poporange!70!black] at (-2.2,4.05) {\footnotesize $V$};
\node[poporange!70!black] at (2.2,3.5) {\footnotesize $2V$};
% Bulles
\foreach \p in {(-2.2,2.4),(-2.3,2.8),(-2.1,3.1)} {\draw[popblue, fill=white] \p circle (0.06);}
\foreach \p in {(2.2,2.3),(2.3,2.7),(2.1,3.0),(2.25,3.3),(2.15,1.9)} {\draw[popblue, fill=white] \p circle (0.06);}
% Flèches sens du courant
\draw[-{Stealth}, very thick, poppink] (-1.4,-2.4) -- (-1.8,-2.4);
\draw[-{Stealth}, very thick, poppink] (1.8,-2.4) -- (1.4,-2.4);
\node[poppink] at (0,-2.05) {\footnotesize sens conventionnel du courant $I$};
% Labels bornes
\node[popdark] at (-2.2,0.75) {\footnotesize $+$};
\node[popdark] at (2.2,0.75) {\footnotesize $-$};
\end{tikzpicture}
\end{center}
\legende{Schéma du voltamètre à eau. L'anode est reliée à la borne $+$ du générateur, la cathode à la borne $-$. Les gaz s'accumulent en haut des tubes gradués renversés : le volume recueilli à la cathode ($2V$) est le double de celui recueilli à l'anode ($V$).}
\end{popfigure}

\courssub{Protocole expérimental : réaliser l'électrolyse de l'eau}

\begin{methode}[Réaliser l'électrolyse de l'eau au laboratoire]
\begin{enumerate}
\item Verser de l'eau dans la cuve du voltamètre, puis ajouter \textbf{quelques gouttes} d'acide sulfurique dilué (avec précaution, l'enseignant manipule l'acide) pour rendre l'eau conductrice.
\item Remplir complètement les deux tubes gradués d'eau acidulée et les renverser au-dessus des électrodes, sans laisser entrer d'air.
\item Relier l'électrode de gauche à la borne $+$ du générateur (anode) et celle de droite à la borne $-$ (cathode).
\item Fermer l'interrupteur et régler le générateur (par exemple \SI{6}{V} à \SI{12}{V} en courant continu).
\item Observer les électrodes et laisser l'expérience durer plusieurs minutes.
\item Comparer les volumes de gaz recueillis dans les deux tubes, puis identifier chaque gaz (voir la partie 2 de la leçon).
\item Ouvrir l'interrupteur en fin d'expérience et ranger le matériel.
\end{enumerate}
\end{methode}

\begin{attention}[Consignes de sécurité]
\begin{itemize}
\item L'acide sulfurique, même dilué, est \textbf{corrosif} : porte des lunettes de protection, ne renverse pas la cuve, et en cas de contact avec la peau, rince abondamment à l'eau.
\item N'utilise \textbf{jamais} le courant du secteur (les prises ENEO de \SI{220}{V}) pour cette expérience : uniquement un générateur de \textbf{courant continu} de faible tension.
\item Les gaz produits sont inflammables ou favorisent la combustion : aucune flamme ne doit approcher les tubes avant le test contrôlé par l'enseignant.
\end{itemize}
\end{attention}

\courssub{Observations expérimentales}

Lorsque le circuit est fermé, voici ce que l'on observe, étape par étape :

\begin{itemize}
\item \textbf{Immédiatement}, de petites \textbf{bulles gazeuses} apparaissent à la surface des deux électrodes : il se dégage un gaz à l'anode et un gaz à la cathode.
\item Les bulles montent et s'accumulent en haut des tubes gradués renversés : le gaz \textbf{chasse le liquide}, dont le niveau descend dans chaque tube.
\item Le dégagement de bulles est \textbf{plus rapide à la cathode} qu'à l'anode.
\item Après quelques minutes, on constate que le tube placé au-dessus de la \textbf{cathode} contient \textbf{environ deux fois plus de gaz} que le tube placé au-dessus de l'anode.
\end{itemize}

\begin{exemplebox}[Un rapport de volumes constant : 2 pour 1]
Après 10 minutes d'électrolyse, un groupe d'élèves du collège de Ngaoundéré mesure :
\begin{itemize}
\item volume de gaz recueilli à la cathode : $V_{\text{cathode}} = \SI{20}{cm^3}$ ;
\item volume de gaz recueilli à l'anode : $V_{\text{anode}} = \SI{10}{cm^3}$.
\end{itemize}
Le rapport des volumes vaut :
\[ \frac{V_{\text{cathode}}}{V_{\text{anode}}} = \frac{20}{10} = 2. \]
Si l'on prolonge l'expérience et qu'on recueille \SI{30}{cm^3} à la cathode, on trouvera environ \SI{15}{cm^3} à l'anode. Quel que soit le temps d'électrolyse, le rapport reste \textbf{constant et égal à 2}. Ce n'est pas un hasard : il révèle la composition même de l'eau, comme nous le verrons dans la suite de la leçon.
\end{exemplebox}

\begin{popfigure}
\begin{center}
\begin{tikzpicture}[scale=0.9, font=\small]
% Tube anode
\draw[draw=popink, thick] (0,0) rectangle (1.4,5);
\draw[fill=popblueL] (0.02,0.02) rectangle (1.38,3.3);
\draw[fill=poporangeL] (0.02,3.3) rectangle (1.38,4.98);
\foreach \y in {0.5,1,1.5,2,2.5,3,3.5,4,4.5} {\draw[popmuted] (0,\y) -- (0.18,\y);}
\node[align=center] at (0.7,-0.5) {\footnotesize Tube de l'\textbf{anode} ($+$)};
\node[poporange!70!black] at (0.7,4.15) {\footnotesize gaz : $V$};
\node[popblue!60!black] at (0.7,1.6) {\footnotesize eau acidulée};
% Tube cathode
\draw[draw=popink, thick] (3.4,0) rectangle (4.8,5);
\draw[fill=popblueL] (3.42,0.02) rectangle (4.78,1.6);
\draw[fill=poporangeL] (3.42,1.6) rectangle (4.78,4.98);
\foreach \y in {0.5,1,1.5,2,2.5,3,3.5,4,4.5} {\draw[popmuted] (3.4,\y) -- (3.58,\y);}
\node[align=center] at (4.1,-0.5) {\footnotesize Tube de la \textbf{cathode} ($-$)};
\node[poporange!70!black] at (4.1,3.3) {\footnotesize gaz : $2V$};
\node[popblue!60!black] at (4.1,0.8) {\footnotesize eau acidulée};
% Double flèche comparaison
\draw[{Stealth}-{Stealth}, very thick, poppink] (1.7,3.3) -- (1.7,5);
\draw[{Stealth}-{Stealth}, very thick, poppink] (5.3,1.6) -- (5.3,5);
\node[poppink, rotate=90] at (1.45,4.15) {\footnotesize 1 volume};
\node[poppink, rotate=90] at (5.55,3.3) {\footnotesize 2 volumes};
\end{tikzpicture}
\end{center}
\legende{Comparaison des niveaux de gaz dans les deux tubes après quelques minutes d'électrolyse : la hauteur de gaz recueillie à la cathode est le double de celle recueillie à l'anode.}
\end{popfigure}

\begin{attention}[Pièges de rédaction]
\begin{itemize}
\item Ne dis pas \og le gaz monte parce qu'il est chaud \fg{} : les gaz montent dans les tubes parce qu'ils sont \textbf{moins denses que le liquide} et parce que les bulles formées aux électrodes sont poussées vers le haut.
\item Ne dis pas \og le liquide s'évapore \fg{} : le niveau du liquide descend dans les tubes parce que le gaz \textbf{le déplace} (il prend sa place).
\item Le rapport 2:1 concerne les \textbf{volumes} de gaz, pas les masses. Et c'est bien à la \textbf{cathode} (borne $-$) que le volume est le plus grand.
\end{itemize}
\end{attention}

\courssub{Conclusion : définition de l'électrolyse}

Que s'est-il passé ? Au départ, la cuve ne contenait que de l'eau (rendue conductrice par l'acide). À l'arrivée, deux gaz différents ont été recueillis, qui n'existaient pas avant l'expérience. Le passage du courant électrique a donc \textbf{décomposé} l'eau en deux corps simples gazeux. Nous identifierons ces gaz dans la partie suivante (test à la bûchette et détonation), mais nous pouvons déjà définir l'opération réalisée.

\begin{definition}[L'électrolyse]
L'\cle{électrolyse} est la \textbf{décomposition d'un corps composé sous l'action du courant électrique}.

L'électrolyse de l'eau est la décomposition de l'eau (corps composé) en deux corps simples gazeux : le \textbf{dihydrogène} (recueilli à la cathode, borne $-$) et le \textbf{dioxygène} (recueilli à l'anode, borne $+$), grâce au passage du courant électrique dans l'eau acidulée.
\end{definition}

\begin{aretenir}[Bilan de la première partie]
\begin{itemize}
\item L'eau pure conduit très mal le courant : on utilise de l'eau \textbf{acidulée}.
\item Le voltamètre à eau comprend une cuve, deux électrodes (anode au $+$, cathode au $-$) et deux tubes gradués renversés.
\item Il se dégage un gaz à chaque électrode ; le volume de gaz à la \textbf{cathode est le double} de celui recueilli à l'anode.
\item L'électrolyse est la décomposition d'un corps composé par le courant électrique.
\end{itemize}
\end{aretenir}

\begin{savaistu}[Une découverte vieille de plus de deux siècles]
L'électrolyse de l'eau a été réalisée pour la première fois en \textbf{1800} par deux savants anglais, \textbf{William Nicholson} et \textbf{Anthony Carlisle}, quelques jours seulement après l'invention de la pile électrique par l'Italien \textbf{Alessandro Volta}. Dès qu'une source de courant continu a existé, les savants se sont empressés de faire passer le courant dans l'eau... et ils ont découvert qu'elle se décomposait en deux gaz ! Aujourd'hui, cette même réaction est au cœur des technologies d'avenir : on l'utilise pour produire du \textbf{dihydrogène \og vert \fg{}}, un carburant propre qui pourrait un jour faire rouler des véhicules sans polluer, y compris au Cameroun.
\end{savaistu}

\begin{exoresolu}[Lecture d'une expérience d'électrolyse]
\textbf{Énoncé.} Un élève de 3ème réalise l'électrolyse de l'eau acidulée avec un voltamètre. Après 15 minutes, il lit sur le tube gradué placé au-dessus de l'électrode reliée à la borne $+$ un volume de gaz égal à \SI{12}{cm^3}.
\begin{enumerate}
\item Comment appelle-t-on l'électrode reliée à la borne $+$ ?
\item Quel volume de gaz l'élève lit-il approximativement dans l'autre tube ? Justifier.
\item Pourquoi a-t-on ajouté de l'acide à l'eau ?
\end{enumerate}
\tcblower
\textbf{Solution détaillée.}
\begin{enumerate}
\item L'électrode reliée à la borne $+$ du générateur s'appelle l'\textbf{anode}.
\item Le volume de gaz recueilli à la cathode est le double de celui recueilli à l'anode :
\[ V_{\text{cathode}} = 2 \times V_{\text{anode}} = 2 \times 12 = \SI{24}{cm^3}. \]
L'élève lit donc environ \SI{24}{cm^3} dans le tube de la cathode.
\item On a ajouté quelques gouttes d'acide parce que \textbf{l'eau pure conduit très mal le courant électrique}. L'acide rend la solution conductrice (elle devient un électrolyte) ; il ne participe pas à la décomposition, c'est bien l'eau qui est électrolysée.
\end{enumerate}
\end{exoresolu}

\begin{exercice}[Pour t'entraîner]
Lors d'une séance de travaux pratiques, un groupe ferme le circuit d'un voltamètre à eau et observe après quelques minutes \SI{8}{cm^3} de gaz dans le tube de l'anode.
\begin{enumerate}
\item Faire le schéma légendé du montage (cuve, électrodes, générateur, tubes, bornes $+$ et $-$).
\item Quel volume de gaz se trouve dans le tube de la cathode ?
\item Citer deux précautions de sécurité à respecter pendant cette expérience.
\end{enumerate}
\end{exercice}

\begin{corrige}[Exercice : pour t'entraîner]
\begin{enumerate}
\item Le schéma doit montrer : une cuve contenant l'eau acidulée, deux électrodes plongées dans le liquide, l'anode reliée à la borne $+$ du générateur de courant continu, la cathode reliée à la borne $-$, et deux tubes gradués renversés au-dessus des électrodes (voir la figure du cours).
\item Le volume à la cathode est le double : $V = 2 \times 8 = \SI{16}{cm^3}$.
\item Deux précautions possibles : porter des lunettes de protection à cause de l'acide ; utiliser un générateur de courant continu de faible tension et jamais le secteur \SI{220}{V} ; ne pas approcher de flamme des tubes remplis de gaz.
\end{enumerate}
\end{corrige}

\coursec{Identification des gaz : tests du dihydrogène et du dioxygène}

Après avoir mis en évidence la décomposition de l'eau par électrolyse (section précédente), nous avons recueilli deux gaz distincts : un volume important à la cathode (borne \cle{−}) et un volume deux fois plus faible à l'anode (borne \cle{+}). Comment être certain de leur nature ? En chimie, on ne devine pas : on \textbf{teste}. Chaque gaz possède une \emph{propriété chimique caractéristique} qui permet de l'identifier sans ambiguïté. C'est ce que nous allons découvrir maintenant, en suivant la démarche d'investigation : observation, hypothèse, expérience, interprétation, conclusion.

\courssub{Test du dihydrogène : la détonation caractéristique}

\begin{experience}[Test du dihydrogène (gaz de la cathode)]
\textbf{Matériel :} tube à essai rempli de gaz recueilli à la cathode (ou tube gradué renversé), support, bûchette en bois, allumette ou briquet, lunettes de protection.
\textbf{Protocole :}
\begin{enumerate}
    \item Maintenir le tube à l'envers (gaz moins dense que l'air, il reste dans le tube).
    \item Allumer la bûchette pour obtenir une flamme vive.
    \item Approcher \textbf{lentement} la flamme de l'ouverture du tube (sans la plonger dedans).
    \item Observer et écouter.
\end{enumerate}
\textbf{Observation :} On entend un \textbf{bruit sec, bref et fort}, comme un petit coup de pistolet ou un « aboiement ». Une flamme bleue fugace peut apparaître à l'orifice du tube. La bûchette continue de brûler normalement.
\textbf{Interprétation :} Le gaz contenu dans le tube a brûlé violemment au contact de la flamme et de l'oxygène de l'air. Cette combustion très rapide produit une onde de choc audible : c'est une \textbf{détonation}.
\textbf{Conclusion :} Ce gaz est le \textbf{dihydrogène} (\ce{H2}). C'est un gaz \cle{combustible}.
\end{experience}

\begin{popfigure}
\begin{tikzpicture}[scale=0.9, every node/.style={font=\small}]
    % Tube à essai gauche (H2)
    \draw[thick, popblue] (0,0) -- (0,4) -- (1,4) -- (1,3.5) arc (180:0:0.5) -- (2,0) -- cycle;
    \draw[thick, popblue] (0.5,0) ellipse (0.5 and 0.15);
    % Gaz H2 (particules fixes pour reproductibilité)
    \foreach \x/\y in {0.4/0.5, 1.2/0.7, 0.7/1.2, 1.5/1.0, 0.3/1.8, 1.6/1.5, 0.9/2.2, 1.3/2.5, 0.5/2.8, 1.4/3.0, 0.6/3.3, 1.1/3.6} {
        \fill[popblue!60] (\x,\y) circle (0.05);
    }
    % Flamme
    \draw[poporange, thick, decorate, decoration={snake, amplitude=1pt, segment length=3pt}] (1.5,3.8) -- (1.5,4.5);
    \fill[poporange!80!yellow] (1.5,4.5) ellipse (0.3 and 0.4);
    \node[poporange, anchor=south] at (1.5,4.9) {\textbf{Flamme}};
    % Onde de choc / Son
    \draw[popred, dashed, thick] (1.5,3.5) circle (1.2);
    \draw[popred, dashed, thick] (1.5,3.5) circle (1.6);
    \node[popred, anchor=north] at (1.5,1.5) {\textbf{« Boum ! » (Détonation)}};
    % Légende tube
    \node[align=center, anchor=north] at (1, -0.5) {Tube cathode ($-$)\\\ce{H2} (2 volumes)};
    
    % Tube à essai droit (O2)
    \begin{scope}[xshift=7cm]
        \draw[thick, popgreen] (0,0) -- (0,4) -- (1,4) -- (1,3.5) arc (180:0:0.5) -- (2,0) -- cycle;
        \draw[thick, popgreen] (0.5,0) ellipse (0.5 and 0.15);
        % Gaz O2 (moins de particules)
        \foreach \x/\y in {0.4/0.6, 1.2/0.9, 0.7/1.5, 1.5/1.3, 0.3/2.1, 1.6/1.8, 0.9/2.6, 1.3/3.0, 0.5/3.4} {
            \fill[popgreen!60] (\x,\y) circle (0.05);
        }
        % Bûchette incandescente (pointe rouge)
        \draw[popdark, thick] (2.5, 2) -- (4.5, 2);
        \fill[popred!80!black] (4.5, 2) circle (0.15);
        \draw[popred, thick, decorate, decoration={snake, amplitude=1.5pt, segment length=4pt}] (4.5, 2) -- (4.5, 2.8);
        \node[popred, anchor=west] at (4.7, 2.4) {\textbf{Rallumage vif}};
        % Bûchette avant (éteinte)
        \draw[popdark, thick, dashed] (2.5, 1) -- (4.5, 1);
        \fill[popmuted] (4.5, 1) circle (0.1);
        \node[popmuted, anchor=west] at (4.7, 1) {Bûchette incandescente};
        % Légende tube
        \node[align=center, anchor=north] at (1, -0.5) {Tube anode ($+$)\\\ce{O2} (1 volume)};
    \end{scope}
\end{tikzpicture}
\legende{À gauche : test du dihydrogène (cathode) → détonation à la flamme. À droite : test du dioxygène (anode) → rallumage d'une bûchette incandescente.}
\end{popfigure}

\begin{propriete}[Test du dihydrogène (\ce{H2})]
Le \textbf{dihydrogène} est un gaz \cle{combustible}. Au contact d'une flamme (source de chaleur) et en présence de dioxygène de l'air, il brûle violemment en produisant une \textbf{détonation} (bruit sec d'« aboiement »). L'équation de la combustion est :
\[ \ce{2 H2 + O2 -> 2 H2O} \quad \text{(réaction exothermique violente)} \]
Ce test est \textbf{spécifique} : aucun autre gaz courant ne donne ce bruit caractéristique.
\end{propriete}

\courssub{Test du dioxygène : le rallumage de la bûchette incandescente}

\begin{experience}[Test du dioxygène (gaz de l'anode)]
\textbf{Matériel :} tube à essai rempli de gaz recueilli à l'anode, support, bûchette en bois, allumette, lunettes de protection.
\textbf{Protocole :}
\begin{enumerate}
    \item Allumer la bûchette, puis souffler sur la flamme pour l'éteindre. La pointe doit rester \textbf{rouge incandescente} (braise) sans flamme visible.
    \item Plonger \textbf{délicatement} la pointe incandescente \textbf{à l'intérieur} du tube (sans toucher les parois humides).
    \item Observer.
\end{enumerate}
\textbf{Observation :} La bûchette \textbf{se rallume instantanément} avec une flamme vive et brillante. Elle brûle beaucoup plus fort que dans l'air ambiant.
\textbf{Interprétation :} Le gaz contenu dans le tube ne brûle pas lui-même (pas de détonation), mais il \textbf{entretient la combustion} de la bûchette (du carbone). Il fournit l'oxygène nécessaire à la réaction : \ce{C + O2 -> CO2}. La concentration en dioxygène dans le tube est bien supérieure à celle de l'air (\SI{21}{\%}).
\textbf{Conclusion :} Ce gaz est le \textbf{dioxygène} (\ce{O2}). C'est un gaz \cle{comburant}.
\end{experience}

\begin{propriete}[Test du dioxygène (\ce{O2})]
Le \textbf{dioxygène} est un gaz \cle{comburant} : il ne brûle pas, mais il \textbf{entretient la combustion}. En présence d'une source de chaleur (bûchette incandescente), il permet le rallumage vif de celle-ci. C'est le test de reconnaissance standard du dioxygène.
\end{propriete}

\begin{attention}[Erreur fréquente : inverser les tests ou les électrodes]
\begin{itemize}
    \item \textbf{Piège 1 :} Passer la flamme dans le tube de l'anode. \textbf{Rien ne se passe} (pas de détonation). On pourrait croire à tort qu'il n'y a pas de gaz.
    \item \textbf{Piège 2 :} Plonger la bûchette incandescente dans le tube de la cathode. \textbf{Rien ne se passe} (le \ce{H2} ne rallume pas la braise, il faut une flamme).
    \item \textbf{Piège 3 (majeur) :} Inverser les bornes du générateur. Le \ce{H2} se dégage alors à l'anode (borne \cle{+}) et le \ce{O2} à la cathode (borne \cle{−}). \textbf{Règle immuable :} Le dihydrogène se dégage \textbf{toujours à la cathode (borne \cle{−})}, car c'est là que les ions \ce{H+} gagnent des électrons (réduction). Le volume de \ce{H2} est \textbf{double} de celui de \ce{O2}.
\end{itemize}
\end{attention}

\courssub{Règle mnémotechnique et tableau de synthèse}

Pour ne jamais confondre les deux tests, retiens cette phrase simple :

\begin{aretenir}[Mémo gaz de l'électrolyse]
\begin{itemize}
    \item Le gaz qui \textbf{« aboie »} (détonation à la flamme) $\rightarrow$ \textbf{Dihydrogène \ce{H2}} (Cathode, borne \cle{−}, 2 volumes).
    \item Le gaz qui \textbf{« ravive »} (rallume la bûchette incandescente) $\rightarrow$ \textbf{Dioxygène \ce{O2}} (Anode, borne \cle{+}, 1 volume).
\end{itemize}
\end{aretenir}

\begin{center}
\begin{tabularx}{\linewidth}{|l|X|X|X|X|}\hline
\rowcolor{popblueL}
\textbf{Électrode} & \textbf{Gaz dégagé} & \textbf{Volume relatif} & \textbf{Test effectué} & \textbf{Observation / Résultat} \\ \hline
\textbf{Cathode} (borne \cle{−}) & \ce{H2} (Dihydrogène) & \textbf{2 volumes} & Approcher une \textbf{flamme} (bûchette allumée) & \textbf{Détonation} (bruit d'aboiement) \\ \hline
\textbf{Anode} (borne \cle{+}) & \ce{O2} (Dioxygène) & \textbf{1 volume} & Plonger une \textbf{bûchette incandescente} (point rouge) & \textbf{Rallumage vif} de la bûchette \\ \hline
\end{tabularx}
\end{center}

\courssub{Consignes de sécurité indispensables}

L'électrolyse de l'eau produit un mélange gazeux potentiellement dangereux. Au laboratoire comme à la maison (groupes électrogènes, batteries de voiture), la rigueur s'impose.

\begin{attention}[Sécurité : le mélange détonant (Knallgas)]
\begin{itemize}
    \item Le mélange \ce{H2} / \ce{O2} (ou \ce{H2} / air) dans de \textbf{grandes proportions} est \textbf{explosif}. Une simple étincelle, une flamme ou de l'électricité statique peut déclencher une violente explosion.
    \item \textbf{Ne jamais} chercher à enflammer le gaz directement à la sortie des tubes s'ils contiennent un gros volume.
    \item \textbf{Ne jamais} boucher les tubes pour « garder le gaz » afin de le faire exploser plus fort : c'est \textbf{interdit} et \textbf{mortel}.
    \item Travailler \textbf{petites quantités} (tubes à essai, seringues de \SI{10}{mL} ou \SI{20}{mL}).
    \item Toujours \textbf{éloigner le visage} des ouvertures des tubes lors des tests.
    \item Aérer le local (hotte ou fenêtres ouvertes) pour éviter l'accumulation de \ce{H2} au plafond (gaz plus léger que l'air).
\end{itemize}
\end{attention}

\begin{savaistu}[Le « Knallgas » et l'histoire de la chimie]
Le mélange détonant hydrogène-oxygène s'appelle \emph{Knallgas} (gaz qui claque) en allemand. C'est en étudiant cette réaction que \textbf{Henry Cavendish} (1766) a montré que l'eau n'est pas un élément, mais un « composé » d'« air inflammable » (H2) et d'« air vital » (O2). Plus tard, \textbf{Lavoisier} répète l'expérience, nomme les gaz et prouve la conservation de la masse : la masse d'eau décomposée égale la masse des gaz formés. Au Cameroun, ce principe est utilisé dans les \textbf{piles à combustible} (projets pilotes pour l'énergie propre) et dans la sécurité des \textbf{batteries plomb-acide} des groupes électrogènes et voitures : la surcharge produit du \ce{H2} et \ce{O2} ; les bouchons à évent évacuent ce mélange pour éviter l'explosion de la batterie.
\end{savaistu}

\courssub{Exercice résolu : identification à l'aveugle}

\begin{exoresolu}[Reconnaissance de deux gaz inconnus]
On dispose de deux tubes à essai fermés, A et B, contenant chacun un gaz incolore et inodore issu d'une électrolyse de l'eau (mais on ne sait plus lequel vient de la cathode ni de l'anode). On réalise les deux tests suivants :
\begin{itemize}
    \item \textbf{Test 1 :} On approche une flamme du tube A $\rightarrow$ \textbf{Détonation}.
    \item \textbf{Test 2 :} On plonge une bûchette incandescente dans le tube B $\rightarrow$ \textbf{Rallumage vif}.
\end{itemize}
\textbf{Identifier le gaz dans chaque tube et dire à quelle électrode il était rattaché.}
\tcblower
\textbf{Solution détaillée :}
\begin{enumerate}
    \item \textbf{Tube A :} Le test à la flamme donne une détonation. \textbf{Propriété :} Seul le dihydrogène (\ce{H2}) donne ce test positif. \textbf{Conclusion :} Tube A = \ce{H2}.
    \item \textbf{Tube B :} Le test à la bûchette incandescente donne un rallumage vif. \textbf{Propriété :} Seul le dioxygène (\ce{O2}) entretient la combustion ainsi. \textbf{Conclusion :} Tube B = \ce{O2}.
    \item \textbf{Rattachement aux électrodes :} Le \ce{H2} se dégage à la \textbf{cathode (borne \cle{−})}. Le \ce{O2} se dégage à l'\textbf{anode (borne \cle{+})}.
    \item \textbf{Vérification des volumes :} Si l'expérience a été bien faite, le volume dans le tube A (cathode) doit être environ le double de celui dans le tube B (anode).
\end{enumerate}
\textbf{Réponse finale :} Tube A = \ce{H2} (cathode, \cle{−}) ; Tube B = \ce{O2} (anode, \cle{+}).
\end{exoresolu}

\begin{methode}[Protocole rigoureux pour identifier un gaz inconnu issu d'une électrolyse]
\begin{enumerate}
    \item \textbf{Préparer :} Lunettes, tubes, bûchettes, allumettes. Vérifier l'aération.
    \item \textbf{Tester le gaz « suspect combustible » (gros volume) :} Approcher une \textbf{flamme} à l'orifice. \begin{itemize} \item Détonation $\rightarrow$ \ce{H2} (Cathode, \cle{−}). \item Rien $\rightarrow$ Pas \ce{H2}.
    \end{itemize}
    \item \textbf{Tester l'autre gaz (petit volume) :} Plonger une \textbf{bûchette incandescente} dans le tube. \begin{itemize} \item Rallumage vif $\rightarrow$ \ce{O2} (Anode, \cle{+}). \item Rien $\rightarrow$ Pas \ce{O2} (problème de montage ou gaz impur).
    \end{itemize}
    \item \textbf{Conclure :} Remplir le tableau de synthèse (électrode, gaz, volume, test, résultat).
    \item \textbf{Ranger :} Éteindre les bûchettes dans l'eau. Vider les tubes à l'air libre (pas dans l'évier).
\end{enumerate}
\end{methode}

\begin{exercice}[Entraînement : QCM et rédaction]
\begin{enumerate}
    \item Lors de l'électrolyse de l'eau, le gaz qui se dégage à la borne positive du générateur est : \quad A. \ce{H2} \quad B. \ce{O2} \quad C. \ce{CO2} \quad D. \ce{Cl2}
    \item Le test de reconnaissance du dioxygène consiste à : \quad A. Approcher une flamme \quad B. Plonger une bûchette incandescente \quad C. Verser de l'eau de chaux \quad D. Sentir l'odeur
    \item \textbf{Rédiger :} On a recueilli \SI{20}{mL} de gaz à la cathode et \SI{10}{mL} à l'anode. Expliquer pourquoi les volumes sont dans ce rapport et écrire l'équation bilan de la réaction qui a lieu à la cathode.
\end{enumerate}
\end{exercice}

\begin{corrige}[Exercice n]
\begin{enumerate}
    \item \textbf{B.} L'anode est la borne \cle{+}, c'est le site de l'oxydation : \ce{2 H2O -> O2 + 4 H+ + 4 e-}. Le dioxygène s'y dégage.
    \item \textbf{B.} Le dioxygène est un comburant : il rallume une bûchette incandescente. La flamme (test A) sert pour le dihydrogène.
    \item \textbf{Explication des volumes :} L'équation bilan de l'électrolyse est \ce{2 H2O(l) -> 2 H2(g) + O2(g)}. Les coefficients stœchiométriques montrent que 2 molécules de \ce{H2} sont produites pour 1 molécule de \ce{O2}. Dans les mêmes conditions de température et pression, les volumes des gaz sont proportionnels aux nombres de molécules (loi d'Avogadro). Donc $V_{\ce{H2}} = 2 \times V_{\ce{O2}}$. Ici \SI{20}{mL} = $2 \times$ \SI{10}{mL} : cohérent.
    \item \textbf{Équation à la cathode (réduction) :} Les ions \ce{H+} (issus de l'autoprotolyse de l'eau ou de l'acide ajouté) gagnent des électrons : \ce{2 H+ + 2 e- -> H2}. En milieu neutre, on écrit souvent : \ce{2 H2O + 2 e- -> H2 + 2 HO-}.
\end{enumerate}
\end{corrige}

Nous savons maintenant identifier formellement les produits de l'électrolyse. L'étape suivante, naturelle, est d'écrire l'équation bilan de cette transformation et de comprendre ce qu'elle nous apprend sur la composition intime de la molécule d'eau : c'est l'objet de la section 3, « Interprétation de l'électrolyse : analyse de l'eau ».

\coursec{Interprétation de l'électrolyse : analyse de l'eau}

\noindent
Dans la section précédente, nous avons identifié les deux gaz produits : le \cle{dihydrogène} (H₂) à la cathode (épreuve de la flamme : « pop » caractéristique) et le \cle{dioxygène} (O₂) à l'anode (épreuve de la braise : rallumage). Nous savons \emph{quoi} se forme. Nous devons maintenant comprendre \emph{ce qui arrive à l'eau} : disparaît-elle ? Se transforme-t-elle ? Quelle est la relation entre les quantités de gaz obtenues ? C'est tout l'objet de l'\textbf{interprétation de l'électrolyse}, qui nous mènera à la notion d'\textbf{analyse d'un corps composé}.

\courssub{Bilan qualitatif : l'eau, un corps pur composé}

\noindent
Reprenons l'observation fondamentale de l'expérience : au départ, la cuve ne contient que de l'\textbf{eau acidifiée} (quelques gouttes d'acide sulfurique) ou de l'eau additionnée de sulfate de sodium (électrolyte inerte) pour permettre le passage du courant. Les électrodes sont inertes (graphite ou platine). À la fin, le niveau de la solution a baissé : l'eau a disparu. Deux \textbf{nouveaux corps} apparaissent : le dihydrogène et le dioxygène. Ni l'un ni l'autre n'existait avant le passage du courant.

\begin{definition}[Analyse d'un corps composé]
On appelle \textbf{analyse} d'un corps composé l'opération qui permet de le décomposer en \textbf{corps simples} (corps purs constitués d'un seul élément chimique). L'électrolyse de l'eau est l'analyse de l'eau : elle décompose l'eau (corps pur composé) en dihydrogène et dioxygène (deux corps simples).
\end{definition}

\noindent
Puisque l'eau se décompose en deux corps simples différents (hydrogène et oxygène), l'eau est bien un \textbf{corps pur composé}. Elle n'est pas un mélange (comme l'air ou l'eau de mer), car un mélange ne se décompose pas en corps simples par une réaction chimique unique, mais par des séparations physiques (distillation, filtration, décantation, etc.).

\begin{attention}[Rôle de l'électrolyte]
L'eau pure ne conduit pas l'électricité. On ajoute un \textbf{électrolyte inerte} (acide sulfurique dilué $\ce{H2SO4}$ ou sulfate de sodium $\ce{Na2SO4}$) qui fournit des ions $\ce{H+}$ / $\ce{SO4^2-}$ ou $\ce{Na+}$ / $\ce{SO4^2-}$ pour conduire le courant \textbf{sans être décomposé} aux électrodes (seule l'eau réagit). \textbf{Attention :} Si l'on utilise du chlorure de sodium (sel de cuisine), c'est l'ion chlorure $\ce{Cl-}$ qui s'oxyde à l'anode en donnant du dichlore $\ce{Cl2}$ (gaz toxique) au lieu du dioxygène !
\end{attention}

\courssub{Bilan quantitatif : un rapport volumique constant}

\noindent
Si l'on mesure soigneusement les volumes de gaz recueillis dans les deux tubes à essai gradués (ou dans les seringues à gaz), on constate une régularité frappante, quel que soit le volume d'eau électrolysé, pourvu que la température et la pression restent constantes.

\begin{propriete}[Rapport volumique lors de l'électrolyse de l'eau]
Lors de l'électrolyse de l'eau, le volume de dihydrogène recueilli à la cathode est \textbf{le double} du volume de dioxygène recueilli à l'anode :
\[
V_{\text{dihydrogène}} = 2 \times V_{\text{dioxygène}} \qquad \text{soit} \qquad \frac{V_{\text{H}_2}}{V_{\text{O}_2}} = 2
\]
Ce rapport $2:1$ est constant (loi des volumes de Gay-Lussac appliquée à cette réaction).
\end{propriete}

\begin{exemplebox}[Calcul de volume]
Lors d'une électrolyse, on recueille \SI{6}{cm^3} de dioxygène à l'anode. Quel est le volume de dihydrogène obtenu à la cathode ?
\tcblower
D'après la propriété : $V_{\text{H}_2} = 2 \times V_{\text{O}_2}$.
$V_{\text{H}_2} = 2 \times \SI{6}{cm^3} = \SI{12}{cm^3}$.
On recueille \SI{12}{cm^3} de dihydrogène.
\end{exemplebox}

\begin{popfigure}
\begin{tikzpicture}
\begin{axis}[
    width=\linewidth,
    height=0.5\linewidth,
    xlabel={Temps $t$ (min)},
    ylabel={Volume $V$ (cm$^3$)},
    xmin=0, xmax=10,
    ymin=0, ymax=22,
    xtick={0,2,4,6,8,10},
    ytick={0,5,10,15,20},
    grid=major,
    grid style={popline},
    legend pos=north west,
    legend style={draw=popdark, fill=popcream},
    axis line style={popdark},
    tick style={popdark},
    label style={popdark},
]
\addplot[popblue, thick, domain=0:10, samples=2] {2*x};
\addlegendentry{$V_{\text{H}_2}$ (cathode) : pente 2}
\addplot[poporange, thick, domain=0:10, samples=2] {1*x};
\addlegendentry{$V_{\text{O}_2}$ (anode) : pente 1}
\addplot[popgreen, dashed, thick, domain=0:10, samples=2] {3*x};
\addlegendentry{$V_{\text{total}} = V_{\text{H}_2} + V_{\text{O}_2}$}
\end{axis}
\end{tikzpicture}
\legende{Évolution des volumes de dihydrogène et de dioxygène en fonction du temps d'électrolyse. La droite du dihydrogène a une pente double de celle du dioxygène. Le volume total de gaz produit augmente linéairement avec le temps.}
\end{popfigure}

\noindent
Ce graphique montre que la production de gaz est proportionnelle au temps (donc à la quantité d'électricité passée, d'après les lois de Faraday que vous verrez en Terminale). La pente double pour H₂ confirme visuellement le rapport $2:1$.

\courssub{Écriture de l'équation-bilan : équilibrage pas à pas}

\noindent
Pour traduire chimiquement ce qui se passe, nous écrivons l'\textbf{équation-bilan} de la transformation. Les réactifs (ce qui disparaît) sont à gauche, les produits (ce qui apparaît) à droite.

\begin{methode}[Équilibrer l'équation de l'électrolyse de l'eau]
\emph{Étape 1 : Écrire les formules des réactifs et des produits.}\\
L'eau : $\ce{H2O}$. Le dihydrogène : $\ce{H2}$. Le dioxygène : $\ce{O2}$.
Équation provisoire : $\ce{H2O -> H2 + O2}$

\emph{Étape 2 : Compter les atomes de chaque élément de part et d'autre de la flèche.}\\
À gauche : 2 H, 1 O. À droite : 2 H (dans H₂), 2 O (dans O₂). L'oxygène n'est pas équilibré.

\emph{Étape 3 : Équilibrer l'oxygène en mettant un coefficient devant l'eau.}\\
Il faut 2 atomes O à gauche $\rightarrow$ mettre \textbf{2} devant $\ce{H2O}$.
$\ce{2 H2O -> H2 + O2}$
Bilan O : 2 à gauche, 2 à droite (OK). Bilan H : 4 à gauche, 2 à droite (pas OK).

\emph{Étape 4 : Équilibrer l'hydrogène en mettant un coefficient devant le dihydrogène.}\\
Il faut 4 atomes H à droite $\rightarrow$ mettre \textbf{2} devant $\ce{H2}$.
$\ce{2 H2O -> 2 H2 + O2}$

\emph{Étape 5 : Vérifier le bilan final.}\\
Atomes H : $2 \times 2 = 4$ à gauche ; $2 \times 2 = 4$ à droite. \checkmark\\
Atomes O : $2 \times 1 = 2$ à gauche ; $1 \times 2 = 2$ à droite. \checkmark
L'équation est équilibrée.
\end{methode}

\noindent
L'équation-bilan définitive de l'électrolyse de l'eau s'écrit donc :
\[
\ce{2 H2O ->[courant electrique] 2 H2 + O2}
\]

\begin{exoresolu}[Équilibrage d'une équation similaire : décomposition de l'eau oxygénée]
L'électrolyse (ou décomposition catalytique) de l'eau oxygénée (peroxyde d'hydrogène, $\ce{H2O2}$) produit de l'eau et du dioxygène. Équilibrer l'équation : $\ce{H2O2 -> H2O + O2}$.
\tcblower
\textbf{Solution :}
1. Formules : $\ce{H2O2 -> H2O + O2}$.
2. Bilan initial : Gauche : 2 H, 2 O. Droite : 2 H, 3 O (1 dans $\ce{H2O}$ + 2 dans $\ce{O2}$).
3. Équilibrer l'oxygène : Mettre \textbf{2} devant $\ce{H2O2}$ pour avoir 4 O à gauche.
   $\ce{2 H2O2 -> H2O + O2}$ (Gauche : 4 H, 4 O).
4. Équilibrer l'hydrogène : 4 H à gauche $\rightarrow$ mettre \textbf{2} devant $\ce{H2O}$.
   $\ce{2 H2O2 -> 2 H2O + O2}$.
5. Vérification finale : O : Gauche $2 \times 2 = 4$, Droite $2 \times 1 + 2 = 4$. H : Gauche 4, Droite 4. \checkmark
\textbf{Équation finale :} $\ce{2 H2O2 -> 2 H2O + O2}$.
\end{exoresolu}

\courssub{Interprétation microscopique : la molécule d'eau « cassée » par le courant}

\noindent
L'équation $\ce{2 H2O -> 2 H2 + O2}$ est une écriture \textbf{macroscopique} (masses, volumes mesurables). Que se passe-t-il à l'échelle de l'infiniment petit, au niveau des \textbf{molécules} et des \textbf{atomes} ? Le courant électrique apporte l'énergie nécessaire pour \textbf{briser les liaisons} à l'intérieur des molécules d'eau. Les atomes libérés se regroupent ensuite différemment pour former les molécules stables des gaz dihydrogène et dioxygène.

\begin{popfigure}
\begin{tikzpicture}[
    atomH/.style={circle, draw=popdark, fill=popblueL, minimum size=0.7cm, inner sep=0, font=\small},
    atomO/.style={circle, draw=popdark, fill=poporangeL, minimum size=0.9cm, inner sep=0, font=\small},
    bond/.style={thick, popdark},
    dblbond/.style={double distance=1.5pt, thick, popdark},
    arrow/.style={->, >=Stealth, very thick, poppurple},
    labelbox/.style={rectangle, draw=popdark, fill=popcream, rounded corners, inner sep=4pt, font=\small, align=center}
]
% Réactifs : 2 H2O
\node[atomO] (O1) at (0, 1.5) {O};
\node[atomH] (H11) at (-0.85, 0.75) {H};
\node[atomH] (H12) at (0.85, 0.75) {H};
\draw[bond] (O1) -- (H11);
\draw[bond] (O1) -- (H12);

\node[atomO] (O2) at (0, -0.5) {O};
\node[atomH] (H21) at (-0.85, -1.25) {H};
\node[atomH] (H22) at (0.85, -1.25) {H};
\draw[bond] (O2) -- (H21);
\draw[bond] (O2) -- (H22);

\node[labelbox, align=center] at (-2.5, 0.5){2 molécules\\d'eau $\ce{H2O}$\\Total : 4 H, 2 O};

% Flèche réaction
\draw[arrow] (2.2, 0.5) -- (4.3, 0.5) node[midway, above, font=\small] {Courant électrique};

% Produits : 2 H2 + 1 O2
% H2 (1)
\node[atomH] (H31) at (5.8, 1.8) {H};
\node[atomH] (H32) at (6.8, 1.8) {H};
\draw[bond] (H31) -- (H32);
% H2 (2)
\node[atomH] (H41) at (5.8, 0.5) {H};
\node[atomH] (H42) at (6.8, 0.5) {H};
\draw[bond] (H41) -- (H42);
% O2 (double liaison)
\node[atomO] (O3) at (6.3, -1) {O};
\node[atomO] (O4) at (7.5, -1) {O};
\draw[dblbond] (O3) -- (O4);

\node[labelbox, align=center] at (9.2, 0.5){2 molécules\\de dihydrogène $\ce{H2}$\\+ 1 molécule\\de dioxygène $\ce{O2}$\\Total : 4 H, 2 O};
\end{tikzpicture}
\legende{Modélisation microscopique de l'électrolyse : 2 molécules d'eau ($\ce{H2O}$) sont décomposées. Les 4 atomes d'hydrogène (bleu) forment 2 molécules $\ce{H2}$ (liaison simple) ; les 2 atomes d'oxygène (orange) forment 1 molécule $\ce{O2}$ (liaison double). \textbf{Les atomes sont conservés, ils se réorganisent.}}
\end{popfigure}

\noindent
Ce schéma illustre deux principes fondamentaux de la chimie (loi de Lavoisier) :
\begin{itemize}
    \item \textbf{Conservation des atomes :} Il y a exactement 4 atomes d'hydrogène et 2 atomes d'oxygène avant et après la transformation. Aucun atome n'est créé, aucun n'est détruit.
    \item \textbf{Réorganisation des atomes :} Les liaisons O–H (dans l'eau) sont rompues. De nouvelles liaisons H–H (dans $\ce{H2}$) et O=O (dans $\ce{O2}$) se forment. C'est cela une \textbf{transformation chimique} : un réarrangement des atomes.
\end{itemize}

\courssub{Transformation chimique ou transformation physique ?}

\noindent
C'est un point crucial pour le BEPC : ne pas confondre \textbf{électrolyse} (transformation chimique) et \textbf{vaporisation / ébullition} (transformation physique).

\begin{attention}[Piège fréquent : Électrolyse $\neq$ Ébullition]
\begin{itemize}
    \item \textbf{Ébullition (vaporisation) :} L'eau liquide devient vapeur d'eau (gaz). \textbf{La molécule $\ce{H2O}$ reste intacte.} C'est un changement d'état (physique). On peut récupérer l'eau liquide par condensation. Formule : $\ce{H2O_{(l)} -> H2O_{(g)}}$.
    \item \textbf{Électrolyse :} L'eau \textbf{disparaît} en tant que substance. De \textbf{nouvelles substances} apparaissent ($\ce{H2}$ et $\ce{O2}$). Les molécules $\ce{H2O}$ sont \textbf{cassées}. C'est une transformation chimique (irréversible par simple changement de température). Formule : $\ce{2 H2O -> 2 H2 + O2}$.
\end{itemize}
\textbf{Astuce mnémotechnique :} « \emph{Bouillir, c'est s'envoler ; électrolyser, c'est se dédoubler.} »
\end{attention}

\noindent
Dans la vie courante au Cameroun, quand vous faites chauffer de l'eau dans une casserole pour le thé ou sur un réchaud à gaz (butane), vous provoquez une vaporisation. Quand vous utilisez un \textbf{générateur d'hydrogène} (parfois testé pour les motos-taxis ou les groupes électrogènes) ou quand une \textbf{bonbonne de gaz} fuit et qu'une étincelle fait exploser le mélange air/gaz (synthèse de l'eau, section suivante), ce sont des transformations chimiques.

\begin{attention}[Sécurité : Gaz inflammables et comburants]
Le dihydrogène est \textbf{très inflammable} (explosif en mélange avec l'air). Le dioxygène est un \textbf{comburant puissant} (active la combustion). \textbf{Ne jamais} approcher une flamme de l'appareil d'électrolyse en fonctionnement. Faire l'épreuve de la flamme (pop) \textbf{uniquement} sur un petit tube à essai prélevé avec précaution, loin de la cuve. Aérer le local.
\end{attention}

\courssub{Conclusion : l'analyse de l'eau}

\noindent
En résumé, l'électrolyse de l'eau est l'expérience de référence qui prouve que l'eau est un \textbf{corps pur composé} d'hydrogène et d'oxygène.
\begin{itemize}
    \item \textbf{Qualitatif :} Décomposition en deux corps simples ($\ce{H2}$, $\ce{O2}$) $\rightarrow$ \textbf{Analyse}.
    \item \textbf{Quantitatif :} Rapport volumique constant $V_{\ce{H2}} / V_{\ce{O2}} = 2 / 1$.
    \item \textbf{Microscopique :} Rupture des liaisons O–H, formation des liaisons H–H et O=O. Conservation stricte des atomes (4 H, 2 O).
    \item \textbf{Écriture :} $\ce{2 H2O -> 2 H2 + O2}$ (équation-bilan équilibrée).
\end{itemize}

\noindent
Cette analyse de l'eau par l'électricité est l'opération inverse de la \textbf{synthèse de l'eau} (réaction du dihydrogène et du dioxygène pour reformer de l'eau), que nous étudierons dans la prochaine section. C'est ce couple analyse/synthèse qui a permis à Lavoisier, à la fin du XVIII\textsuperscript{e} siècle, de comprendre la composition de l'eau et de fonder la chimie moderne.

\begin{savaistu}[Lavoisier et la fin du « phlogistique »]
Avant Lavoisier (1783), les chimistes croyaient au « phlogistique », une substance imaginaire qui quitterait les corps lors de la combustion. En décomposant l'eau (analyse) et en la reformant (synthèse) avec une précision de masse parfaite, Lavoisier a prouvé que l'eau n'est pas un élément, mais un composé d'hydrogène (« qui engendre l'eau ») et d'oxygène (« qui engendre l'acide »). Il a ainsi tué la théorie du phlogistique. Au Cameroun, l'eau de la SNEC ou des sources (Tangui, Supermont) a partout la même composition : $\ce{H2O}$, fruit de cette histoire scientifique.
\end{savaistu}

\coursec{Électrolyse et synthèse de l'eau}

\courssub{Mise en évidence expérimentale : l'électrolyse de l'eau}

Pour comprendre de quoi l'eau est faite, nous allons la \emph{décomposer} à l'aide de l'électricité. C'est l'expérience fondatrice de la chimie de l'eau.

\begin{experience}[Électrolyse de l'eau acidifiée]
\textbf{Matériel :} un bac d'électrolyse (ou un verre haut), deux électrodes en graphite (ou en platine), une pile \SI{9}{V} ou un générateur basse tension (\SI{12}{V} max), des fils, deux tubes à essai, de l'eau additionnée de quelques gouttes d'acide sulfurique dilué (ou de sulfate de sodium \ce{Na2SO4}) pour la rendre conductrice.

\textbf{Protocole :}
\begin{enumerate}
\item Remplir le bac et les deux tubes à essai avec la solution conductrice. Inverser les tubes à essai sur les électrodes pour les remplir complètement (pas de bulle d'air).
\item Brancher les électrodes à la pile : l'électrode reliée au \textbf{plus (+)} est l'\cle{anode}, celle reliée au \textbf{moins (-)} est la \cle{cathode}.
\item Laisser le courant passer pendant quelques minutes.
\item Observer le volume de gaz recueilli dans chaque tube à essai.
\item Retirer délicatement les tubes (en les gardant à l'envers) pour tester les gaz.
\end{enumerate}

\textbf{Observations :}
\begin{itemize}
\item Des bulles se forment sur \textbf{les deux électrodes} dès le passage du courant.
\item Le volume de gaz recueilli à la \textbf{cathode (-)} est \textbf{double} de celui recueilli à l'anode (+). On note environ \textbf{2 volumes pour 1 volume}.
\item Si on coupe le courant, la production de gaz s'arrête immédiatement.
\end{itemize}

\textbf{Interprétation :} Le courant électrique a décomposé l'eau. Il se forme un gaz à chaque électrode. Le rapport de volumes 2:1 est une première indication sur la composition de l'eau.
\end{experience}

\begin{popfigure}
\begin{tikzpicture}[scale=0.9]
% Bac
\draw[thick, popdark] (-3,-0.5) -- (-3,4) -- (3,4) -- (3,-0.5);
\draw[popblue!30] (-2.8,0.2) rectangle (2.8,3.5);
% Électrodes
\draw[thick, popdark] (-1.5,4) -- (-1.5,5.5) node[above, font=\small] {Anode (+)};
\draw[thick, popdark] (1.5,4) -- (1.5,5.5) node[above, font=\small] {Cathode (-)};
\fill[popmuted] (-1.6,3.5) rectangle (-1.4,4);
\fill[popmuted] (1.4,3.5) rectangle (1.6,4);
% Bulles
\foreach \y in {1,1.5,2,2.5,3} {
  \fill[poporange!50] (-1.7,\y) circle (0.08);
  \fill[poporange!50] (-1.3,\y) circle (0.06);
  \fill[popblue!50] (1.3,\y) circle (0.08);
  \fill[popblue!50] (1.7,\y) circle (0.06);
  \fill[popblue!50] (1.5,\y+0.15) circle (0.05);
}
% Tubes à essai
\draw[thick, popink] (-2,3.5) -- (-2,6.5) arc (180:0:0.5) -- (-1,3.5);
\node[font=\small, popblue] at (-1.5,5.5) {Gaz 1};
\draw[thick, popink] (1,3.5) -- (1,7.5) arc (180:0:0.5) -- (2,3.5);
\node[font=\small, popblue] at (1.5,6.2) {Gaz 2 (2$\times$ vol)};
% Pile
\draw[fill=popgold!30] (-4,1) rectangle (-3.2,3) node[midway, font=\small, rotate=90] {Pile 9V};
\draw[thick, -latex] (-3.2,2) -- (-3,2);
\draw[thick, -latex] (-3,1.5) -- (-3.2,1.5);
% Légendes gaz
\node[font=\small, align=center] at (-4.5,5) {Anode (+)\\\ce{O2} (1 vol)};
\node[font=\small, align=center] at (4.5,6) {Cathode (-)\\\ce{H2} (2 vol)};
\end{tikzpicture}
\legende{Montage de l'électrolyse de l'eau acidifiée. Le gaz se forme à l'anode (+) et à la cathode (-). Le volume à la cathode est double de celui à l'anode.}
\end{popfigure}

\courssub{Identification des gaz : tests du dihydrogène et du dioxygène}

Recueillir des gaz ne suffit pas : il faut les \emph{identifier} par des tests caractéristiques, réalisés immédiatement après l'électrolyse.

\begin{methode}[Test du dihydrogène \ce{H2} (gaz de la cathode)]
\begin{enumerate}
\item Approcher une \textbf{flamme} (allumette ou bec Bunsen) de l'embouchure du tube à essai contenant le gaz de la cathode.
\item \textbf{Observation :} on entend un \textbf{bruit sec et bref} (un « pop » ou \cle{détonation} légère) et on voit une flamme bleutée à l'intérieur du tube.
\item \textbf{Conclusion :} ce gaz est le \cle{dihydrogène} (\ce{H2}). C'est un gaz \emph{inflammable} qui brûle sans flamme visible forte, mais avec détonation en mélange avec l'air.
\end{enumerate}
\end{methode}

\begin{methode}[Test du dioxygène \ce{O2} (gaz de l'anode)]
\begin{enumerate}
\item Allumer une \textbf{baguette incandescente} (baguette de bois qui braise, sans flamme).
\item La plonger rapidement dans le tube à essai contenant le gaz de l'anode.
\item \textbf{Observation :} la baguette \textbf{se rallume} instantanément avec une flamme vive.
\item \textbf{Conclusion :} ce gaz est le \cle{dioxygène} (\ce{O2}). C'est un gaz \emph{comburant} (qui entretient la combustion).
\end{enumerate}
\end{methode}

\begin{aretenir}[Identification des gaz de l'électrolyse]
\begin{center}
\begin{tabularx}{\linewidth}{|l|X|X|}\hline
\rowcolor{popblueL}
\textbf{Gaz} & \textbf{Électrode} & \textbf{Test positif} \\ \hline
Dihydrogène \ce{H2} & Cathode (-) & Détonation à la flamme (bruit « pop ») \\ \hline
Dioxygène \ce{O2} & Anode (+) & Rallumage d'une baguette incandescente \\ \hline
\end{tabularx}
\end{center}
\end{aretenir}

\courssub{Interprétation de l'électrolyse : analyse de l'eau}

L'expérience montre qu'en faisant passer un courant électrique dans l'eau, on obtient deux corps simples : le dihydrogène et le dioxygène. C'est une \cle{transformation chimique} : l'eau a disparu, de nouvelles substances sont apparues.

\begin{definition}[Analyse de l'eau]
L'\cle{analyse de l'eau} est sa décomposition en corps simples (dihydrogène et dioxygène) sous l'action d'un courant électrique. C'est l'\cle{électrolyse de l'eau}.
Son équation-bilan est :
\[ \ce{2H2O ->[electricite] 2H2 + O2} \]
\end{definition}

\begin{propriete}[Bilan de l'électrolyse]
\begin{itemize}
\item \textbf{Réactif :} l'eau (\ce{H2O}), corps \emph{composé}.
\item \textbf{Produits :} le dihydrogène (\ce{H2}) et le dioxygène (\ce{O2}), corps \emph{simples}.
\item \textbf{Énergie :} il faut \textbf{apporter} de l'énergie électrique pour que la réaction ait lieu. On dit que la réaction est \cle{endothermique} (au sens large : absorption d'énergie).
\item \textbf{Volumes gazeux :} 2 volumes de dihydrogène pour 1 volume de dioxygène (loi des volumes gazeux d'Avogadro).
\end{itemize}
\end{propriete}

\begin{attention}[Écriture de l'équation-bilan]
\begin{itemize}
\item La flèche \ce{->} est surmontée de la condition : \ce{->[electricite]} ou \ce{->[courant electrique]}. N'écris pas « chaleur » ni « flamme ».
\item Les coefficients stœchiométriques (2 devant \ce{H2O} et \ce{H2}, 1 devant \ce{O2}) sont \textbf{obligatoires} pour respecter la conservation des atomes (4 atomes H, 2 atomes O de chaque côté).
\item L'eau est liquide (\ce{(l)}), les gaz sont gazeux (\ce{(g)}) : \ce{2H2O(l) -> 2H2(g) + O2(g)}.
\end{itemize}
\end{attention}

\courssub{La synthèse de l'eau}

Dans la section précédente, nous avons \emph{démonté} l'eau : grâce à l'électrolyse, nous avons cassé les molécules d'eau et recueilli deux gaz, le dihydrogène \ce{H2} et le dioxygène \ce{O2}. Cette opération s'appelle l'\cle{analyse} de l'eau. Une question se pose alors naturellement : peut-on faire le chemin inverse ? Peut-on \emph{remonter} l'eau, c'est-à-dire fabriquer de l'eau à partir de ces deux gaz ? C'est exactement ce que nous allons découvrir : c'est la \cle{synthèse} de l'eau.

\courssub{Principe : l'opération inverse de l'analyse}

Imagine que tu démontes un jouet : tu obtiens des pièces séparées. Si tu remontes les pièces dans le bon ordre, tu retrouves le jouet. En chimie, c'est pareil :

\begin{itemize}
\item l'\cle{analyse} de l'eau (électrolyse) \emph{décompose} l'eau en dihydrogène et dioxygène : \ce{2H2O -> 2H2 + O2} ;
\item la \cle{synthèse} de l'eau \emph{recompose} l'eau en combinant le dihydrogène et le dioxygène : \ce{2H2 + O2 -> 2H2O}.
\end{itemize}

\begin{definition}[Synthèse de l'eau]
La \cle{synthèse de l'eau} est la formation d'eau par combinaison chimique du dihydrogène et du dioxygène. Son équation-bilan est :
\[ \ce{2H2 + O2 -> 2H2O} \]
C'est la réaction \emph{inverse} de l'électrolyse de l'eau.
\end{definition}

Plus généralement, la \cle{synthèse} d'un corps composé est la formation de ce corps composé à partir des corps simples qui le constituent. Ici, l'eau (corps composé) est synthétisée à partir du dihydrogène et du dioxygène (corps simples).

\courssub{Expérience : le mélange détonant}

\begin{experience}[Synthèse de l'eau dans un eudiomètre]
\textbf{Matériel :} un \cle{eudiomètre} (tube gradué fermé à une extrémité, muni de deux fils électriques traversant sa paroi pour produire une étincelle), du dihydrogène et du dioxygène (par exemple ceux recueillis lors de l'électrolyse), un générateur d'étincelles (bobine de Ruhmkorff ou piezo), un écran de protection, lunettes de sécurité.

\textbf{Protocole :}
\begin{enumerate}
\item Introduire dans l'eudiomètre \textbf{2 volumes de dihydrogène} et \textbf{1 volume de dioxygène} (par exemple \SI{20}{cm^3} de \ce{H2} et \SI{10}{cm^3} de \ce{O2}).
\item Fermer le tube et le placer \textbf{derrière l'écran de protection}.
\item Produire une étincelle entre les deux fils électriques.
\item Observer les parois internes du tube.
\end{enumerate}

\textbf{Observations :}
\begin{itemize}
\item on entend une \textbf{détonation} (un bruit sec et violent) : la réaction est explosive ;
\item une \textbf{buée} apparaît sur les parois froides du tube, puis de fines \textbf{gouttelettes de liquide} s'y déposent ;
\item si les proportions étaient exactement de 2 volumes de \ce{H2} pour 1 volume de \ce{O2}, \textbf{aucun gaz ne subsiste} : la réaction est totale, sans résidu.
\end{itemize}

\textbf{Interprétation :} le dihydrogène et le dioxygène ont réagi ensemble pour former un liquide. Reste à identifier ce liquide.
\end{experience}

\begin{popfigure}
\begin{tikzpicture}[scale=0.95]
% Tube de l'eudiomètre
\draw[thick, popink] (0,0) -- (0,5.5) arc (180:0:0.9) -- (1.8,0);
\draw[thick, popink] (0,0) -- (1.8,0);
% Graduations
\foreach \y in {1,2,3,4,5} {\draw[popmuted] (0,\y) -- (0.25,\y);}
\node[left, popmuted, font=\small] at (0,1) {10};
\node[left, popmuted, font=\small] at (0,2) {20};
\node[left, popmuted, font=\small] at (0,3) {30};
\node[left, popmuted, font=\small] at (0,4) {40};
\node[left, popmuted, font=\small] at (0,5) {50};
% Gaz H2 (2 volumes) et O2 (1 volume) - AVANT réaction
\fill[popblue!20] (0.05,1.7) rectangle (1.75,5.2);
\fill[poporange!25] (0.05,0.05) rectangle (1.75,1.7);
\node[popblue, font=\small] at (0.9,3.6) {2 volumes \ce{H2}};
\node[poporange, font=\small] at (0.9,0.9) {1 volume \ce{O2}};
% Fils électriques et étincelle
\draw[thick, popdark] (-1.2,0.35) -- (0.6,0.35);
\draw[thick, popdark] (-1.2,0.75) -- (0.55,0.75);
\draw[popgold, thick, decorate, decoration={zigzag, segment length=2.5, amplitude=1.2}] (0.6,0.35) -- (0.9,0.55) -- (0.55,0.75);
\node[popdark, font=\small, align=center] at (-1.9,0.55) {étincelle};
% Gouttelettes d'eau sur les parois (APRÈS réaction)
\foreach \p in {(0.15,4.6),(1.6,4.2),(0.2,3.1),(1.62,2.6),(0.18,2.2),(1.55,1.8)} {\fill[popblue] \p circle (0.04);}
\node[popblue, font=\small, align=left] at (3.8,4.4) {gouttelettes\\ d'eau};
\draw[->, popblue] (3.0,4.3) -- (1.7,4.2);
% Détonation
\node[popgold, font=\small\bfseries] at (3.8,0.6) {détonation !};
\draw[->, popgold] (2.8,0.6) -- (1.2,0.55);
\end{tikzpicture}
\legende{Schéma de l'eudiomètre : le mélange de 2 volumes de dihydrogène et 1 volume de dioxygène est enflammé par une étincelle ; après la détonation, des gouttelettes d'eau apparaissent sur les parois froides.}
\end{popfigure}

\courssub{Identification du produit formé : c'est de l'eau}

En science, observer un liquide ne suffit pas : il faut prouver son identité. Deux tests permettent de montrer que le liquide obtenu est bien de l'eau :

\begin{enumerate}
\item \textbf{Test au sulfate de cuivre anhydre :} quelques gouttes du liquide déposées sur du sulfate de cuivre \emph{anhydre} (poudre blanche) le font \textbf{bleuir} (formation de sulfate de cuivre hydraté \ce{CuSO4.5H2O}). C'est le test caractéristique de l'eau.
\item \textbf{Mesure du point d'ébullition :} le liquide bout à \SI{100}{\celsius} sous pression atmosphérique normale (\SI{1013}{hPa}), comme l'eau pure.
\end{enumerate}

\begin{propriete}[Le mélange détonant]
Le mélange de \textbf{2 volumes de dihydrogène} et \textbf{1 volume de dioxygène} est appelé \cle{mélange détonant}. Enflammé par une étincelle ou une flamme, il réagit de façon explosive et sa combustion donne \textbf{exclusivement de l'eau} :
\[ \ce{2H2 + O2 -> 2H2O} \]
Cette réaction est très \cle{exothermique} : elle libère une grande quantité d'énergie (chaleur et lumière), ce qui explique la détonation.
\end{propriete}

\courssub{Les proportions stœchiométriques : 2 volumes de H2 pour 1 volume de O2}

Regardons l'équation-bilan de près : \ce{2H2 + O2 -> 2H2O}. Elle nous dit que \textbf{2 « doses » de dihydrogène réagissent avec exactement 1 « dose » de dioxygène}. Comme, pour les gaz, les volumes sont proportionnels au nombre de molécules (loi d'Avogadro), on retient :

\begin{aretenir}[Proportions de la synthèse de l'eau]
La synthèse de l'eau est \textbf{totale et sans résidu} uniquement si l'on mélange \textbf{2 volumes de dihydrogène pour 1 volume de dioxygène}. Ce sont les \cle{proportions stœchiométriques}. Si l'un des deux gaz est en excès, il reste du gaz non réagi après la réaction.
\end{aretenir}

\begin{definition}[Réactif limitant et réactif en excès]
Dans un mélange non stœchiométrique, le \cle{réactif limitant} est celui qui est entièrement consommé et qui détermine la quantité de produit formé. L'autre réactif est en \cle{excès} : il en reste après la réaction.
\end{definition}

\begin{exemplebox}[Un mélange mal proportionné]
Si l'on mélange \SI{30}{cm^3} de \ce{H2} avec \SI{30}{cm^3} de \ce{O2} : les \SI{30}{cm^3} de \ce{H2} ne peuvent réagir qu'avec \SI{15}{cm^3} de \ce{O2} (proportion 2/1). Après la réaction, il reste $30 - 15 = \SI{15}{cm^3}$ de dioxygène en excès. Le dihydrogène, lui, a entièrement disparu : on dit qu'il est le \cle{réactif limitant}.
\end{exemplebox}

\begin{exoresolu}[Volume de gaz restant après la synthèse]
Énoncé : on introduit dans un eudiomètre \SI{40}{cm^3} de dihydrogène et \SI{30}{cm^3} de dioxygène, puis on produit une étincelle. La réaction est-elle totale ? Quel gaz reste-t-il, et en quel volume ?
\tcblower
Solution détaillée :
\begin{enumerate}
\item La réaction exige 2 volumes de \ce{H2} pour 1 volume de \ce{O2}.
\item Avec \SI{40}{cm^3} de \ce{H2}, le volume de \ce{O2} qui peut réagir est : $\dfrac{40}{2} = \SI{20}{cm^3}$.
\item Or on dispose de \SI{30}{cm^3} de \ce{O2} : le dioxygène est donc en excès.
\item Volume de dioxygène restant : $30 - 20 = \SI{10}{cm^3}$.
\end{enumerate}
\textbf{Conclusion :} seuls \SI{20}{cm^3} de \ce{O2} réagissent avec les \SI{40}{cm^3} de \ce{H2}. Il reste \SI{10}{cm^3} de dioxygène non réagi. Le dihydrogène est le réactif limitant.
\end{exoresolu}

\courssub{Analyse et synthèse : deux réactions inverses}

Nous pouvons maintenant boucler la boucle. L'eau peut être \emph{décomposée} par l'électricité (analyse) puis \emph{recomposée} par une étincelle (synthèse) :

\begin{popfigure}
\begin{tikzpicture}[scale=1]
\node[draw, thick, rounded corners, fill=popblue!20, minimum width=2.8cm, minimum height=1.2cm, font=\bfseries] (eau) at (0,0) {EAU \ce{H2O(l)}};
\node[draw, thick, rounded corners, fill=poporange!20, minimum width=4.5cm, minimum height=1.2cm, font=\bfseries, align=center] (gaz) at (8,0) {DIHYDROGÈNE + DIOXYGÈNE\\ \ce{2H2(g)} + \ce{O2(g)}};
\draw[->, very thick, popblue] (eau.north) .. controls (2.5,1.8) and (5.5,1.8) .. node[above, font=\small, align=center]{ÉLECTROLYSE (analyse)\\ énergie électrique absorbée} (gaz.north);
\draw[->, very thick, poporange] (gaz.south) .. controls (5.5,-1.8) and (2.5,-1.8) .. node[below, font=\small, align=center]{SYNTHÈSE (mélange détonant)\\ énergie libérée (explosion)} (eau.south);
\end{tikzpicture}
\legende{Schéma-bilan en boucle : l'électrolyse décompose l'eau en dihydrogène et dioxygène ; la synthèse reconstitue l'eau à partir de ces deux gaz. Les deux réactions sont inverses l'une de l'autre.}
\end{popfigure}

\begin{aretenir}[Énergie des réactions inverses (savoir admis)]
L'électrolyse et la synthèse de l'eau sont deux réactions \emph{inverses}. On admet que \textbf{l'énergie libérée par la synthèse de l'eau est égale à l'énergie absorbée par l'électrolyse} d'une même quantité d'eau. L'énergie dépensée pour « casser » l'eau est exactement celle que l'on récupère quand on la « refabrique ».
\end{aretenir}

\courssub{Formule et modèle de la molécule d'eau}

L'analyse et la synthèse nous ont montré que l'eau est composée d'hydrogène et d'oxygène. Mais comment les atomes sont-ils arrangés ?

\begin{propriete}[Formule chimique de l'eau]
La formule de l'eau est \ce{H2O}. Cela signifie qu'une molécule d'eau contient \textbf{2 atomes d'hydrogène} (H) et \textbf{1 atome d'oxygène} (O).
\end{propriete}

\begin{experience}[Modèle de la molécule d'eau (représentation de Lewis)]
On représente les électrons de la couche externe (électrons de valence) par des points ou des traits.
\begin{itemize}
\item L'atome d'oxygène (O) a 6 électrons de valence. Il lui en manque 2 pour atteindre la configuration stable (8 électrons, règle de l'octet).
\item Chaque atome d'hydrogène (H) a 1 électron de valence. Il lui en manque 1 pour atteindre la configuration stable du gaz rare hélium (2 électrons, règle du duet).
\item L'oxygène partage un électron avec chaque hydrogène : il se forme \textbf{2 liaisons covalentes simples} \ce{O-H}.
\end{itemize}
\end{experience}

\begin{popfigure}
\begin{tikzpicture}[scale=1.2]
% Modèle de Lewis
\node[font=\small] at (0,1.5) {Représentation de Lewis :};
\draw[thick, popdark] (-1.5,0) circle (0.6);
\node at (-1.5,0) {O};
\foreach \angle/\label in {60/., 120/., 240/., 300/.} {
  \draw[popmuted] (-1.5,0) -- ++(\angle:0.8) node[font=\tiny] {\label};
}
\draw[thick, popblue] (-1.5,0) -- ++(-60:0.8) node[below right, font=\small] {H};
\draw[thick, popblue] (-1.5,0) -- ++(-120:0.8) node[below left, font=\small] {H};
% Liaisons covalentes (traits)
\draw[very thick, popblue] (-1.5,0) -- ++(-60:0.6);
\draw[very thick, popblue] (-1.5,0) -- ++(-120:0.6);
\node[popblue, font=\small] at (-1.5,-1.5) {\ce{H-O-H}};

% Géométrie coudée (modèle VSEPR simplifié)
\node[font=\small] at (4,1.5) {Géométrie (modèle VSEPR) :};
\draw[thick, popdark] (4,0) circle (0.5);
\node at (4,0) {O};
\draw[very thick, popblue] (4,0) -- ++(105:1) node[above right, font=\small] {H};
\draw[very thick, popblue] (4,0) -- ++(-105:1) node[below right, font=\small] {H};
\draw[dashed, popmuted] (4,0) -- ++(105:1.3);
\draw[dashed, popmuted] (4,0) -- ++(-105:1.3);
\draw[<->, popmuted] (5.05,0.95) -- (5.05,-0.95) node[midway, right, font=\small] {$\approx 104,5^\circ$};
\node[popdark, font=\small, align=center] at (4,-1.8) {Molécule \textbf{coudée}\\(angle \SI{104,5}{\degree})};
\end{tikzpicture}
\legende{À gauche : représentation de Lewis (partage d'électrons). À droite : géométrie réelle de la molécule d'eau, coudée avec un angle de \SI{104,5}{\degree}. Les deux paires non liantes sur l'oxygène repoussent les liaisons O-H.}
\end{popfigure}

\begin{aretenir}[Molécule d'eau \ce{H2O}]
\begin{itemize}
\item \textbf{Formule :} \ce{H2O}.
\item \textbf{Liaisons :} 2 liaisons covalentes simples \ce{O-H}.
\item \textbf{Géométrie :} \textbf{coudée} (angle \SI{104,5}{\degree}).
\item \textbf{Polarité :} L'oxygène attire plus fort les électrons que l'hydrogène. La molécule est \textbf{polaire} (charge partielle négative $\delta^-$ sur O, charges partielles positives $\delta^+$ sur H). C'est cette polarité qui explique le pouvoir solvant exceptionnel de l'eau.
\end{itemize}
\end{aretenir}

\courssub{Applications et sécurité}

\begin{attention}[La synthèse de l'eau est explosive !]
La réaction entre le dihydrogène et le dioxygène est \textbf{extrêmement violente}. Elle ne doit être réalisée :
\begin{itemize}
\item qu'avec de \textbf{très petits volumes} de gaz (quelques centimètres cubes) ;
\item que \textbf{derrière un écran de protection}, avec des lunettes de sécurité ;
\item que par le professeur ou sous sa surveillance directe.
\end{itemize}
Ne cherche jamais à enflammer du dihydrogène chez toi : c'est le même type de danger qu'une fuite de gaz dans une cuisine (gaz butane/propane + air) !
\end{attention}

\begin{savaistu}[Cavendish et Lavoisier : l'eau n'est pas un élément]
Pendant des siècles, on a cru que l'eau était un \emph{élément}, c'est-à-dire un corps simple indécomposable. En 1781, le physicien anglais Henry \cle{Cavendish} fit exploser un mélange de « gaz inflammable » (dihydrogène) et d'air et obtint... de l'eau ! Le chimiste français Antoine de \cle{Lavoisier} reprit et interpréta ces expériences : par la \textbf{synthèse} puis l'\textbf{analyse} de l'eau, il démontra que l'eau est un \emph{corps composé} d'hydrogène et d'oxygène. C'est d'ailleurs Lavoisier qui baptisa l'hydrogène, mot qui signifie littéralement « qui engendre l'eau ». Aujourd'hui encore, la réaction \ce{2H2 + O2 -> 2H2O} est utilisée dans les fusées spatiales : les moteurs de la fusée Ariane (étage Vulcain) brûlent du dihydrogène et du dioxygène liquides, et leur seul rejet est... de la vapeur d'eau ! Au Cameroun, la SONARA produit de l'essence et du gasoil, mais l'hydrogène pourrait devenir un carburant d'avenir « propre » (pile à combustible).
\end{savaistu}

\begin{exercice}[À toi de jouer]
On réalise la synthèse de l'eau avec \SI{50}{cm^3} de dihydrogène et \SI{20}{cm^3} de dioxygène.
\begin{enumerate}
\item Écris l'équation-bilan de la réaction.
\item Quel est le réactif limitant ? Quel volume de gaz reste-t-il après la réaction ?
\item Cite deux tests permettant d'identifier le liquide obtenu.
\end{enumerate}
\end{exercice}

\begin{corrige}[Exercice : synthèse de l'eau]
\begin{enumerate}
\item \ce{2H2 + O2 -> 2H2O}.
\item Il faut 2 volumes de \ce{H2} pour 1 volume de \ce{O2}. Avec \SI{20}{cm^3} de \ce{O2}, seuls $2 \times 20 = \SI{40}{cm^3}$ de \ce{H2} peuvent réagir. Le dioxygène est donc le réactif limitant. Il reste $50 - 40 = \SI{10}{cm^3}$ de dihydrogène en excès.
\item Le sulfate de cuivre anhydre blanc bleuit au contact du liquide ; le liquide bout à \SI{100}{\celsius} sous pression normale. C'est de l'eau.
\end{enumerate}
\end{corrige}

\coursec{Formule et modèle de la molécule d'eau}

Dans la section précédente, nous avons vu que la synthèse de l'eau est possible : en faisant réagir le dihydrogène avec le dioxygène (grâce à une étincelle), on obtient de l'eau. Nous savons donc maintenant deux choses : l'électrolyse \emph{décompose} l'eau en deux gaz, et la combustion \emph{recompose} l'eau à partir de ces deux mêmes gaz. Une question naturelle se pose alors : combien d'atomes d'hydrogène et d'atomes d'oxygène se cachent dans une seule molécule d'eau ? Autrement dit, quelle est la \cle{formule} de l'eau ? C'est ce que nous allons découvrir, en suivant exactement le raisonnement des savants du XIX\textsuperscript{e} siècle.

\courssub{Des volumes de gaz à la formule : la démarche d'investigation}

\textbf{Observation : les volumes mesurés lors des expériences}\par

Reprenons nos résultats expérimentaux et regardons-les avec des yeux de chercheur.

\begin{itemize}
\item Lors de l'\cle{électrolyse} de l'eau (analyse), on recueille dans le tube au-dessus de la cathode un volume de dihydrogène \textbf{deux fois plus grand} que le volume de dioxygène recueilli au-dessus de l'anode : $V(\ce{H2}) = 2 \times V(\ce{O2})$.
\item Lors de la \cle{synthèse} de l'eau (eudiomètre), il faut mélanger \textbf{2 volumes de dihydrogène pour 1 volume de dioxygène} pour que la réaction soit complète, sans gaz restant.
\end{itemize}

Les deux expériences donnent le même rapport : \textbf{2 pour 1}. Ce n'est pas un hasard. Formulons une hypothèse : ce rapport de volumes doit correspondre au rapport des nombres d'atomes dans la molécule d'eau.

\textbf{Interprétation : des volumes aux atomes}\par

Un principe fondamental de la chimie des gaz (loi d'Avogadro, admise à ton niveau) affirme que, dans les mêmes conditions de température et de pression, \textbf{des volumes égaux de gaz contiennent le même nombre de molécules}. Le rapport des volumes est donc aussi le rapport des nombres de molécules.

Raisonnons pas à pas :
\begin{enumerate}
\item 2 volumes de \ce{H2} réagissent avec 1 volume de \ce{O2} ;
\item donc 2 molécules de \ce{H2} réagissent avec 1 molécule de \ce{O2} ;
\item chaque molécule \ce{H2} apporte 2 atomes d'hydrogène : 2 molécules \ce{H2} apportent donc $2 \times 2 = 4$ atomes H ;
\item la molécule \ce{O2} apporte 2 atomes d'oxygène ;
\item ces 4 atomes H et 2 atomes O se retrouvent dans l'eau formée : on obtient 2 molécules d'eau, chacune contenant donc $4 \div 2 = 2$ atomes H et $2 \div 2 = 1$ atome O.
\end{enumerate}

\begin{propriete}[Formule de la molécule d'eau]
La molécule d'eau, de \cle{formule brute} \ce{H2O}, est constituée de \textbf{deux atomes d'hydrogène} et d'\textbf{un atome d'oxygène} liés entre eux. L'analyse (électrolyse) et la synthèse de l'eau montrent toutes deux que l'eau résulte de la combinaison de \textbf{2 volumes de dihydrogène pour 1 volume de dioxygène}, ce qui correspond au rapport de 2 atomes H pour 1 atome O.
\end{propriete}

\begin{methode}[Déduire une formule à partir des volumes de gaz]
Pour trouver la formule d'un composé formé à partir de gaz, suis ces étapes :
\begin{enumerate}
\item Mesure ou lis le \textbf{rapport volumique expérimental} des gaz qui réagissent (ici : 2 volumes de \ce{H2} pour 1 volume de \ce{O2}).
\item Traduis ce rapport en \textbf{nombres de molécules} : le rapport des volumes est égal au rapport des nombres de molécules (2 molécules \ce{H2} pour 1 molécule \ce{O2}).
\item Compte les \textbf{atomes} apportés par chaque molécule (attention : \ce{H2} contient 2 atomes H, \ce{O2} contient 2 atomes O).
\item Répartis ces atomes dans les molécules du produit formé : tu obtiens la formule, ici \ce{H2O}.
\end{enumerate}
\textbf{Piège de l'étape 3} : ne confonds jamais le nombre de molécules avec le nombre d'atomes. « 2 volumes de \ce{H2} » signifie 2 molécules de \ce{H2}, soit 4 atomes d'hydrogène au total.
\end{methode}

\begin{popfigure}
\begin{center}
\begin{tikzpicture}[scale=1, every node/.style={font=\small}]
% Flacon H2 (2 volumes)
\draw[thick, popblue] (0,0) rectangle (1.6,2.4);
\fill[popblueL] (0,0) rectangle (1.6,2.4);
\node[popdark] at (0.8,2.7) {\textbf{2 volumes de \ce{H2}}};
\node[popdark] at (0.8,1.2) {\ce{H2} + \ce{H2}};
\node[popmuted] at (0.8,-0.4) {4 atomes H};
% Flacon O2 (1 volume)
\draw[thick, poporange] (2.6,0) rectangle (3.4,2.4);
\fill[poporangeL] (2.6,0) rectangle (3.4,2.4);
\node[popdark] at (3.0,2.7) {\textbf{1 volume de \ce{O2}}};
\node[popdark] at (3.0,1.2) {\ce{O2}};
\node[popmuted] at (3.0,-0.4) {2 atomes O};
% Flèche
\draw[-{Stealth[length=3mm]}, very thick, popdark] (4.0,1.2) -- (5.2,1.2);
\node[popdark] at (4.6,1.55) {synthèse};
% Molécules d'eau
\foreach \xc in {6.0, 8.0}{
  \fill[poporange] (\xc,1.2) circle (0.42);
  \fill[popblue] (\xc-0.42,1.62) circle (0.24);
  \fill[popblue] (\xc+0.42,1.62) circle (0.24);
  \node[white, font=\footnotesize\bfseries] at (\xc,1.2) {O};
  \node[white, font=\tiny\bfseries] at (\xc-0.42,1.62) {H};
  \node[white, font=\tiny\bfseries] at (\xc+0.42,1.62) {H};
}
\node[popdark] at (7.0,2.7) {\textbf{2 molécules d'eau}};
\node[popdark] at (7.0,-0.4) {formule de chacune : \ce{H2O}};
\end{tikzpicture}
\end{center}
\legende{Correspondance entre les volumes expérimentaux et la formule : 2 volumes de dihydrogène (4 atomes H) et 1 volume de dioxygène (2 atomes O) donnent 2 molécules d'eau. Chaque molécule d'eau contient donc 2 atomes H et 1 atome O : \ce{H2O}.}
\end{popfigure}

\courssub{Le modèle moléculaire de l'eau}

Pour « voir » une molécule, les chimistes utilisent des \cle{modèles moléculaires} : chaque atome est représenté par une boule de couleur et de taille différentes, et les boules se touchent lorsque les atomes sont liés.

\begin{definition}[Modèle moléculaire]
Le \cle{modèle moléculaire} d'un corps est une représentation de sa molécule par des boules : chaque boule figure un atome, sa couleur indique la nature de l'atome, et les boules qui se touchent représentent des atomes liés. Par convention au collège, l'hydrogène est une petite boule (souvent blanche ou bleue) et l'oxygène une grosse boule (souvent rouge ou orange).
\end{definition}

Le modèle de la molécule d'eau est donc formé de \textbf{deux petites boules H} fixées sur \textbf{une grosse boule O}. En réalité, les deux atomes d'hydrogène ne sont pas placés exactement à l'opposé l'un de l'autre : la molécule d'eau a une forme \textbf{coudée} (comme un « V »), avec un angle d'environ $105°$ entre les deux liaisons. Ce détail est un complément hors programme, mais il explique des propriétés étonnantes de l'eau.

\begin{popfigure}
\begin{center}
\begin{tikzpicture}[scale=1.3]
% Molécule d'eau coudée
\fill[poporange] (0,0) circle (0.55);
\node[white, font=\bfseries] at (0,0) {O};
\fill[popblue] (-0.62,0.62) circle (0.30);
\node[white, font=\footnotesize\bfseries] at (-0.62,0.62) {H};
\fill[popblue] (0.62,0.62) circle (0.30);
\node[white, font=\footnotesize\bfseries] at (0.62,0.62) {H};
% Angle
\draw[popmuted, dashed] (0,0) -- (-0.62,0.62);
\draw[popmuted, dashed] (0,0) -- (0.62,0.62);
\draw[popmuted] (0,0.42) arc (90:135:0.42);
\draw[popmuted] (0,0.42) arc (90:45:0.42);
\node[popmuted, font=\footnotesize] at (0,1.25) {environ $105°$};
\node[popdark, font=\small] at (0,-1.0) {Molécule d'eau \ce{H2O} (forme coudée)};
% Légende
\fill[popblue] (2.6,0.55) circle (0.16);
\node[right, font=\small] at (2.85,0.55) {atome d'hydrogène (H)};
\fill[poporange] (2.6,-0.15) circle (0.24);
\node[right, font=\small] at (2.85,-0.15) {atome d'oxygène (O)};
\end{tikzpicture}
\end{center}
\legende{Modèle moléculaire de l'eau : deux petites boules (atomes d'hydrogène H) liées à une grosse boule (atome d'oxygène O). La forme coudée (angle d'environ $105°$) est donnée à titre de complément.}
\end{popfigure}

\begin{savaistu}[Pourquoi la glace flotte-t-elle ?]
La forme coudée de la molécule d'eau explique beaucoup de ses propriétés étonnantes. À cause de cet angle, les molécules d'eau s'attirent fortement les unes les autres et, en gelant, elles s'organisent en une structure très aérée, pleine de vides. Résultat : la glace est \textbf{moins dense} que l'eau liquide, et elle flotte ! C'est pour cela que, en saison fraîche dans les pays froids, la surface des lacs gèle mais que les poissons survivent dans l'eau liquide en dessous. C'est aussi pour cela qu'une bouteille d'eau pleine peut éclater au congélateur : la glace prend plus de place que l'eau liquide. Presque tous les autres corps font l'inverse (leur solide coule dans leur liquide) : l'eau est une exception précieuse pour la vie.
\end{savaistu}

\courssub{L'eau, corps pur composé ; \ce{H2} et \ce{O2}, corps purs simples}

L'électrolyse nous a permis de décomposer l'eau : elle n'est donc pas un élément, contrairement à ce que croyaient les Anciens qui la classaient parmi les quatre « éléments » (eau, air, terre, feu).

\begin{definition}[Corps pur simple et corps pur composé]
Un \cle{corps pur simple} est un corps pur dont les molécules ne contiennent qu'\textbf{une seule sorte d'atome}. Exemples : le dihydrogène \ce{H2} (2 atomes H) et le dioxygène \ce{O2} (2 atomes O).
Un \cle{corps pur composé} est un corps pur dont les molécules contiennent \textbf{plusieurs sortes d'atomes}. Exemple : l'eau \ce{H2O} (atomes H \emph{et} O).
\end{definition}

Comme l'eau peut être décomposée en deux corps purs simples (électrolyse) et reformée à partir d'eux (synthèse), on conclut rigoureusement : \textbf{l'eau est un corps pur composé}, formé des éléments hydrogène et oxygène. Le dihydrogène et le dioxygène, eux, ne peuvent pas être décomposés en corps plus simples : ce sont des corps purs simples.

\begin{attention}[Ne confonds pas atome, molécule et élément]
\begin{itemize}
\item On dit « l'\textbf{élément} hydrogène » (symbole H) mais « le \textbf{corps pur simple} dihydrogène » (formule \ce{H2}). Le mot « hydrogène » tout seul est ambigu : précise toujours de quoi tu parles.
\item L'eau ne contient \textbf{pas} de molécules \ce{H2} ni \ce{O2} : elle contient des molécules \ce{H2O}, c'est-à-dire des atomes H et O \emph{liés}. Les propriétés de l'eau n'ont rien à voir avec celles des gaz qui la constituent : \ce{H2} brûle, \ce{O2} entretient la combustion, mais l'eau... éteint le feu !
\end{itemize}
\end{attention}

\courssub{Lecture de l'équation bilan dans les deux sens}

L'équation bilan qui résume toute cette leçon s'écrit :
\[ \ce{2H2O <=> 2H2 + O2} \]

Elle se lit dans les deux sens, ce que traduit la double flèche :

\begin{itemize}
\item \textbf{Sens direct (analyse, électrolyse)} : 2 molécules d'eau se décomposent en 2 molécules de dihydrogène et 1 molécule de dioxygène. C'est la décomposition de l'eau par le courant électrique.
\item \textbf{Sens inverse (synthèse)} : 2 molécules de dihydrogène réagissent avec 1 molécule de dioxygène pour former 2 molécules d'eau. C'est la combustion du dihydrogène dans le dioxygène.
\end{itemize}

\begin{aretenir}[Vérification de la conservation des atomes]
Dans une équation bilan équilibrée, le nombre d'atomes de chaque élément est le même des deux côtés. Vérifie pour \ce{2H2O <=> 2H2 + O2} :
\begin{itemize}
\item Hydrogène : $2 \times 2 = 4$ atomes à gauche ; $2 \times 2 = 4$ atomes à droite ;
\item Oxygène : $2 \times 1 = 2$ atomes à gauche ; $1 \times 2 = 2$ atomes à droite.
\end{itemize}
Les coefficients \textbf{2}, \textbf{2} et \textbf{1} placés devant les formules (appelés coefficients stœchiométriques) garantissent cette conservation. On ne modifie jamais les indices à l'intérieur d'une formule (le petit 2 de \ce{H2O}) pour équilibrer : on ne touche qu'aux coefficients placés devant.
\end{aretenir}

\begin{attention}[Erreurs fréquentes sur la formule de l'eau]
\begin{itemize}
\item \textbf{Écrire HO} : c'est faux. Les expériences donnent 2 volumes de \ce{H2} pour 1 volume de \ce{O2}, donc 2 atomes H pour 1 atome O : \ce{H2O}.
\item \textbf{Écrire \ce{H2O2}} : c'est la formule de l'\textbf{eau oxygénée} (peroxyde d'hydrogène), un corps \emph{différent} utilisé comme désinfectant pour nettoyer les plaies. L'eau oxygénée possède un atome d'oxygène en trop et des propriétés très différentes : ne confonds jamais \ce{H2O} et \ce{H2O2}.
\item \textbf{Écrire \ce{H2O2} pour équilibrer l'équation} : interdit ! Pour équilibrer, on place des coefficients \emph{devant} les formules (\ce{2H2O}), on ne change jamais les indices.
\end{itemize}
\end{attention}

\courssub{Complément utile au BEPC : la masse molaire de l'eau}

\emph{Ce paragraphe dépasse légèrement le strict programme de la leçon, mais il te prépare efficacement aux calculs du BEPC et de la classe de seconde.}

\begin{definition}[Masse molaire]
La \cle{masse molaire} d'un corps, notée $M$, est la masse d'une mole de ce corps ; elle s'exprime en grammes par mole (\SI{}{g/mol}). Pour une molécule, on additionne les masses molaires atomiques de tous ses atomes : $M(\ce{H}) = \SI{1}{g/mol}$ et $M(\ce{O}) = \SI{16}{g/mol}$.
\end{definition}

Pour l'eau :
\[ M(\ce{H2O}) = 2 \times M(\ce{H}) + M(\ce{O}) = 2 \times 1 + 16 = \SI{18}{g/mol} \]

Ce résultat permet un calcul très concret : dans \SI{18}{g} d'eau, il y a \SI{2}{g} d'élément hydrogène et \SI{16}{g} d'élément oxygène. En proportions massiques :
\[ \text{hydrogène : } \frac{2}{18} \approx 11\,\% \qquad ; \qquad \text{oxygène : } \frac{16}{18} \approx 89\,\% \]

Retiens bien cette surprise : l'eau contient \textbf{deux fois plus d'atomes} d'hydrogène que d'oxygène, mais presque \textbf{neuf fois plus de masse} d'oxygène, car l'atome d'oxygène est 16 fois plus lourd que l'atome d'hydrogène. Les rapports en volume (ou en nombre d'atomes) et les rapports en masse ne sont donc pas les mêmes !

\begin{exoresolu}[Proportions en masse dans l'eau]
\textbf{Énoncé.} Lors de l'électrolyse de \SI{36}{g} d'eau, quelle masse d'élément oxygène et quelle masse d'élément hydrogène obtient-on ? On donne $M(\ce{H2O}) = \SI{18}{g/mol}$, $M(\ce{H}) = \SI{1}{g/mol}$, $M(\ce{O}) = \SI{16}{g/mol}$.
\tcblower
\textbf{Solution détaillée.}

\textbf{Étape 1 : proportions massiques dans l'eau.} Dans une mole d'eau (\SI{18}{g}), l'oxygène représente \SI{16}{g} et l'hydrogène \SI{2}{g}.

\textbf{Étape 2 : application à \SI{36}{g}.} \SI{36}{g} d'eau correspondent à $36 \div 18 = 2$ moles d'eau.
\begin{align*}
m(\text{oxygène}) &= 2 \times 16 = \SI{32}{g} \\
m(\text{hydrogène}) &= 2 \times 2 = \SI{4}{g}
\end{align*}

\textbf{Étape 3 : vérification (conservation de la masse).} $32 + 4 = 36$~g : la masse totale est bien conservée, conformément à la loi de Lavoisier.

\textbf{Conclusion.} L'électrolyse de \SI{36}{g} d'eau produit \SI{32}{g} de dioxygène et \SI{4}{g} de dihydrogène. Remarque que le dioxygène est 8 fois plus lourd que le dihydrogène, alors que son \emph{volume} recueilli est 2 fois plus petit : ne confonds jamais rapport de masses et rapport de volumes.
\end{exoresolu}

\begin{exercice}[À toi de jouer]
\begin{enumerate}
\item Recopie et complète : la molécule d'eau, de formule \ldots, est constituée de \ldots atomes d'hydrogène et de \ldots atome(s) d'oxygène. L'eau est un corps pur \ldots (simple ou composé ?).
\item Lors d'une électrolyse, on recueille \SI{20}{cm^3} de dihydrogène. Quel volume de dioxygène recueille-t-on au même moment ? Justifie.
\item Pourquoi est-il faux d'écrire la formule de l'eau \ce{HO} ? Argumente à partir des volumes expérimentaux.
\item (Complément) Calcule la masse d'élément hydrogène contenue dans \SI{90}{g} d'eau.
\end{enumerate}
\end{exercice}

\begin{corrige}[Exercice « À toi de jouer »]
\begin{enumerate}
\item La molécule d'eau, de formule \ce{H2O}, est constituée de \textbf{deux} atomes d'hydrogène et d'\textbf{un} atome d'oxygène. L'eau est un corps pur \textbf{composé} (deux sortes d'atomes).
\item Le rapport volumique est $V(\ce{H2}) = 2 \times V(\ce{O2})$, donc $V(\ce{O2}) = 20 \div 2 = \SI{10}{cm^3}$.
\item L'analyse et la synthèse montrent qu'il faut \textbf{2 volumes} de dihydrogène pour \textbf{1 volume} de dioxygène, soit 2 molécules \ce{H2} (4 atomes H) pour 1 molécule \ce{O2} (2 atomes O), formant 2 molécules d'eau. Chaque molécule d'eau contient donc 2 atomes H et 1 atome O : \ce{H2O} et non \ce{HO}.
\item \SI{90}{g} d'eau représentent $90 \div 18 = 5$ moles. Masse d'hydrogène : $5 \times 2 = \SI{10}{g}$.
\end{enumerate}
\end{corrige}

\begin{aretenir}[L'essentiel de la section]
\begin{itemize}
\item L'analyse et la synthèse de l'eau montrent que l'eau résulte de la combinaison de \textbf{2 volumes de dihydrogène pour 1 volume de dioxygène}.
\item La molécule d'eau a pour formule \cle{\ce{H2O}} : 2 atomes d'hydrogène liés à 1 atome d'oxygène ; son modèle est formé de 2 petites boules H sur une grosse boule O (forme coudée en réalité).
\item L'eau est un \cle{corps pur composé} ; \ce{H2} et \ce{O2} sont des \cle{corps purs simples}.
\item L'équation \ce{2H2O <=> 2H2 + O2} se lit dans les deux sens : électrolyse (analyse) et combustion (synthèse).
\item Complément : $M(\ce{H2O}) = \SI{18}{g/mol}$ ; en masse, l'eau contient environ 11\,\% d'hydrogène et 89\,\% d'oxygène.
\end{itemize}
\end{aretenir}

\coursec{Applications et sécurité}

Nous venons de découvrir que la molécule d'eau s'écrit \ce{H2O} et qu'elle peut être \textbf{décomposée} par l'électrolyse ou \textbf{formée} par la synthèse. Ces deux transformations sont inverses l'une de l'autre : \ce{2H2O ->[electrolyse] 2H2 + O2} et \ce{2H2 + O2 ->[flammme] 2H2O}. Mais au-delà du laboratoire, ces réactions sont au cœur de grands enjeux industriels, énergétiques et de sécurité. Voyons comment la chimie que tu viens d'apprendre s'applique dans la vie réelle, ici au Cameroun et ailleurs.

\courssub{Applications industrielles de l'électrolyse}

L'électrolyse de l'eau n'est pas qu'une expérience de classe : c'est un procédé industriel majeur.

\begin{experience}[Production industrielle de dihydrogène et de dioxygène]
\textbf{Contexte.} L'industrie a besoin de grandes quantités de \ce{H2} (pour fabriquer de l'ammoniac \ce{NH3}, des engrais, raffiner le pétrole) et de \ce{O2} (pour la sidérurgie, les hôpitaux, la combustion assistée).

\textbf{Procédé.} On utilise de grands \textbf{électrolyseurs} industriels. L'eau est rendue conductrice (on ajoute de la soude \ce{NaOH}, comme au laboratoire). Les électrodes sont en nickel ou en métaux spéciaux.

\textbf{Observation.} Un courant électrique de forte intensité traverse les cuves. Le \ce{H2} est recueilli à la \textbf{cathode} (pôle $-$), le \ce{O2} à l'\textbf{anode} (pôle $+$). Les gaz sont ensuite purifiés, séchés, puis comprimés dans des bouteilles ou transportés par canalisations.

\textbf{Interprétation.} C'est exactement la même réaction qu'au laboratoire : \ce{2H2O -> 2H2 + O2}, mais réalisée en continu et optimisée pour obtenir des gaz très purs.

\textbf{Conclusion.} L'électrolyse de l'eau est la voie la plus propre pour produire du \ce{H2} dit « vert », à condition que l'électricité utilisée soit d'origine renouvelable (hydroélectricité, solaire, éolien).
\end{experience}

\begin{savaistu}[Le dihydrogène vert et le potentiel du Cameroun]
Le \textbf{dihydrogène vert} est produit par électrolyse de l'eau en utilisant \textbf{uniquement de l'électricité d'origine renouvelable}. Sa production ne rejette pas de \ce{CO2}. Le Cameroun possède un atout majeur : un \textbf{potentiel hydroélectrique énorme} (barrages de Song Loulou, Edéa, Memve'ele, Nachtigal) et un ensoleillement exceptionnel dans le Grand Nord. Demain, nos barrages et nos fermes solaires pourront alimenter de grands électrolyseurs pour produire de l'hydrogène vert : pour nos engrais, pour l'industrie, pour le transport lourd (camions, bus à pile à combustible) et même pour l'exportation. C'est une piste sérieuse pour la transition énergétique de notre pays.
\end{savaistu}

\medskip
D'autres électrolyses sont tout aussi importantes dans l'industrie :

\begin{itemize}
    \item \textbf{Électrolyse de la saumure (eau salée, \ce{NaCl} dissous) :} c'est la base de la fabrication du chlore \ce{Cl2} et de la soude \ce{NaOH}. Le mélange de ces deux produits donne l'\textbf{eau de Javel}, désinfectant utilisé dans toutes les maisons camerounaises.
    \item \textbf{Métallurgie (simple mention) :} l'électrolyse permet d'obtenir l'aluminium à partir de son minerai, et de purifier le cuivre.
\end{itemize}

\courssub{Le dihydrogène : vecteur d'énergie du futur}

Le \ce{H2} n'est pas une source d'énergie primaire (on ne le trouve pas à l'état pur dans la nature en grande quantité), mais un \textbf{vecteur d'énergie} : on stocke de l'électricité sous forme de \ce{H2} (par électrolyse) pour la restituer plus tard.

\begin{exemplebox}[La pile à combustible : l'électrolyse inversée]
Une \textbf{pile à combustible} réalise la \textbf{synthèse de l'eau} de manière contrôlée pour produire directement de l'électricité.

\textbf{Fonctionnement simplifié :}
\begin{enumerate}
    \item On alimente un côté de la pile en dihydrogène \ce{H2}, l'autre côté en dioxygène \ce{O2} (pris dans l'air).
    \item Les deux gaz réagissent sans flamme, grâce à des électrodes spéciales séparées par une membrane.
    \item La réaction \ce{2H2 + O2 -> 2H2O} libère de l'énergie, récupérée sous forme de \textbf{courant électrique} (et un peu de chaleur).
\end{enumerate}
\textbf{Avantages :} bon rendement, aucune émission polluante (le seul produit est de l'eau !), fonctionnement silencieux.

\textbf{Applications :} voitures à hydrogène (Toyota Mirai), bus, camions, groupes électrogènes de secours pour hôpitaux, drones à longue endurance.
\end{exemplebox}

\begin{savaistu}[Le chalumeau oxyhydrique : la synthèse qui soude]
Lorsque l'on brûle un mélange de \ce{H2} et \ce{O2} dans les proportions de la synthèse (2 volumes de \ce{H2} pour 1 volume de \ce{O2}), la flamme atteint environ \textbf{\SI{2800}{\celsius}}. C'est le \textbf{chalumeau oxyhydrique}. Il permet de souder ou de couper l'acier, le cuivre, le laiton avec une grande précision. Sa particularité ? La combustion ne produit \textbf{que de l'eau} (vapeur) : pas de \ce{CO2}, pas de suie. Attention : la flamme est presque invisible en plein jour (danger !) et sa température extrême impose des protections lourdes.
\end{savaistu}

\courssub{Utilisations du dioxygène}

Le \ce{O2} produit industriellement (par électrolyse ou à partir de l'air) est vital :

\begin{itemize}
    \item \textbf{Santé (hôpitaux) :} réanimation, oxygénothérapie pour les détresses respiratoires (asthme, prématurés). Les bouteilles d'oxygène médical sont indispensables dans les hôpitaux de Yaoundé, Douala et partout au Cameroun.
    \item \textbf{Soudure oxyacétylénique :} le \ce{O2} combiné à l'acétylène (\ce{C2H2}) donne une flamme à environ \SI{3200}{\celsius}, standard chez les soudeurs, serruriers et garagistes de nos quartiers.
    \item \textbf{Sidérurgie :} on souffle de l'oxygène pur dans la fonte liquide pour la transformer en acier.
    \item \textbf{Traitement de l'eau :} oxygénation des bassins d'épuration (les bactéries qui nettoient l'eau ont besoin d'oxygène).
\end{itemize}

\courssub{Consignes de sécurité au laboratoire}

L'électrolyse de l'eau produit deux gaz dangereux : \ce{H2} (très inflammable, explosif en mélange avec l'air) et \ce{O2} (comburant : il ne brûle pas, mais il accélère violemment toute combustion). La rigueur est impérative.

\begin{methode}[Fiche de sécurité : électrolyse de l'eau au collège]
\textbf{Avant de commencer :}
\begin{enumerate}
    \item Vérifier le montage \textbf{hors tension} : électrodes bien fixées, tubes à essai bien emboîtés, pas de fuite visible.
    \item S'assurer que la tension du générateur ne dépasse pas \SI{12}{V} (\SI{6}{V} ou \SI{9}{V} suffisent). \textbf{Jamais le secteur \SI{220}{V} directement.}
    \item Mettre les \textbf{lunettes de protection} (obligatoires).
    \item Aérer le local (ouvrir fenêtres et portes) : le \ce{H2} est beaucoup plus léger que l'air, il monte et s'accumule au plafond.
\end{enumerate}
\textbf{Pendant l'expérience :}
\begin{enumerate}
    \item Brancher le générateur \textbf{en dernier}.
    \item Observer la formation des bulles. Ne jamais approcher de flamme (briquet, bec Bunsen) de l'appareil en fonctionnement.
    \item Ne pas laisser l'électrolyse tourner longtemps sans surveillance.
\end{enumerate}
\textbf{Pour tester les gaz (après coupure du courant !) :}
\begin{enumerate}
    \item Couper le courant \textbf{avant} d'enlever les tubes à essai.
    \item Garder le tube renversé (ouverture vers le bas) pour le \ce{H2}, plus léger que l'air.
    \item Approcher la flamme \textbf{doucement} au bord du tube (test du \ce{H2} : « pouf » sec ; test du \ce{O2} : tison ardent qui se rallume).
    \item Tester de \textbf{petits volumes} (quelques millilitres). Jamais un tube plein de \ce{H2} : l'explosion peut casser le verre et blesser.
\end{enumerate}
\textbf{Après :} rincer le matériel, ranger les lunettes, éteindre le générateur.
\end{methode}

\begin{attention}[Pièges mortels avec le dihydrogène]
\begin{itemize}
    \item \textbf{Invisible et inodore :} on ne voit ni ne sent le \ce{H2}. Un dégagement dans un local fermé crée un mélange explosif avec l'air.
    \item \textbf{Légèreté extrême :} sa masse volumique est d'environ \SI{0,09}{g/L} (contre \SI{1,29}{g/L} pour l'air). Il s'accumule \textbf{sous les plafonds}. Une simple étincelle (interrupteur, moteur de réfrigérateur, électricité statique) suffit à déclencher l'explosion.
    \item \textbf{Jamais de flamme près de l'électrolyseur en marche.} Le « pouf » du test est une \textbf{mini-explosion contrôlée} sur un tout petit volume. Un volume plus grand provoque une déflagration violente.
    \item \textbf{Pas de bouchon sur les tubes à essai} pendant l'électrolyse : la surpression transformerait le tube en projectile de verre.
\end{itemize}
\end{attention}

\courssub{Sécurité domestique et industrielle}

Les principes du laboratoire s'appliquent à la maison et en usine.

\begin{aretenir}[Règles d'or pour les gaz inflammables (dihydrogène, butane, GPL)]
\begin{itemize}
    \item \textbf{Ventilation permanente :} toute pièce contenant une bouteille de gaz doit être aérée. Le \ce{H2} monte (aération en partie haute) ; le butane et le propane, plus lourds que l'air, descendent (aération en partie basse).
    \item \textbf{Stockage :} bouteilles à l'extérieur, à l'ombre, verticales, bien fermées. Jamais dans une cave, un garage fermé ou un placard sous l'évier.
    \item \textbf{Détection de fuite :} une odeur désagréable est ajoutée au gaz butane et au GPL pour signaler les fuites. Pour vérifier un raccord : eau savonneuse (des bulles = fuite). \textbf{Jamais de flamme pour chercher une fuite !}
    \item \textbf{Électricité :} pas d'interrupteur ni d'appareil produisant des étincelles dans une zone où un gaz inflammable peut s'accumuler.
\end{itemize}
\end{aretenir}

\courssub{Bilan récapitulatif de la leçon : analyse et synthèse}

Tout ce que nous avons fait tient en un cycle magnifique. Regarde ce schéma : il résume toute la leçon.

\begin{popfigure}
\begin{tikzpicture}[
    node distance=2.5cm and 3cm,
    box/.style={draw=popdark, thick, rounded corners=8pt, fill=popblueL, minimum width=3.5cm, minimum height=1.8cm, align=center, font=\small},
    arrow/.style={-Stealth, thick, popblue},
    testbox/.style={draw=poporange, thick, rounded corners=4pt, fill=poporangeL, minimum width=2.8cm, minimum height=1.2cm, align=center, font=\scriptsize},
    gasbox/.style={draw=popgreen, thick, rounded corners=4pt, fill=popgreenL, minimum width=2.8cm, minimum height=1.4cm, align=center, font=\small},
]
% Nœud central
\node[box, align=center] (water){\textbf{EAU LIQUIDE}\\\ce{H2O (l)}};

% Gaz produits
\node[gasbox, above left=of water, xshift=-0.5cm, align=center] (H2){\textbf{Dihydrogène}\\\ce{H2 (g)}\\(2 volumes)};
\node[gasbox, above right=of water, xshift=0.5cm, align=center] (O2){\textbf{Dioxygène}\\\ce{O2 (g)}\\(1 volume)};

% Flèches Électrolyse
\draw[arrow] (water.north west) -- node[above left, sloped, font=\scriptsize, text=popdark] {\textbf{ÉLECTROLYSE} (courant)} (H2.south east);
\draw[arrow] (water.north east) -- node[above right, sloped, font=\scriptsize, text=popdark] {\textbf{ÉLECTROLYSE} (courant)} (O2.south west);

% Anode / Cathode labels (corrects : H2 à la cathode -, O2 à l'anode +)
\node[font=\scriptsize, popblue, anchor=north west] at (H2.south west) {Cathode ($-$)};
\node[font=\scriptsize, popblue, anchor=north east] at (O2.south east) {Anode ($+$)};

% Flèches Synthèse
\draw[arrow, poporange] (H2.north east) -- node[above, font=\scriptsize, text=popdark] {\textbf{SYNTHÈSE} (flammme)} ++(2.2,0) coordinate (midtop);
\draw[arrow, poporange] (O2.north west) -- (midtop);
\draw[arrow, poporange] (midtop) -- node[right, font=\scriptsize, text=popdark] {énergie libérée} (water.north);

% Boîtes tests
\node[testbox, left=of H2, xshift=-0.3cm, align=center] (testH2){\textbf{Test \ce{H2}}\\Approcher une flamme\\$\to$ détonation « pouf »};
\node[testbox, right=of O2, xshift=0.3cm, align=center] (testO2){\textbf{Test \ce{O2}}\\Approcher un tison ardent\\$\to$ il se rallume};

% Flèches vers tests
\draw[dashed, poporange] (H2.west) -- (testH2.east);
\draw[dashed, poporange] (O2.east) -- (testO2.west);

% Équation bilan
\node[below=1.2cm of water, font=\small, align=center, text=popdark] {
    \textbf{Bilan global :} \quad \ce{2 H2O (l) ->[electrolyse] 2 H2 (g) + O2 (g)}\\
    \ce{2 H2 (g) + O2 (g) ->[flammme] 2 H2O (l)} + énergie
};

% Rappels sécurité en bas
\node[draw=poppink, thick, rounded corners=4pt, fill=poppinkL, font=\scriptsize, align=center, minimum width=3.2cm] at (-4,-7.2) {\textbf{Inflammable}\\pas de flamme};
\node[draw=poporange, thick, rounded corners=4pt, fill=poporangeL, font=\scriptsize, align=center, minimum width=3.2cm] at (0,-7.2) {\textbf{Danger électrique}\\tension $\leq$ \SI{12}{V}};
\node[draw=popblue, thick, rounded corners=4pt, fill=popblueL, font=\scriptsize, align=center, minimum width=3.2cm] at (4,-7.2) {\textbf{Lunettes}\\obligatoires};
\end{tikzpicture}
\legende{Schéma récapitulatif de la leçon : l'eau au centre ; l'électrolyse (flèches bleues) la décompose en \ce{H2} (cathode, pôle $-$) et \ce{O2} (anode, pôle $+$), identifiés par leurs tests (boîtes orange) ; la synthèse (flèches oranges) les recombine en libérant de l'énergie. En bas : les trois règles de sécurité indispensables.}
\end{popfigure}

\begin{exoresolu}[Bilan de matière lors d'une électrolyse]
\textbf{Énoncé.} On électrolyse \SI{36}{g} d'eau. On donne : $M(\ce{H}) = \SI{1}{g/mol}$, $M(\ce{O}) = \SI{16}{g/mol}$.
\begin{enumerate}
    \item Calculer la masse molaire de l'eau, puis la quantité de matière d'eau électrolysée.
    \item En utilisant l'équation bilan \ce{2H2O -> 2H2 + O2}, déterminer la quantité de matière de dihydrogène et de dioxygène produites.
    \item En déduire la masse de dihydrogène et la masse de dioxygène obtenues. Vérifier la conservation de la masse.
\end{enumerate}
\tcblower
\textbf{Solution détaillée.}
\begin{enumerate}
    \item $M(\ce{H2O}) = 2 \times 1 + 16 = \SI{18}{g/mol}$. Donc $n(\ce{H2O}) = \dfrac{36}{18} = \mathbf{2,0 \text{ mol}}$.
    \item D'après l'équation \ce{2H2O -> 2H2 + O2} : 2 moles d'eau donnent 2 moles de \ce{H2} et 1 mole de \ce{O2}. Donc $n(\ce{H2}) = \mathbf{2,0 \text{ mol}}$ et $n(\ce{O2}) = \mathbf{1,0 \text{ mol}}$.
    \item $M(\ce{H2}) = \SI{2}{g/mol}$ donc $m(\ce{H2}) = 2,0 \times 2 = \mathbf{4 \text{ g}}$. $M(\ce{O2}) = \SI{32}{g/mol}$ donc $m(\ce{O2}) = 1,0 \times 32 = \mathbf{32 \text{ g}}$. Vérification : $4 + 32 = \SI{36}{g}$, c'est bien la masse d'eau de départ. \textbf{La masse se conserve} au cours de la transformation chimique (loi de Lavoisier).
\end{enumerate}
\end{exoresolu}

\begin{exercice}[Applications et sécurité — Entraînement BEPC]
\begin{enumerate}
    \item \textbf{QCM :} Lequel de ces mélanges gazeux est le plus dangereux dans un local mal ventilé ?
    \begin{itemize}
        \item[A.] \ce{O2} pur au fond de la pièce.
        \item[B.] \ce{H2} pur accumulé sous le plafond.
        \item[C.] Air normal.
        \item[D.] \ce{CO2} au niveau du sol.
    \end{itemize}
    \item \textbf{Vrai / Faux :} Justifie chaque réponse.
    \begin{itemize}
        \item[a.] L'électrolyse de l'eau produit du dioxygène à la cathode.
        \item[b.] La pile à combustible consomme du dihydrogène et du dioxygène pour produire de l'électricité et de l'eau.
        \item[c.] Pour tester le dihydrogène, on approche une flamme d'un tube rempli aux trois quarts de ce gaz.
        \item[d.] Le chalumeau oxyhydrique brûle un mélange \ce{H2}/\ce{O2} et ne rejette que de l'eau.
    \end{itemize}
    \item \textbf{Calcul :} On souhaite produire \SI{4,0}{g} de dihydrogène par électrolyse de l'eau. On donne $M(\ce{H2}) = \SI{2}{g/mol}$, $M(\ce{H2O}) = \SI{18}{g/mol}$ et le volume molaire $V_m = \SI{22,4}{L/mol}$.
    \begin{itemize}
        \item[a.] Quelle quantité de matière de \ce{H2} cela représente-t-il ?
        \item[b.] Quelle masse d'eau faut-il électrolyser ?
        \item[c.] Quel volume de \ce{O2} est produit simultanément ?
    \end{itemize}
    \item \textbf{Rédaction :} Explique à un camarade de 4ème pourquoi on branche le générateur \textbf{en dernier} et on le coupe \textbf{en premier} lors d'une électrolyse de l'eau. Utilise les mots : \ce{H2}, explosif, étincelle, sécurité.
\end{enumerate}
\end{exercice}

\begin{corrige}[Exercice Applications et sécurité]
\begin{enumerate}
    \item \textbf{B.} Le \ce{H2} est beaucoup plus léger que l'air : il monte et s'accumule sous le plafond. Mélangé à l'air, il est explosif : une simple étincelle (interrupteur) peut tout faire exploser. L'\ce{O2} ne brûle pas seul (c'est un comburant), le \ce{CO2} n'est pas inflammable.
    \item \textbf{Vrai / Faux :}
    \begin{itemize}
        \item[a.] \textbf{Faux.} Le \ce{O2} se forme à l'\textbf{anode} (pôle $+$). C'est le \ce{H2} qui se forme à la cathode (pôle $-$).
        \item[b.] \textbf{Vrai.} C'est la transformation inverse de l'électrolyse : \ce{2H2 + O2 -> 2H2O}, avec production d'énergie électrique.
        \item[c.] \textbf{Faux, et dangereux !} On teste de \textbf{petits volumes} (quelques millilitres). Un tube presque plein de \ce{H2} exploserait violemment au contact de la flamme (bris de verre, brûlures).
        \item[d.] \textbf{Vrai.} La combustion \ce{2H2 + O2 -> 2H2O} ne produit que de l'eau, avec une flamme à environ \SI{2800}{\celsius}.
    \end{itemize}
    \item \textbf{Calcul :}
    \begin{itemize}
        \item[a.] $n(\ce{H2}) = \dfrac{4,0}{2} = \mathbf{2,0 \text{ mol}}$.
        \item[b.] D'après \ce{2H2O -> 2H2 + O2}, $n(\ce{H2O}) = n(\ce{H2}) = \SI{2,0}{mol}$. Donc $m(\ce{H2O}) = 2,0 \times 18 = \mathbf{36 \text{ g}}$.
        \item[c.] $n(\ce{O2}) = \dfrac{1}{2} n(\ce{H2}) = \SI{1,0}{mol}$. Donc $V(\ce{O2}) = 1,0 \times 22,4 = \mathbf{\SI{22,4}{L}}$.
    \end{itemize}
    \item \textbf{Rédaction (exemple) :} « On branche le générateur en dernier pour éviter qu'une étincelle au moment du branchement n'enflamme le dihydrogène \ce{H2} qui se forme dès que le courant passe, car le mélange \ce{H2}/air est explosif. On coupe le courant en premier pour arrêter la production de \ce{H2} avant de manipuler les tubes pour les tests. Ainsi, aucun gaz inflammable nouveau n'est produit pendant qu'on approche la flamme. C'est la règle de sécurité de base : courant coupé = pas de gaz nouveau = pas d'explosion. »
\end{enumerate}
\end{corrige}

\medskip
\textbf{Félicitations !} Tu viens de boucler cette leçon d'électrolyse et de synthèse de l'eau. Tu sais maintenant : décrire et interpréter l'électrolyse de l'eau, identifier les gaz \ce{H2} et \ce{O2} par leurs tests, écrire les équations d'analyse et de synthèse de l'eau, et surtout \textbf{manipuler en sécurité}. Ces bases te serviront tout au long de l'année et pour le BEPC. Révise tes fiches \texttt{À retenir}, refais les \texttt{Exercices résolus} sans regarder la correction, et tu seras prêt pour la suite ! \textbf{Courage, le BEPC approche et tu as les armes pour réussir.}

\newpage

\coursec{Activité d'intégration}

\coursec{Électrolyse et synthèse de l'eau}

\courssub{Objectifs de l'activité}

À la fin de cette activité, tu seras capable de :

\begin{itemize}
\item exploiter des mesures expérimentales d'électrolyse de l'eau et tracer des courbes $V = f(t)$ ;
\item montrer que le rapport des volumes de gaz recueillis aux deux électrodes est constant et vaut $2$ ;
\item identifier le \cle{dihydrogène} et le \cle{dioxygène} à l'aide de leurs tests caractéristiques ;
\item écrire et équilibrer les équations-bilan de l'\cle{électrolyse} (analyse) et de la \cle{synthèse} de l'eau ;
\item interpréter un mélange gazeux qui réagit et identifier un gaz résiduel ;
\item en déduire la formule brute de l'eau \ce{H2O} et son modèle moléculaire ;
\item rédiger des consignes de sécurité adaptées à la manipulation de gaz inflammables et de produits corrosifs.
\end{itemize}

\courssub{Situation déclenchante : le mystère des deux gaz}

\textbf{Enquête au laboratoire du lycée de Nkongsamba.}

Le professeur de PCT de 3\textsuperscript{e} C retrouve sur sa table le compte rendu \emph{incomplet} d'un groupe d'élèves. La veille, le groupe avait réalisé l'électrolyse d'eau rendue conductrice par quelques gouttes d'acide sulfurique dilué (eau \cle{acidulée}), à l'aide d'un électrolyseur à deux tubes gradués renversés sur les électrodes. Le tube A est relié à la borne $-$ (cathode) du générateur, le tube B à la borne $+$ (anode). Malheureusement, les élèves sont partis sans identifier les gaz, sans conclure et sans afficher les consignes de sécurité. Le professeur te confie le dossier : à toi de mener l'enquête !

\textbf{Document 1 : les mesures du groupe.} Volumes de gaz relevés toutes les \SI{2}{min} :

\begin{center}
\begin{tabularx}{\linewidth}{|l|X|X|X|X|X|X|}
\hline
\rowcolor{popblueL} Temps $t$ (\si{\minute}) & 0 & 2 & 4 & 6 & 8 & 10 \\
\hline Volume $V_A$ dans le tube A (\si{\cubic\centi\meter}) & 0 & 4 & 8 & 12 & 16 & 20 \\
\hline Volume $V_B$ dans le tube B (\si{\cubic\centi\meter}) & 0 & 2 & 4 & 6 & 8 & 10 \\
\hline
\end{tabularx}
\end{center}

\textbf{Document 2 : les tests d'identification.} Le groupe avait prélevé trois échantillons anonymes de gaz dans des flacons numérotés et noté les observations :

\begin{center}
\begin{tabularx}{\linewidth}{|l|X|X|}
\hline
\rowcolor{popblueL} Flacon & Test réalisé & Observation \\
\hline 1 & Approche d'une flamme (allumette) & Légère détonation (\og aboiement \fg) \\
\hline 2 & Bûchette incandescente (point rouge, sans flamme) & La bûchette se rallume et brille vivement \\
\hline 3 & Flamme puis bûchette incandescente & Aucun effet observable \\
\hline
\end{tabularx}
\end{center}

\textbf{Document 3 : l'expérience du mélange (eudiomètre).} Un élève rapporte : \og Nous avons introduit dans un eudiomètre \SI{30}{cm^3} du gaz du tube A et \SI{20}{cm^3} du gaz du tube B. Après le passage d'une étincelle, une détonation s'est produite, de la buée est apparue sur les parois froides, et il reste \SI{5}{cm^3} d'un seul gaz. \fg

\textbf{Problème à résoudre :} quels sont les deux gaz recueillis ? Dans quelles proportions se forment-ils et se combinent-ils ? Quelle formule et quel modèle de la molécule d'eau ces résultats permettent-ils d'établir ?

\courssub{Consignes de l'activité}

\begin{enumerate}
\item Sur papier millimétré (ou dans ton cahier), trace sur le même graphique les deux courbes $V_A = f(t)$ et $V_B = f(t)$. Que remarques-tu sur l'allure de ces courbes ?
\item Calcule le rapport $\dfrac{V_A}{V_B}$ à chaque date ($t = \SI{2}{\minute}, \SI{4}{\minute}, \SI{6}{\minute}, \SI{8}{\minute}, \SI{10}{\minute}$). Que constates-tu ? Conclus sur la proportion des volumes de gaz dégagés.
\item À l'aide du document 2, identifie le gaz de chaque flacon en justifiant chaque réponse. Attribue ensuite chaque gaz au tube A ou au tube B, en précisant l'électrode correspondante (cathode ou anode).
\item Écris et équilibre l'équation-bilan de l'électrolyse de l'eau. Pourquoi dit-on que l'électrolyse est une \emph{analyse} de l'eau ?
\item Dans le document 3, calcule les volumes de gaz qui ont réellement réagi. Vérifie qu'ils sont dans la proportion $2$ volumes de dihydrogène pour $1$ volume de dioxygène. Identifie la nature des \SI{5}{cm^3} de gaz restant et explique l'origine de la buée observée.
\item Écris et équilibre l'équation-bilan de la synthèse de l'eau. Déduis-en la formule brute de la molécule d'eau et représente son modèle de Lewis (représentation de Lewis).
\item Rédige quatre consignes de sécurité que le groupe aurait dû afficher près de l'électrolyseur.
\end{enumerate}

\courssub{Ressources : rappels de cours}

\begin{aretenir}[Ce qu'il faut avoir en tête]
\begin{itemize}
\item \textbf{Électrolyse de l'eau acidulée :} \ce{2H2O ->[electrolyse] 2H2 + O2}. Le dihydrogène se dégage à la \cle{cathode} (borne $-$), le dioxygène à l'\cle{anode} (borne $+$).
\item \textbf{Proportion des volumes :} $V(\ce{H2}) = 2 \times V(\ce{O2})$ (loi des volumes combinés de Gay-Lussac).
\item \textbf{Tests des gaz :} le dihydrogène \ce{H2} produit une \cle{détonation} (aboiement) au contact d'une flamme ; le dioxygène \ce{O2} \cle{rallume} une bûchette incandescente.
\item \textbf{Synthèse de l'eau :} \ce{2H2 + O2 ->[etincelle] 2H2O}, réaction \cle{explosive} déclenchée par une étincelle ou une flamme ; la buée est de l'eau liquide qui se condense.
\item \textbf{Conservation de la matière :} au cours d'une transformation chimique, les atomes se conservent ; il faut donc \emph{équilibrer} les équations-bilan (mêmes nombres d'atomes de chaque élément de part et d'autre de la flèche).
\end{itemize}
\end{aretenir}

\begin{propriete}[Modèle de la molécule d'eau \ce{H2O}]
La molécule d'eau est un corps composé moléculaire. Sa formule brute est \ce{H2O}.
\begin{itemize}
\item \textbf{Représentation de Lewis :} l'atome d'oxygène (6 électrons de valence) forme deux liaisons covalentes simples avec deux atomes d'hydrogène (1 électron de valence chacun). Il reste deux doublets non liants sur l'oxygène.
\item \textbf{Géométrie (modèle VSEPR) :} la répulsion entre les 4 paires d'électrons (2 liantes + 2 non liantes) donne une géométrie \textbf{coudée} (angle \ce{H-O-H} $\approx \SI{104,5}{\degree}$).
\item \textbf{Polarité :} l'oxygène étant plus électronégatif que l'hydrogène, la molécule est \textbf{polaire} (dipôle permanent), ce qui explique ses propriétés physiques exceptionnelles (solvant, point d'ébullition élevé).
\end{itemize}
\end{propriete}

\begin{savaistu}[L'hydrogène au Cameroun et dans le monde]
\begin{itemize}
\item \textbf{Stockage de l'énergie :} le dihydrogène \ce{H2} est un vecteur d'énergie propre (sa combustion ne rejette que de l'eau). On envisage de l'utiliser pour les piles à combustible (voitures, groupes électrogènes silencieux) et pour stocker l'énergie solaire ou éolienne excédentaire.
\item \textbf{Industrie :} à la \textbf{SONARA} (Limbe), l'hydrogène est utilisé pour le désulfurage des carburants (essence, gasoil) afin de réduire la pollution atmosphérique.
\item \textbf{Sécurité au quotidien :} les bonbonnes de gaz butane/propane (gaz de cuisine) sont odorisées (mercaptans) pour détecter les fuites. Le dihydrogène, lui, est inodore et incolore : d'où l'interdiction absolue de flamme près d'un électrolyseur !
\end{itemize}
\end{savaistu}

\courssub{Démarche et corrigé commenté}

\begin{exoresolu}[Le mystère des deux gaz : corrigé complet]
Résous pas à pas les sept consignes de l'enquête.
\tcblower
\textbf{Étape 1 : tracé des courbes $V = f(t)$.}
On choisit les échelles : abscisses (temps) : \SI{1}{cm} pour \SI{2}{\minute} ; ordonnées (volume) : \SI{1}{cm} pour \SI{2}{cm^3}.
Les points du tube A : $(0;0)$, $(2;4)$, $(4;8)$, $(6;12)$, $(8;16)$, $(10;20)$.
Les points du tube B : $(0;0)$, $(2;2)$, $(4;4)$, $(6;6)$, $(8;8)$, $(10;10)$.
Les points de chaque série sont alignés avec l'origine : les deux courbes sont des \textbf{droites passant par l'origine}. Le volume de gaz dégagé est donc proportionnel à la durée de l'électrolyse (loi de Faraday).

\begin{popfigure}
\begin{center}
\begin{tikzpicture}
\begin{axis}[width=0.9\linewidth,height=6.5cm,xlabel={Temps $t$ (\si{\minute})},ylabel={Volume $V$ (\si{\cubic\centi\meter})},xmin=0,xmax=11,ymin=0,ymax=22,xtick={0,2,4,6,8,10},ytick={0,4,8,12,16,20},grid=both,grid style={popline!40},legend pos=north west,legend style={font=\small,fill=white,draw=popline}]
\addplot[popcurve,color=popblue,mark=*,thick] coordinates {(0,0)(2,4)(4,8)(6,12)(8,16)(10,20)};
\addlegendentry{Tube A (cathode, borne $-$)}
\addplot[popcurve,color=poporange,mark=square*,thick] coordinates {(0,0)(2,2)(4,4)(6,6)(8,8)(10,10)};
\addlegendentry{Tube B (anode, borne $+$)}
\end{axis}
\end{tikzpicture}
\end{center}
\legende{Courbes $V=f(t)$ : droites passant par l'origine. La pente du tube A (cathode) est le double de celle du tube B (anode).}
\end{popfigure}

\textbf{Étape 2 : rapport des volumes $V_A / V_B$.}
\begin{align*}
t=\SI{2}{\minute} : \frac{V_A}{V_B} &= \frac{4}{2}=2 \quad & t=\SI{4}{\minute} : \frac{8}{4}&=2 \quad & t=\SI{6}{\minute} : \frac{12}{6}&=2 \\
t=\SI{8}{\minute} : \frac{16}{8}&=2 \quad & t=\SI{10}{\minute} : \frac{20}{10}&=2
\end{align*}
Le rapport est \textbf{constant et égal à 2} à tout instant. Le tube A (cathode) recueille toujours deux fois plus de gaz que le tube B (anode). C'est la signature volumique de l'électrolyse de l'eau : $V(\ce{H2}) = 2\,V(\ce{O2})$.

\textbf{Étape 3 : identification des gaz (Document 2).}
\begin{itemize}
\item \textbf{Flacon 1 :} détonation (\og aboiement \fg) à la flamme $\rightarrow$ test caractéristique du \textbf{dihydrogène \ce{H2}}.
\item \textbf{Flacon 2 :} bûchette incandescente qui se rallume $\rightarrow$ test caractéristique du \textbf{dioxygène \ce{O2}}.
\item \textbf{Flacon 3 :} aucun effet $\rightarrow$ ni \ce{H2} ni \ce{O2} (air résiduel ou gaz inerte).
\end{itemize}
\textbf{Attribution aux tubes :} Le dihydrogène se dégage à la cathode (borne $-$) et en volume double. C'est donc le gaz du \textbf{tube A}. Le dioxygène se dégage à l'anode (borne $+$). C'est le gaz du \textbf{tube B}.

\textbf{Étape 4 : équation-bilan de l'électrolyse (analyse de l'eau).}
\[ \ce{2H2O ->[electrolyse] 2H2 + O2} \]
\textbf{Vérification de l'équilibre :} Réactifs : 4 H, 2 O. Produits : 4 H (dans 2 \ce{H2}), 2 O (dans \ce{O2}). \textbf{Équilibrée.}
L'électrolyse \emph{décompose} l'eau en ses éléments constitutifs (dihydrogène et dioxygène) sous l'effet de l'énergie électrique : c'est une \textbf{analyse} de l'eau. Elle prouve que l'eau est un corps \emph{composé} d'hydrogène et d'oxygène.

\textbf{Étape 5 : étude du mélange dans l'eudiomètre (Document 3).}
La synthèse consomme les gaz dans la proportion stœchiométrique $2$ volumes \ce{H2} pour $1$ volume \ce{O2}.
\begin{itemize}
\item Volume de \ce{H2} introduit (tube A) : $V_{\ce{H2}} = \SI{30}{cm^3}$.
\item Volume de \ce{O2} introduit (tube B) : $V_{\ce{O2}} = \SI{20}{cm^3}$.
\item Volume de \ce{O2} nécessaire pour réagir avec tout le \ce{H2} : $V_{\ce{O2},\text{nécessaire}} = \frac{30}{2} = \SI{15}{cm^3}$.
\end{itemize}
Comme $V_{\ce{O2},\text{introduit}} (\SI{20}{cm^3}) > V_{\ce{O2},\text{nécessaire}} (\SI{15}{cm^3})$, le \textbf{dioxygène est en excès}.
Volume de \ce{O2} consommé : \SI{15}{cm^3}. Volume de \ce{O2} restant (gaz résiduel) : $20 - 15 = \SI{5}{cm^3}$. Cela correspond exactement au résidu observé.
\textbf{Nature du gaz résiduel :} \textbf{dioxygène \ce{O2}} (on pourrait le vérifier en y plongeant une bûchette incandescente : elle se rallumerait).
\textbf{Buée observée :} c'est de l'eau liquide formée par la réaction exothermique, qui se condense sur les parois froides de l'eudiomètre. Preuve que le produit est bien l'eau.

\textbf{Étape 6 : synthèse de l'eau, formule brute et modèle moléculaire.}
\[ \ce{2H2 + O2 ->[etincelle] 2H2O} \]
Deux volumes de dihydrogène réagissent avec un volume de dioxygène pour donner deux volumes de vapeur d'eau (loi des gaz parfaits, même $T$, $P$).
D'après l'équation équilibrée : 2 molécules de \ce{H2} (soit 4 atomes H) + 1 molécule de \ce{O2} (soit 2 atomes O) $\rightarrow$ 2 molécules d'eau.
Donc \textbf{1 molécule d'eau contient 2 atomes d'hydrogène et 1 atome d'oxygène} : formule brute \ce{H2O}.
\textbf{Modèle de Lewis :}
\begin{center}
\begin{tikzpicture}[scale=0.8, every node/.style={font=\small}]
% Atome O
\node[circle, draw=popdark, thick, minimum size=1.2cm, fill=popblueL] (O) at (0,0) {\textbf{O}};
% Doublets non liants
\node[circle, draw=poppurple, fill=poppurpleL, minimum size=0.6cm] (LP1) at (-0.8,0.8) {$\times\,2$};
\node[circle, draw=poppurple, fill=poppurpleL, minimum size=0.6cm] (LP2) at (0.8,-0.8) {$\times\,2$};
% Liaisons H-O
\draw[thick, popdark] (O) -- ++(-120:1.5) node[circle, draw=popgreen, thick, minimum size=0.8cm, fill=popgreenL, anchor=center] {\textbf{H}};
\draw[thick, popdark] (O) -- ++(-60:1.5) node[circle, draw=popgreen, thick, minimum size=0.8cm, fill=popgreenL, anchor=center] {\textbf{H}};
% Légende
\node[align=left, font=\small] at (3.5,0) {Représentation de Lewis :\\2 liaisons covalentes O-H\\2 doublets non liants sur O\\Géométrie : coudée ($\approx 104,5^\circ$)};
\end{tikzpicture}
\end{center}

\textbf{Étape 7 : consignes de sécurité (affichage obligatoire).}
\begin{enumerate}
\item \textbf{Interdiction formelle de toute flamme, étincelle ou cigarette} à proximité de l'électrolyseur : le dihydrogène est très inflammable et le mélange \ce{H2}/\ce{O2} (gaz de détonation) est \textbf{explosif} (risque de projection de verre et d'acide).
\item \textbf{Port des EPI obligatoires :} lunettes de protection (risque de projection d'acide ou d'éclats de verre) et blouse en coton (protection contre l'acide sulfurique dilué).
\item \textbf{Manipulation de l'eau acidulée :} verser l'acide dans l'eau (jamais l'inverse pour éviter les projections), porter des gants, ne jamais goûter, se laver les mains après manipulation.
\item \textbf{Gestion des gaz :} ne jamais boucher les tubes gradués pendant le dégagement des gaz (risque de surpression et d'éclatement) ; débrancher le générateur \textbf{avant} de retirer les tubes ou de faire les tests.
\item \textbf{Électricité et eau :} vérifier l'état des fils (pas de dénudé), ne pas poser le générateur sur le plan de travail mouillé, tenir les prises à l'écart des bacs à eau.
\end{enumerate}
\end{exoresolu}

\begin{attention}[Pièges fréquents au BEPC]
\begin{itemize}
\item \textbf{Confusion des électrodes :} le dihydrogène (le plus abondant, volume double) se dégage à la borne $-$ (\textbf{cathode}). L'anode (borne $+$) donne le dioxygène.
\item \textbf{Inversion des tests :} détonation (bruit sec) = \ce{H2} ; rallumement bûchette = \ce{O2}. Pas l'inverse !
\item \textbf{Excès de réactif :} dans la synthèse, le gaz en excès n'est pas forcément le dihydrogène. Compare toujours les volumes introduits à la proportion $2/1$ (\ce{H2}/\ce{O2}) avant de conclure.
\item \textbf{Équilibrage :} n'oublie pas les coefficients stœchiométriques devant les formules (\ce{2H2O}, \ce{2H2}, \ce{2H2O}). Les indices dans les formules (\ce{H2}, \ce{O2}, \ce{H2O}) ne se modifient \textbf{jamais}.
\item \textbf{État physique :} à température ambiante, l'eau formée est liquide (buée = condensation). L'équation-bilan s'écrit souvent \ce{2H2 + O2 -> 2H2O(l)} ou \ce{2H2O(g)} si chaud.
\end{itemize}
\end{attention}

\courssub{Bilan de l'activité}

Cette enquête t'a permis de mobiliser l'ensemble des savoir-faire de la leçon dans une situation authentique de laboratoire. En traçant les courbes $V=f(t)$, tu as établi expérimentalement que les volumes de gaz dégagés sont proportionnels au temps (loi de Faraday) et que leur rapport est constant : $V(\ce{H2}) = 2\,V(\ce{O2})$. Les tests caractéristiques (détonation pour \ce{H2}, rallumement d'une bûchette incandescente pour \ce{O2}) ont permis d'identifier chaque gaz et de l'attribuer à la bonne électrode (cathode/anode). L'étude du mélange dans l'eudiomètre a montré que la synthèse de l'eau consomme les gaz dans la même proportion $2/1$, le gaz résiduel étant le réactif en excès (ici le dioxygène), et que la buée observée est bien de l'eau liquide. Analyse (électrolyse) et synthèse concordent parfaitement : l'eau est un corps composé de formule brute \ce{H2O}, de géométrie coudée (angle \SI{104,5}{\degree}), polaire. Tu as enfin rédigé des consignes de sécurité rigoureuses, compétence indispensable dès qu'on manipule des gaz inflammables, des produits corrosifs et du matériel électrique au laboratoire.

\newpage

\coursec{Méthodes et exercices résolus}

\coursec{Électrolyse et synthèse de l'eau}

\noindent
\textbf{Situation déclenchante :} Au Cameroun, l'eau du robinet (ENEO/SNEC) ou de forage semble pure. Pourtant, si on y fait passer un courant électrique, elle se décompose en deux gaz : l'un fait « pouf » avec une flamme, l'autre ravive une braise. L'eau est-elle un corps pur ou un corps composé ? Comment produire du dihydrogène, carburant propre pour les moteurs de demain (piles à combustible, fusées) ? Nous allons décomposer l'eau par \cle{électrolyse} puis la reformer par \cle{synthèse}.

\begin{experience}[TP : Électrolyse de l'eau acidulée]
\textbf{Matériel :} générateur de courant continu (0--12~V), deux électrodes en graphite (mines de crayon) ou platine, cuve (bécher), eau distillée + quelques gouttes d'acide sulfurique dilué (ou soude), deux tubes à essai, bac pneumatique, allumettes, bûchettes, lunettes de protection.
\begin{enumerate}
\item \textbf{Montage :} Remplis la cuve d'eau acidulée. Plonge les deux électrodes. Remplis les deux tubes à essai d'eau acidulée, retourne-les au-dessus de chaque électrode (displacement d'eau). Relie les électrodes au générateur (repère le pôle $+$ et le pôle $-$).
\item \textbf{Mise en route :} Ferme le circuit. Observe la formation de bulles sur chaque électrode.
\item \textbf{Recueil :} Laisse les gaz chasser l'eau des tubes. Note le volume relatif dans chaque tube.
\item \textbf{Tests (après ouverture du circuit et retrait des tubes, ouverture vers le bas) :}
      \begin{itemize}
      \item Tube au pôle $-$ : approche une allumette enflammée $\rightarrow$ \textbf{détonation} (bruit sec).
      \item Tube au pôle $+$ : présente une bûchette incandescente (sans flamme) $\rightarrow$ elle \textbf{se rallume}.
      \end{itemize}
\end{enumerate}
\textbf{Interprétation :} Le gaz du pôle $-$ (volume double) est le \cle{dihydrogène} \ce{H2} ; celui du pôle $+$ est le \cle{dioxygène} \ce{O2}. L'eau s'est décomposée sous l'effet du courant électrique.
\end{experience}

\courssub{Méthode 1 : Identifier les gaz produits lors de l'électrolyse de l'eau}

\begin{methode}[Reconnaître le dihydrogène et le dioxygène]
Pour identifier les gaz recueillis dans les deux tubes de l'électrolyseur, suis toujours la même démarche :
\begin{enumerate}
\item \textbf{Observer les volumes} : le tube qui contient le \cle{plus grand volume} de gaz (le double de l'autre) contient le \cle{dihydrogène} \ce{H2}. C'est aussi le tube placé au-dessus de l'électrode reliée au pôle \textbf{négatif} (cathode) du générateur.
\item \textbf{Tester le dihydrogène} : approche une \textbf{allumette enflammée} de l'ouverture du tube. On entend une \textbf{détonation} (un petit « aboiement ») : c'est le test du dihydrogène.
\item \textbf{Tester le dioxygène} : présente une \textbf{bûchette (ou allumette) à peine incandescente} (rougeoyante, sans flamme) à l'ouverture de l'autre tube. Elle \textbf{se rallume} et brûle vivement : c'est le test du dioxygène \ce{O2} (tube au-dessus du pôle \textbf{positif}, anode).
\item \textbf{Conclure} en écrivant la relation des volumes : $V(\ce{H2}) = 2 \times V(\ce{O2})$.
\end{enumerate}
\begin{attention}[Réflexe sécurité]
Le mélange dihydrogène-dioxygène est \textbf{explosif}. On ne rapproche jamais une flamme de l'électrolyseur en fonctionnement ; on teste les gaz \emph{après} avoir retiré les tubes (ouverture vers le bas car \ce{H2} est plus léger que l'air).
\end{attention}
\end{methode}

\begin{exoresolu}[Identification des gaz d'une électrolyse]
Lors de l'électrolyse de l'eau additionnée de quelques gouttes d'acide, on recueille \SI{20}{cm^3} de gaz dans le tube A (au-dessus de l'électrode reliée au pôle négatif) et \SI{10}{cm^3} dans le tube B (au-dessus du pôle positif).
\begin{enumerate}
\item Quel est le gaz contenu dans chaque tube ? Justifier.
\item Décrire le test permettant d'identifier le gaz du tube A et le résultat attendu.
\end{enumerate}
\tcblower
\textbf{1. Identification des gaz.}

On compare les volumes recueillis :
\[ V_A = \SI{20}{cm^3} = 2 \times \SI{10}{cm^3} = 2 \times V_B \]
Le tube A contient un volume \textbf{double} de celui du tube B. Or, lors de l'électrolyse de l'eau, le dihydrogène se dégage à la \textbf{cathode} (pôle négatif) avec un volume double de celui du dioxygène dégagé à l'\textbf{anode} (pôle positif).
\begin{itemize}
\item Le tube A (pôle négatif, grand volume) contient le \textbf{dihydrogène} \ce{H2}.
\item Le tube B (pôle positif, petit volume) contient le \textbf{dioxygène} \ce{O2}.
\end{itemize}

\textbf{2. Test du gaz du tube A.}

On retire le tube A en le maintenant ouverture vers le bas (le dihydrogène est plus léger que l'air), puis on approche une \textbf{allumette enflammée} de son ouverture.

\textbf{Résultat attendu} : on entend une \textbf{détonation} caractéristique (petit « aboiement »).

\textbf{Conclusion} : la détonation confirme que le tube A contient bien du dihydrogène \ce{H2}.
\end{exoresolu}

\courssub{Méthode 2 : Schématiser et légender le montage de l'électrolyse de l'eau}

\begin{methode}[Dessiner l'électrolyseur sans erreur]
\begin{enumerate}
\item Dessine le \textbf{générateur de courant continu} en repérant clairement le pôle $+$ et le pôle $-$.
\item Relie chaque pôle à une \textbf{électrode} (tige de graphite ou de platine) plongée dans la cuve contenant l'\textbf{eau acidulée} (eau + quelques gouttes d'acide ou de soude pour rendre l'eau conductrice).
\item Place au-dessus de chaque électrode un \textbf{tube à essai renversé} rempli d'eau au départ, qui recueille le gaz par déplacement d'eau.
\item Indique par des flèches le \textbf{sens du courant} (du $+$ vers le $-$ à l'extérieur du générateur) et le nom de chaque gaz : \ce{H2} au pôle $-$, \ce{O2} au pôle $+$.
\item Vérifie la cohérence des niveaux de gaz : le niveau d'eau descend \textbf{deux fois plus} dans le tube du dihydrogène.
\end{enumerate}
\end{methode}

\begin{exoresolu}[Schéma de l'électrolyse]
Réaliser le schéma légendé du montage permettant de réaliser l'électrolyse de l'eau au laboratoire, en précisant la nature des gaz recueillis.
\tcblower
\textbf{Étape 1 : liste des éléments indispensables} : générateur de courant continu, fils de connexion, deux électrodes, cuve d'eau acidulée, deux tubes à essai renversés.

\textbf{Étape 2 : schéma.}

\begin{popfigure}
\begin{center}
\begin{tikzpicture}[scale=0.9]
% Cuve
\draw[thick,popblue] (-2.6,0) -- (-2.6,2.6) (2.6,0) -- (2.6,2.6) (-2.6,0) -- (2.6,0);
\fill[popblueL,opacity=0.5] (-2.55,0.05) rectangle (2.55,2.2);
\node[popdark] at (0,0.35) {\small eau acidul\'ee};
% Electrodes
\draw[thick,popdark,fill=popmuted] (-1.3,0.1) rectangle (-1.1,2.1);
\draw[thick,popdark,fill=popmuted] (1.1,0.1) rectangle (1.3,2.1);
% Tubes renverses
\draw[thick,popdark] (-1.7,2.0) -- (-1.7,4.4) -- (-0.7,4.4) -- (-0.7,2.0);
\draw[thick,popdark] (0.7,2.0) -- (0.7,4.4) -- (1.7,4.4) -- (1.7,2.0);
% Gaz H2 (grand volume)
\fill[poporangeL] (-1.65,1.2) rectangle (-0.75,4.35);
\node[popdark] at (-1.2,3.6) {\small \ce{H2}};
% Gaz O2 (petit volume)
\fill[popgreenL] (0.75,3.2) rectangle (1.65,4.35);
\node[popdark] at (1.2,3.9) {\small \ce{O2}};
% Generateur
\draw[thick,popdark,fill=popcream] (-0.5,5.6) rectangle (0.5,6.4);
\node[popdark] at (0,6.0) {\small G};
\node[popdark] at (-0.75,6.0) {\small $-$};
\node[popdark] at (0.75,6.0) {\small $+$};
% Fils
\draw[thick,popdark] (-1.2,4.4) -- (-1.2,6.0) -- (-0.5,6.0);
\draw[thick,popdark] (1.2,4.4) -- (1.2,6.0) -- (0.5,6.0);
% Annotations
\node[popdark,align=center] at (-3.3,5.2) {\small cathode\\ \small (p\^ole $-$)};
\draw[->,popdark] (-2.6,5.1) -- (-1.35,4.6);
\node[popdark,align=center] at (3.3,5.2) {\small anode\\ \small (p\^ole $+$)};
\draw[->,popdark] (2.6,5.1) -- (1.35,4.6);
\end{tikzpicture}
\end{center}
\legende{Schéma de l'électrolyse de l'eau : le dihydrogène (volume double) se recueille à la cathode, le dioxygène à l'anode.}
\end{popfigure}

\textbf{Étape 3 : vérification.} Le tube de gauche (pôle $-$) contient un volume de \ce{H2} double de celui de \ce{O2} à droite (pôle $+$) : le schéma est cohérent.

\textbf{Conclusion} : le montage comporte un générateur de courant continu, deux électrodes plongées dans l'eau acidulée et deux tubes renversés recueillant \ce{H2} au pôle $-$ et \ce{O2} au pôle $+$.
\end{exoresolu}

\courssub{Méthode 3 : Écrire et équilibrer l'équation bilan de l'électrolyse}

\begin{methode}[Équilibrer une équation chimique pas à pas]
\begin{enumerate}
\item Écris les \textbf{formules correctes} des réactifs (à gauche) et des produits (à droite) : eau \ce{H2O}, dihydrogène \ce{H2}, dioxygène \ce{O2}.
\item Compte les \textbf{atomes de chaque élément} de part et d'autre de la flèche.
\item Ajuste les \textbf{coefficients stœchiométriques} (nombres placés \emph{devant} les formules) en commençant par l'élément présent dans le moins de formules. On ne modifie \textbf{jamais} les indices à l'intérieur d'une formule.
\item Vérifie la conservation : même nombre d'atomes de chaque élément des deux côtés.
\item Précise les \textbf{conditions} au-dessus de la flèche si nécessaire (courant électrique) et l'état des espèces si demandé.
\end{enumerate}
\end{methode}

\begin{exoresolu}[Équation bilan de l'électrolyse de l'eau]
Écrire et équilibrer l'équation bilan de l'électrolyse de l'eau. Vérifier la conservation des atomes.
\tcblower
\textbf{Étape 1 : formules des réactifs et produits.}
\[ \ce{H2O -> H2 + O2} \quad \text{(non équilibrée)} \]

\textbf{Étape 2 : comptage des atomes.}
\begin{itemize}
\item À gauche : 2 atomes d'hydrogène, 1 atome d'oxygène.
\item À droite : 2 atomes d'hydrogène, 2 atomes d'oxygène.
\end{itemize}
L'oxygène n'est pas conservé : il faut équilibrer.

\textbf{Étape 3 : ajustement des coefficients.} On place le coefficient 2 devant \ce{H2O} pour avoir 2 atomes d'oxygène à gauche :
\[ \ce{2H2O -> H2 + O2} \]
Il y a alors 4 atomes d'hydrogène à gauche ; on place le coefficient 2 devant \ce{H2} :
\[ \ce{2H2O ->[courant electrique] 2H2 + O2} \]

\textbf{Étape 4 : vérification.}
\begin{itemize}
\item Hydrogène : $2 \times 2 = 4$ atomes à gauche ; $2 \times 2 = 4$ atomes à droite. \checkmark
\item Oxygène : $2 \times 1 = 2$ atomes à gauche ; $1 \times 2 = 2$ atomes à droite. \checkmark
\end{itemize}

\textbf{Conclusion} : l'équation bilan équilibrée de l'électrolyse de l'eau est $\ce{2H2O -> 2H2 + O2}$ : deux molécules d'eau sont décomposées en deux molécules de dihydrogène et une molécule de dioxygène.
\end{exoresolu}

\courssub{Méthode 4 : Exploiter les volumes de gaz (analyse quantitative)}

\begin{methode}[Utiliser la relation $V(\ce{H2}) = 2\,V(\ce{O2})$]
\begin{enumerate}
\item Repère dans l'énoncé le volume de gaz connu et le tube concerné (pôle $-$ : \ce{H2} ; pôle $+$ : \ce{O2}).
\item Applique la relation des volumes issue de l'équation bilan $\ce{2H2O -> 2H2 + O2}$ :
\[ V(\ce{H2}) = 2 \times V(\ce{O2}) \qquad \text{ou} \qquad V(\ce{O2}) = \dfrac{V(\ce{H2})}{2} \]
\item Calcule en respectant les \textbf{unités} (les deux volumes dans la même unité : \si{cm^3} ou \si{mL}, avec $1~\si{cm^3} = 1~\si{mL}$).
\item Rédige la réponse avec une phrase complète et l'unité.
\end{enumerate}
\end{methode}

\begin{exoresolu}[Calcul de volumes de gaz]
Au cours d'une électrolyse de l'eau, on recueille \SI{36}{cm^3} de dihydrogène à la cathode.
\begin{enumerate}
\item Quel volume de dioxygène recueille-t-on simultanément à l'anode ?
\item Quel volume total de gaz a été produit ?
\end{enumerate}
\tcblower
\textbf{1. Volume de dioxygène.}

D'après l'équation bilan $\ce{2H2O -> 2H2 + O2}$, le volume de dihydrogène est le double de celui du dioxygène :
\[ V(\ce{O2}) = \dfrac{V(\ce{H2})}{2} = \dfrac{36}{2} = \SI{18}{cm^3} \]

\textbf{2. Volume total de gaz.}
\[ V_{\text{total}} = V(\ce{H2}) + V(\ce{O2}) = 36 + 18 = \SI{54}{cm^3} \]

\textbf{Conclusion} : on recueille \SI{18}{cm^3} de dioxygène à l'anode, soit un volume total de gaz de \SI{54}{cm^3} (soit \SI{54}{mL}).
\end{exoresolu}

\courssub{Méthode 5 : Réaliser et interpréter la synthèse de l'eau}

\begin{experience}[TP : Synthèse de l'eau (réaction inverse)]
\textbf{Matériel :} flacon sec à large col (ou tube à essai grand modèle) avec bouchon percé, tuyau pour introduction des gaz, source de \ce{H2} (réaction Zn/HCl) et \ce{O2} (eau oxygénée/MnO2 ou bonbonne), allumette, sulfate de cuivre anhydre blanc, lunettes.
\begin{enumerate}
\item \textbf{Mélange :} Introduis \SI{20}{cm^3} de \ce{H2} puis \SI{10}{cm^3} de \ce{O2} (rapport 2/1) dans le flacon sec. Bouchonne.
\item \textbf{Amorçage :} Approche une flamme (ou étincelle) du bouchon.
\item \textbf{Observation :} \textbf{Détonation} violente. Les parois se couvrent de \textbf{buée} (gouttelettes liquides).
\item \textbf{Identification :} Verse quelques gouttes du liquide sur du sulfate de cuivre anhydre blanc $\rightarrow$ il devient \textbf{bleu}.
\end{enumerate}
\textbf{Interprétation :} Le liquide est de l'eau. La réaction est $\ce{2H2 + O2 -> 2H2O}$. L'eau est un corps composé d'hydrogène et d'oxygène.
\end{experience}

\begin{methode}[La synthèse de l'eau : protocole et conclusion]
\begin{enumerate}
\item \textbf{Mélanger} dans un flacon (ou une eudiomètre) du dihydrogène et du dioxygène, idéalement dans les proportions $2$ volumes de \ce{H2} pour $1$ volume de \ce{O2}.
\item \textbf{Enflammer} le mélange (ou provoquer une étincelle) : on observe une \textbf{détonation} violente.
\item \textbf{Observer} les parois du récipient : elles se couvrent de \textbf{buée} (fines gouttelettes liquides).
\item \textbf{Identifier} le liquide : il trouble le sulfate de cuivre anhydre blanc qui devient bleu $\Rightarrow$ c'est de l'\textbf{eau}.
\item \textbf{Conclure} : l'eau est un corps \textbf{composé} formé d'hydrogène et d'oxygène. Écris l'équation : $\ce{2H2 + O2 -> 2H2O}$, qui est la réaction \textbf{inverse} de l'électrolyse.
\end{enumerate}
\end{methode}

\begin{exoresolu}[Synthèse de l'eau au laboratoire]
On introduit dans un flacon sec \SI{20}{cm^3} de dihydrogène et \SI{10}{cm^3} de dioxygène, puis on approche une flamme. Une détonation se produit et les parois du flacon se couvrent de buée.
\begin{enumerate}
\item Quel est le liquide formé ? Comment l'identifier expérimentalement ?
\item Écrire l'équation bilan de la réaction.
\item Que peut-on conclure sur la nature chimique de l'eau ?
\end{enumerate}
\tcblower
\textbf{1. Nature et identification du liquide.}

Le liquide formé est de l'\textbf{eau}. Pour l'identifier, on verse quelques gouttes sur du \textbf{sulfate de cuivre anhydre} (poudre blanche) : celui-ci \textbf{bleuit}, ce qui prouve la présence d'eau.

\textbf{2. Équation bilan.}

Les réactifs sont le dihydrogène \ce{H2} et le dioxygène \ce{O2} ; le produit est l'eau \ce{H2O}. En équilibrant (4 atomes H et 2 atomes O de chaque côté) :
\[ \ce{2H2 + O2 -> 2H2O} \]
Les proportions de l'expérience ($20~\si{cm^3}$ de \ce{H2} pour $10~\si{cm^3}$ de \ce{O2}) correspondent exactement au rapport $2:1$ de l'équation : tout le mélange réagit.

\textbf{3. Conclusion.}

L'eau peut être \textbf{décomposée} (électrolyse) en hydrogène et oxygène, et \textbf{reformée} (synthèse) à partir de ces deux éléments : l'eau est donc un \textbf{corps composé}, constitué des éléments hydrogène et oxygène.

\textbf{Conclusion générale} : la buée est de l'eau (test au sulfate de cuivre anhydre positif) ; la réaction $\ce{2H2 + O2 -> 2H2O}$ prouve que l'eau est un corps composé.
\end{exoresolu}

\courssub{Méthode 6 : Décrire la molécule d'eau et appliquer les consignes de sécurité}

\begin{methode}[Formule et modèle de la molécule d'eau]
\begin{enumerate}
\item Retiens la \textbf{formule brute} : \ce{H2O} — une molécule d'eau contient \textbf{2 atomes d'hydrogène} et \textbf{1 atome d'oxygène}.
\item Dans le \textbf{modèle compact}, représente chaque atome par une boule : une grosse boule (oxygène) accolée à deux petites boules (hydrogène). La molécule est \textbf{coudée} (les trois centres ne sont pas alignés, angle $\approx 104,5^\circ$).
\item Pour la \textbf{sécurité} en électrolyse : utiliser une \textbf{faible tension} continue ($\le \SI{12}{V}$), ne jamais approcher de flamme pendant le dégagement des gaz, tester les gaz seulement après retrait des tubes, et manipuler l'acide (ou la soude) avec des lunettes et des gants.
\end{enumerate}
\end{methode}

\begin{exoresolu}[Modèle moléculaire et sécurité]
\begin{enumerate}
\item Donner la composition de la molécule d'eau et représenter son modèle compact.
\item Citer deux consignes de sécurité à respecter lors de l'électrolyse de l'eau, en justifiant.
\end{enumerate}
\tcblower
\textbf{1. La molécule d'eau.}

Sa formule est \ce{H2O} : elle est formée de \textbf{deux atomes d'hydrogène} liés à \textbf{un atome d'oxygène}.

\begin{popfigure}
\begin{center}
\begin{tikzpicture}
\shade[ball color=popblue] (0,0) circle (0.55);
\node[white] at (0,0) {\small O};
\shade[ball color=poporange] (-0.75,0.45) circle (0.32);
\node[white] at (-0.75,0.45) {\tiny H};
\shade[ball color=poporange] (0.75,0.45) circle (0.32);
\node[white] at (0.75,0.45) {\tiny H};
\end{tikzpicture}
\end{center}
\legende{Modèle compact de la molécule d'eau \ce{H2O} : un atome d'oxygène (grosse boule) et deux atomes d'hydrogène (petites boules), molécule coudée.}
\end{popfigure}

\textbf{2. Consignes de sécurité.}
\begin{itemize}
\item \textbf{Ne jamais approcher une flamme de l'électrolyseur en fonctionnement} : le mélange dihydrogène-dioxygène qui se dégage est \textbf{explosif} et pourrait détoner violemment.
\item \textbf{Porter des lunettes de protection} et manipuler l'acide (ou la soude) ajouté à l'eau avec précaution : ces produits sont \textbf{corrosifs}.
\end{itemize}

\textbf{Conclusion} : la molécule d'eau \ce{H2O} comporte 2 atomes d'hydrogène et 1 atome d'oxygène ; l'électrolyse exige de tenir les flammes à distance (gaz explosif) et de se protéger des produits corrosifs.
\end{exoresolu}

\begin{savaistu}[L'hydrogène, énergie de demain ?]
Au Cameroun, la \textbf{SONARA} (Société Nationale de Raffinage) utilise de l'hydrogène pour désulfurer les carburants. Demain, le \cle{dihydrogène} produit par électrolyse de l'eau (avec l'électricité solaire ou hydraulique d'ENEO) pourrait alimenter des \textbf{piles à combustible} pour des voitures, motos-taxis ou groupes électrogènes \textbf{sans émission de CO2} (seule de l'eau en sortie). C'est la filière \emph{hydrogène vert} : stocker l'énergie renouvelable sous forme chimique.
\end{savaistu}

\begin{attention}[Les 5 erreurs qui coûtent des points au BEPC]
\begin{enumerate}
\item \textbf{Inverser les électrodes} : le dihydrogène se dégage au pôle \textbf{négatif} (cathode), le dioxygène au pôle \textbf{positif} (anode). Astuce : « \ce{H2} au moins, \ce{O2} au plus ».
\item \textbf{Confondre les tests des gaz} : le dihydrogène \textbf{détone} au contact d'une \emph{flamme} ; le dioxygène \textbf{rallume} une bûchette \emph{incandescente} (sans flamme). Ne jamais écrire « le dioxygène explose ».
\item \textbf{Oublier la proportion des volumes} : $V(\ce{H2}) = 2\,V(\ce{O2})$. Écrire que les volumes sont égaux, ou inverser le rapport, fait perdre tous les points du calcul.
\item \textbf{Modifier les indices pour équilibrer} : on n'écrit jamais \ce{H2O2} pour équilibrer l'électrolyse ! On place des coefficients \emph{devant} les formules : $\ce{2H2O -> 2H2 + O2}$. Vérifie toujours le comptage final des atomes.
\item \textbf{Oublier les unités et la conclusion} : un volume sans unité (\si{cm^3} ou \si{mL}) est sanctionné, et chaque exercice doit se terminer par une phrase de conclusion répondant à la question posée.
\end{enumerate}
\end{attention}

\newpage

\coursec{Exercices}

\courssub{Je vérifie mes connaissances}

\begin{exercice}[QCM : l'électrolyse de l'eau]
Pour chaque affirmation, choisis la bonne réponse. Une seule réponse est correcte.
\begin{enumerate}
\item Pour réaliser l'électrolyse de l'eau, on ajoute quelques gouttes d'acide ou de soude à l'eau parce que :
\begin{enumerate}
\item l'eau pure conduit très mal le courant électrique ;
\item l'eau pure est trop froide ;
\item l'acide donne de la couleur au gaz.
\end{enumerate}
\item Le dihydrogène se dégage :
\begin{enumerate}
\item à l'anode (électrode reliée à la borne $+$ du générateur) ;
\item à la cathode (électrode reliée à la borne $-$ du générateur) ;
\item au fond de la cuve.
\end{enumerate}
\item Le volume de dihydrogène recueilli est :
\begin{enumerate}
\item égal au volume de dioxygène ;
\item deux fois plus petit que le volume de dioxygène ;
\item deux fois plus grand que le volume de dioxygène.
\end{enumerate}
\item L'équation de la synthèse de l'eau est :
\begin{enumerate}
\item \ce{2H2O -> 2H2 + O2} ;
\item \ce{2H2 + O2 -> 2H2O} ;
\item \ce{H2 + O -> H2O}.
\end{enumerate}
\end{enumerate}
\end{exercice}

\begin{exercice}[Vrai ou faux, avec justification]
Réponds par Vrai ou Faux, puis justifie ta réponse en une ou deux phrases.
\begin{enumerate}
\item L'eau pure conduit bien le courant électrique.
\item Le dihydrogène se dégage à l'anode lors de l'électrolyse de l'eau.
\item Le dioxygène rallume une bûchette incandescente (qui ne présente plus que des points rouges).
\item La synthèse de l'eau est une transformation endothermique.
\item L'électrolyse de l'eau est une transformation chimique.
\end{enumerate}
\end{exercice}

\begin{exercice}[Définitions à compléter]
Recopie et complète les phrases suivantes avec les mots ou expressions : \emph{analyse}, \emph{synthèse}, \emph{cathode}, \emph{anode}, \emph{électrolyse}, \emph{eudiomètre}.
\begin{enumerate}
\item On appelle \ldots\ldots\ldots\ldots la décomposition d'un corps composé par le courant électrique.
\item L'électrolyse de l'eau est une \ldots\ldots\ldots\ldots de l'eau, car on obtient les corps simples qui la constituent.
\item La réaction entre le dihydrogène et le dioxygène pour former de l'eau est une \ldots\ldots\ldots\ldots de l'eau.
\item Le dioxygène se recueille à l'\ldots\ldots\ldots\ldots et le dihydrogène à la \ldots\ldots\ldots\ldots.
\item Le tube gradué dans lequel on réalise la synthèse de l'eau s'appelle un \ldots\ldots\ldots\ldots.
\end{enumerate}
\end{exercice}

\begin{exercice}[Calcul direct de volumes]
Lors de l'électrolyse de l'eau, on recueille \SI{16}{cm^3} de dihydrogène dans le tube placé sur l'une des électrodes.
\begin{enumerate}
\item Quel volume de dioxygène recueille-t-on dans l'autre tube ? Justifie.
\item À quelle électrode (anode ou cathode) chaque gaz est-il dégagé ?
\end{enumerate}
\end{exercice}

\courssub{Je m'entraîne}

\begin{exercice}[Équilibrer et reconnaître les équations]
\begin{enumerate}
\item Équilibre les deux équations suivantes :
\[\ce{H2O -> H2 + O2} \qquad ; \qquad \ce{H2 + O2 -> H2O}\]
\item Laquelle de ces deux équations correspond à l'\emph{analyse} de l'eau ? Laquelle correspond à la \emph{synthèse} de l'eau ? Justifie tes réponses.
\item Pourquoi dit-on que ces deux transformations sont inverses l'une de l'autre ?
\end{enumerate}
\end{exercice}

\begin{exercice}[Le voltamètre à légender]
Le schéma ci-dessous représente un voltamètre utilisé pour l'électrolyse de l'eau.
\begin{center}
\begin{tikzpicture}[scale=0.9]
\draw[thick] (-3,0) -- (-3,3.5) (3,0) -- (3,3.5) (-3,0) -- (3,0);
\fill[popblueL] (-2.95,0.1) rectangle (2.95,2.2);
\draw[thick] (-2,2.2) -- (-2,4.2) -- (-0.9,4.2) -- (-0.9,2.2);
\draw[thick] (0.9,2.2) -- (0.9,4.2) -- (2,4.2) -- (2,2.2);
\fill[popblue!40] (-1.95,2.25) rectangle (-0.95,3.1);
\fill[popblue!40] (0.95,2.25) rectangle (1.95,2.65);
\draw[thick,popdark] (-1.45,0.2) -- (-1.45,2.2);
\draw[thick,popdark] (1.45,0.2) -- (1.45,2.2);
\draw[thick] (-1.45,0.2) -- (-1.45,-0.6) -- (-0.5,-0.6);
\draw[thick] (1.45,0.2) -- (1.45,-0.6) -- (0.5,-0.6);
\draw[thick] (-0.5,-0.75) -- (-0.5,-0.45);
\draw[thick] (0.5,-0.85) -- (0.5,-0.35);
\node at (-0.5,-0.3) {\scriptsize $-$};
\node at (0.5,-0.2) {\scriptsize $+$};
\node at (0,-1.3) {\scriptsize A};
\node at (-1.45,4.5) {\scriptsize B};
\node at (1.45,4.5) {\scriptsize C};
\node at (2.6,1.2) {\scriptsize D};
\end{tikzpicture}
\end{center}
\begin{enumerate}
\item Légende le schéma en identifiant : A (appareil qui fournit le courant), B et C (nom du gaz recueilli dans chaque tube), D (liquide contenu dans la cuve).
\item Donne le nom de l'électrode placée sous le tube B et celui de l'électrode placée sous le tube C.
\item Explique pourquoi le niveau du liquide descend plus vite dans le tube B que dans le tube C.
\end{enumerate}
\end{exercice}

\begin{exercice}[Mélange dans l'eudiomètre]
On introduit dans un eudiomètre \SI{50}{cm^3} de dihydrogène et \SI{50}{cm^3} de dioxygène. On provoque une étincelle électrique : la réaction a lieu et il se forme de l'eau.
\begin{enumerate}
\item Écris l'équation équilibrée de la réaction.
\item En utilisant les proportions volumiques de cette équation, montre que l'un des deux gaz est en excès.
\item Quel gaz reste-t-il dans l'eudiomètre après la réaction, et en quel volume ? Justifie ta réponse par un calcul.
\item Propose un test expérimental permettant d'identifier le gaz restant.
\end{enumerate}
\end{exercice}

\begin{exercice}[Sécurité : le test du dihydrogène]
Voici, dans le désordre, une liste de gestes proposés par des élèves pour identifier le gaz recueilli lors de l'électrolyse de l'eau. Certains gestes sont dangereux ou inutiles :
\begin{enumerate}
\item[a)] approcher une allumette enflammée de l'ouverture du tube ;
\item[b)] sentir directement le gaz avec le nez au-dessus du tube ;
\item[c)] boucher le tube avec le pouce dès la fin de l'électrolyse et le retourner ouverture vers le bas ;
\item[d)] verser de l'eau de chaux dans le tube ;
\item[e)] porter des lunettes de protection et tenir le tube à bout de bras ;
\item[f)] noter l'observation : une petite détonation (aboiement) caractérise le dihydrogène.
\end{enumerate}
\begin{enumerate}
\item Barre les gestes dangereux ou inutiles en justifiant ton choix.
\item Rédige, dans le bon ordre, le protocole correct et sécurisé du test du dihydrogène.
\end{enumerate}
\end{exercice}

\courssub{Je me prépare au BEPC}

\begin{exercice}[Type BEPC : l'électrolyse de l'eau au laboratoire]
Au laboratoire du collège, un professeur réalise l'électrolyse de l'eau additionnée de quelques gouttes de soude, à l'aide d'un voltamètre à électrodes de platine alimenté par un générateur de courant continu.
\begin{enumerate}
\item Définis le terme \emph{électrolyse}.
\item Pourquoi ajoute-t-on quelques gouttes de soude à l'eau ?
\item Fais le schéma légendé du montage (générateur, cuve, électrodes, tubes à essai renversés).
\item On observe un dégagement gazeux sur chaque électrode. Après quelques minutes, on recueille \SI{24}{cm^3} de gaz dans le tube placé au-dessus de la cathode.
\begin{enumerate}
\item Quel est ce gaz ? Décris le test qui permet de l'identifier et donne le résultat observé.
\item Calcule le volume de gaz recueilli au-dessus de l'anode. Quel est ce gaz ? Décris son test d'identification.
\end{enumerate}
\item Écris l'équation équilibrée de l'électrolyse de l'eau.
\item Pourquoi dit-on que l'électrolyse de l'eau est une \emph{analyse} de l'eau ? Que prouve-t-elle sur la nature de l'eau ?
\end{enumerate}
\end{exercice}

\begin{exercice}[Type BEPC : la synthèse de l'eau et le gaz des bonbonnes]
Dans certaines régions du Cameroun, on envisage d'utiliser le dihydrogène comme combustible pour remplacer le gaz butane des bonbonnes. Pour comprendre, un groupe d'élèves de 3ème réalise la synthèse de l'eau dans un eudiomètre contenant \SI{30}{cm^3} de dihydrogène et \SI{20}{cm^3} de dioxygène. On donne : dans les mêmes conditions de température et de pression, les volumes des gaz réagissent dans les proportions indiquées par l'équation équilibrée.
\begin{enumerate}
\item Écris l'équation équilibrée de la synthèse de l'eau.
\item Quel est le rôle de l'étincelle électrique dans cette expérience ?
\item Détermine, par le calcul, le volume de dioxygène nécessaire pour réagir exactement avec \SI{30}{cm^3} de dihydrogène.
\item En déduire le gaz qui reste dans l'eudiomètre après la réaction et son volume.
\item La synthèse de l'eau est-elle une transformation exothermique ou endothermique ? Justifie en citant une observation faite pendant l'expérience.
\item Explique pourquoi le dihydrogène est un combustible intéressant pour l'environnement, en t'appuyant sur la nature du produit de sa combustion.
\end{enumerate}
\end{exercice}

\begin{exercice}[Problème ouvert : corriger une erreur]
Lors d'une évaluation, un élève écrit : \og Faire bouillir de l'eau dans une marmite, c'est la décomposer en dihydrogène et en dioxygène, car on voit des bulles de gaz se former. \fg
\begin{enumerate}
\item Rappelle la différence entre une transformation physique et une transformation chimique.
\item À quelle catégorie appartient l'ébullition de l'eau ? Justifie.
\item Décris une expérience simple (avec le matériel du laboratoire ou de la cuisine) permettant de montrer que la vapeur d'eau est encore de l'eau.
\item Rédige un court texte argumenté (5 à 8 lignes) destiné à cet élève pour corriger son erreur, en utilisant les notions d'analyse de l'eau, de changement d'état et de molécule d'eau.
\end{enumerate}
\end{exercice}



\newpage

\coursec{Corrigés des exercices}

\begin{corrige}[Exercice 1 — QCM : l'électrolyse de l'eau]
\begin{enumerate}
\item Réponse \textbf{a} : l'eau pure est un très mauvais conducteur électrique ; l'ajout de quelques gouttes d'acide sulfurique dilué ou de solution de soude la rend conductrice en fournissant des ions.
\item Réponse \textbf{b} : le dihydrogène se dégage à la \cle{cathode}, électrode reliée à la borne $-$ du générateur.
\item Réponse \textbf{c} : le volume de dihydrogène est \textbf{deux fois plus grand} que celui de dioxygène, car l'équation \ce{2H2O -> 2H2 + O2} montre qu'il se forme 2 volumes de \ce{H2} pour 1 volume de \ce{O2} (dans les mêmes conditions de température et de pression).
\item Réponse \textbf{b} : la synthèse de l'eau s'écrit \ce{2H2 + O2 -> 2H2O}.
\end{enumerate}
\textbf{Point de vigilance} : ne pas inverser les électrodes. Astuce mnémotechnique : \og \textbf{H}ydrogène à la \textbf{c}athode, \textbf{O}xygène à l'\textbf{a}node \fg.
\end{corrige}

\begin{corrige}[Exercice 2 — Vrai ou faux, avec justification]
\begin{enumerate}
\item \textbf{Faux.} L'eau pure conduit très mal le courant électrique (c'est un isolant) ; c'est pourquoi on y ajoute quelques gouttes d'acide sulfurique dilué ou de solution de soude pour réaliser l'électrolyse.
\item \textbf{Faux.} Le dihydrogène se dégage à la \cle{cathode} (borne $-$) ; c'est le dioxygène qui se dégage à l'\cle{anode} (borne $+$).
\item \textbf{Vrai.} Le dioxygène entretient les combustions : il rallume une bûchette incandescente. C'est son test d'identification caractéristique.
\item \textbf{Faux.} La synthèse de l'eau dégage beaucoup de chaleur (flammèche, échauffement de l'eudiomètre) : c'est une transformation \cle{exothermique}.
\item \textbf{Vrai.} Il y a disparition d'un corps (l'eau) et apparition de corps nouveaux (dihydrogène et dioxygène) : c'est une transformation chimique.
\end{enumerate}
\end{corrige}

\begin{corrige}[Exercice 3 — Définitions à compléter]
\begin{enumerate}
\item On appelle \textbf{électrolyse} la décomposition d'un corps composé par le passage du courant électrique.
\item L'électrolyse de l'eau est une \textbf{analyse} de l'eau, car on obtient les corps simples qui la constituent (dihydrogène et dioxygène).
\item La réaction entre le dihydrogène et le dioxygène pour former de l'eau est une \textbf{synthèse} de l'eau.
\item Le dioxygène se recueille à l'\textbf{anode} et le dihydrogène à la \textbf{cathode}.
\item Le tube gradué dans lequel on réalise la synthèse de l'eau s'appelle un \textbf{eudiomètre}.
\end{enumerate}
\textbf{Point de vigilance} : \emph{analyse} = on décompose un corps composé en éléments ; \emph{synthèse} = on fabrique un corps composé à partir de ses éléments. Ne pas confondre les deux termes dans la rédaction.
\end{corrige}

\begin{corrige}[Exercice 4 — Calcul direct de volumes]
\begin{enumerate}
\item L'équation de l'électrolyse \ce{2H2O -> 2H2 + O2} montre que le volume de dihydrogène est le double du volume de dioxygène (loi des volumes gazeux). Donc :
\[ V(\ce{O2}) = \frac{V(\ce{H2})}{2} = \frac{\SI{16}{cm^3}}{2} = \SI{8}{cm^3}. \]
On recueille \SI{8}{cm^3} de dioxygène.
\item Le dihydrogène (\SI{16}{cm^3}) se dégage à la \cle{cathode} (borne $-$) ; le dioxygène (\SI{8}{cm^3}) se dégage à l'\cle{anode} (borne $+$).
\end{enumerate}
\textbf{Point de vigilance} : c'est le \emph{dihydrogène} qui est deux fois plus abondant ; on divise par 2 (et non on ne multiplie pas) pour trouver le volume de dioxygène.
\end{corrige}

\begin{corrige}[Exercice 5 — Équilibrer et reconnaître les équations]
\begin{enumerate}
\item Équations équilibrées :
\[\ce{2H2O -> 2H2 + O2} \qquad ; \qquad \ce{2H2 + O2 -> 2H2O}\]
Vérification : dans chaque cas, on compte 4 atomes d'hydrogène et 2 atomes d'oxygène de chaque côté de la flèche.
\item \ce{2H2O -> 2H2 + O2} est l'\textbf{analyse} de l'eau : on part d'un seul corps composé (l'eau) et on obtient les corps simples qui la constituent. \ce{2H2 + O2 -> 2H2O} est la \textbf{synthèse} de l'eau : on fabrique l'eau à partir de ses éléments.
\item Les réactifs de l'une sont les produits de l'autre : l'analyse décompose l'eau en dihydrogène et dioxygène, la synthèse reconstitue l'eau à partir de ces deux gaz. Elles sont donc inverses l'une de l'autre.
\end{enumerate}
\textbf{Point de vigilance} : pour équilibrer, on modifie uniquement les coefficients stœchiométriques placés devant les formules, jamais les indices ($\ce{H2O}$ ne devient jamais $\ce{H2O2}$).
\end{corrige}

\begin{corrige}[Exercice 6 — Le voltamètre à légender]
\begin{enumerate}
\item \textbf{A} : générateur de courant continu ; \textbf{B} : dihydrogène (tube au-dessus de l'électrode reliée à la borne $-$) ; \textbf{C} : dioxygène (tube au-dessus de l'électrode reliée à la borne $+$) ; \textbf{D} : eau additionnée de quelques gouttes d'acide sulfurique dilué ou de solution de soude (eau acidulée ou alcaline).
\item Sous le tube B se trouve la \cle{cathode} (reliée à la borne $-$) ; sous le tube C se trouve l'\cle{anode} (reliée à la borne $+$).
\item Le niveau descend plus vite dans le tube B car il s'y dégage deux fois plus de gaz : le volume de dihydrogène recueilli est le double du volume de dioxygène, d'après l'équation \ce{2H2O -> 2H2 + O2}.
\end{enumerate}

\begin{center}
\begin{tikzpicture}[scale=0.75, every node/.style={font=\small}]
% Cuve
\draw[thick, fill=popblue!10] (0,0) rectangle (6,-3.5);
% Niveau eau initial
\draw[popblue!50, line width=2pt] (0,-0.5) -- (6,-0.5);
% Générateur
\draw[thick] (3,0) -- (3,1);
\draw[thick, fill=white] (2,1) rectangle (4,1.8);
\node at (3,1.4) {G};
\node[popred] at (2.3,1.4) {$-$};
\node[popgreen] at (3.7,1.4) {$+$};
% Fils
\draw[thick, popred] (2,1) -- (2,-0.5);
\draw[thick, popgreen] (4,1) -- (4,-0.5);
% Électrodes
\draw[thick, fill=gray!50] (1.8,-0.5) rectangle (2.2,-3);
\draw[thick, fill=gray!50] (3.8,-0.5) rectangle (4.2,-3);
% Tubes renversés
\draw[thick] (1.5,-0.5) -- (1.5,2) -- (2.5,2) -- (2.5,-0.5);
\draw[thick] (3.5,-0.5) -- (3.5,2) -- (4.5,2) -- (4.5,-0.5);
% Eau résiduelle dans tubes
\fill[popblue!20] (1.5,-0.5) rectangle (2.5,0.2);
\fill[popblue!20] (3.5,-0.5) rectangle (4.5,0.2);
% Gaz collectés
\fill[poporange!40] (1.5,0.2) rectangle (2.5,2);
\fill[popgreen!40] (3.5,0.2) rectangle (4.5,1.1);
% Légendes gaz
\node[poporange, font=\bfseries] at (2,1.3) {\ce{H2}};
\node[popgreen, font=\bfseries] at (4,0.8) {\ce{O2}};
% Légendes électrodes
\node[popred, anchor=north] at (2,-3.3) {Cathode ($-$)};
\node[popgreen, anchor=north] at (4,-3.3) {Anode ($+$)};
% Légende solution
\node[anchor=north] at (3,-3.8) {Eau acidulée (ou alcaline)};
\end{tikzpicture}
\end{center}
\legende{Schéma du voltamètre à eau (électrolyse) : le tube au-dessus de la cathode (borne $-$), plus rempli, contient le dihydrogène ; l'autre contient le dioxygène.}

\textbf{Point de vigilance} : sur le schéma, le tube le plus rempli de gaz (niveau de liquide le plus bas) est toujours celui du dihydrogène, côté borne $-$.
\end{corrige}

\begin{corrige}[Exercice 7 — Mélange dans l'eudiomètre]
\begin{enumerate}
\item Équation équilibrée de la synthèse : \ce{2H2 + O2 -> 2H2O}.
\item D'après l'équation, 2 volumes de dihydrogène réagissent avec 1 volume de dioxygène. Pour brûler \SI{50}{cm^3} de dihydrogène, il faut donc $50/2 = \SI{25}{cm^3}$ de dioxygène. Or on dispose de \SI{50}{cm^3} de dioxygène : le dioxygène est en excès, le dihydrogène est le réactif limitant.
\item Il reste du dioxygène. Volume restant :
\[ V_{\text{restant}}(\ce{O2}) = \SI{50}{cm^3} - \SI{25}{cm^3} = \SI{25}{cm^3}. \]
Il reste \SI{25}{cm^3} de dioxygène (tout le dihydrogène a été consommé).
\item On approche une bûchette incandescente (présentant encore des points rouges) de l'ouverture du tube : si elle se rallume, le gaz restant est bien du dioxygène.
\end{enumerate}
\textbf{Point de vigilance} : toujours identifier d'abord le réactif limitant (ici \ce{H2}) avant de calculer le volume restant ; ne pas soustraire les volumes totaux entre eux.
\end{corrige}

\begin{corrige}[Exercice 8 — Sécurité : le test du dihydrogène]
\begin{enumerate}
\item Gestes à barrer :
\begin{itemize}
\item \textbf{b)} Sentir directement le gaz : \textbf{interdit}, certains gaz sont toxiques ou asphyxiants.
\item \textbf{d)} Verser de l'eau de chaux : c'est le test du dioxyde de carbone (\ce{CO2}), elle ne permet pas d'identifier le dihydrogène (inutile et sale).
\end{itemize}
\item Protocole correct, dans l'ordre :
\begin{enumerate}
\item \textbf{e)} Porter des lunettes de protection et tenir le tube à bout de bras.
\item \textbf{c)} Boucher le tube avec le pouce dès la fin de l'électrolyse et le retourner ouverture vers le bas (le dihydrogène, moins dense que l'air, reste ainsi piégé dans le tube).
\item \textbf{a)} Approcher une allumette enflammée de l'ouverture du tube.
\item \textbf{f)} Noter l'observation : une petite détonation sèche (aboiement) caractérise la présence de dihydrogène.
\end{enumerate}
\end{enumerate}
\textbf{Point de vigilance} : le tube se tient ouverture vers le \emph{bas} pour le dihydrogène (gaz plus léger que l'air) ; c'est l'inverse pour un gaz plus dense que l'air (ex. \ce{CO2}).
\end{corrige}

\begin{corrige}[Exercice 9 — Type BEPC : l'électrolyse de l'eau au laboratoire]
\begin{enumerate}
\item L'\cle{électrolyse} est la décomposition d'un corps composé par le passage d'un courant électrique.
\item L'eau pure conduit très mal le courant électrique (très peu d'ions \ce{H+} et \ce{HO-}) ; la soude (ou l'acide) la rend conductrice en apportant des ions mobiles, ce qui permet le passage du courant.
\item Schéma du montage (voltamètre) :
\begin{center}
\begin{tikzpicture}[scale=0.7, every node/.style={font=\small}]
% Cuve
\draw[thick, fill=popblue!10] (0,0) rectangle (6,-3.5);
% Niveau eau
\draw[popblue!50, line width=2pt] (0,-0.5) -- (6,-0.5);
% Générateur
\draw[thick] (3,0) -- (3,1);
\draw[thick, fill=white] (2,1) rectangle (4,1.8);
\node at (3,1.4) {G};
\node[popred] at (2.3,1.4) {$-$};
\node[popgreen] at (3.7,1.4) {$+$};
% Fils
\draw[thick, popred] (2,1) -- (2,-0.5);
\draw[thick, popgreen] (4,1) -- (4,-0.5);
% Électrodes
\draw[thick, fill=gray!50] (1.8,-0.5) rectangle (2.2,-3);
\draw[thick, fill=gray!50] (3.8,-0.5) rectangle (4.2,-3);
% Tubes
\draw[thick] (1.5,-0.5) -- (1.5,2) -- (2.5,2) -- (2.5,-0.5);
\draw[thick] (3.5,-0.5) -- (3.5,2) -- (4.5,2) -- (4.5,-0.5);
% Eau dans tubes
\fill[popblue!20] (1.5,-0.5) rectangle (2.5,0.2);
\fill[popblue!20] (3.5,-0.5) rectangle (4.5,0.2);
% Gaz
\fill[poporange!40] (1.5,0.2) rectangle (2.5,2);
\fill[popgreen!40] (3.5,0.2) rectangle (4.5,1.1);
% Légendes
\node[poporange, font=\bfseries] at (2,1.5) {\ce{H2}};
\node[popgreen, font=\bfseries] at (4,0.8) {\ce{O2}};
\node[popred, anchor=north] at (2,-3.3) {Cathode ($-$)};
\node[popgreen, anchor=north] at (4,-3.3) {Anode ($+$)};
\node[anchor=north] at (3,-3.8) {Eau + soude (ou acide)};
\end{tikzpicture}
\end{center}
\legende{Voltamètre : générateur courant continu (bornes $+$/$-$), cathode ($-$), anode ($+$), eau additionnée de soude, tube \ce{H2} (cathode), tube \ce{O2} (anode). Tubes renversés remplis d'eau au départ.}

\item \textbf{a)} Le gaz recueilli à la cathode est le \cle{dihydrogène}. Test : on approche une allumette enflammée de l'ouverture du tube ; on entend une petite détonation caractéristique (aboiement).
\textbf{b)} Le volume de dihydrogène étant le double de celui de dioxygène :
\[ V(\ce{O2}) = \frac{V(\ce{H2})}{2} = \frac{\SI{24}{cm^3}}{2} = \SI{12}{cm^3}. \]
Le gaz recueilli à l'anode est le \cle{dioxygène}. Test : il rallume une bûchette incandescente présentant encore des points rouges.
\item Équation équilibrée de l'électrolyse : \ce{2H2O -> 2H2 + O2}.
\item C'est une \cle{analyse} car on décompose l'eau en les corps simples qui la constituent (dihydrogène et dioxygène). Elle prouve que l'eau est un \cle{corps composé}, formé des éléments chimiques hydrogène et oxygène.
\end{enumerate}
\textbf{Barème indicatif (sur 10)} : définition (1 pt) ; rôle de la soude (1 pt) ; schéma légendé complet (2 pts) ; gaz et test à la cathode (1,5 pt) ; calcul du volume et test à l'anode (1,5 pt) ; équation équilibrée (1,5 pt) ; analyse et nature de l'eau (1,5 pt).

\textbf{Points de vigilance} : exiger la mention \og courant continu \fg{} pour le générateur ; le schéma doit montrer les tubes renversés remplis d'eau au départ ; la conclusion \og l'eau est un corps composé \fg{} est attendue explicitement à la dernière question.
\end{corrige}

\begin{corrige}[Exercice 10 — Type BEPC : la synthèse de l'eau et le gaz des bonbonnes]
\begin{enumerate}
\item Équation équilibrée de la synthèse : \ce{2H2 + O2 -> 2H2O}.
\item L'étincelle électrique fournit l'énergie d'activation nécessaire pour déclencher (amorcer) la réaction ; sans elle, le mélange \ce{H2}/\ce{O2} reste stable à température ambiante.
\item D'après l'équation, 2 volumes de \ce{H2} réagissent avec 1 volume de \ce{O2}. Donc :
\[ V(\ce{O2}) = \frac{V(\ce{H2})}{2} = \frac{\SI{30}{cm^3}}{2} = \SI{15}{cm^3}. \]
Il faut \SI{15}{cm^3} de dioxygène pour brûler exactement \SI{30}{cm^3} de dihydrogène.
\item On disposait de \SI{20}{cm^3} de dioxygène ; il en reste $20 - 15 = \SI{5}{cm^3}$. Le gaz restant est le dioxygène, en volume de \SI{5}{cm^3} (tout le dihydrogène a été consommé, c'était le réactif limitant).
\item C'est une transformation \cle{exothermique} : elle dégage de la chaleur. On observe expérimentalement une flammèche (ou une détonation) et l'échauffement des parois de l'eudiomètre ; de fines gouttelettes d'eau liquide apparaissent sur les parois froides (condensation de la vapeur d'eau formée).
\item La combustion du dihydrogène ne produit que de l'eau (\ce{2H2 + O2 -> 2H2O}), contrairement au butane (gaz des bonbonnes) qui dégage du dioxyde de carbone (\ce{CO2}), un gaz à effet de serre. Le dihydrogène ne pollue donc pas l'air lors de sa combustion : c'est un combustible propre, intéressant pour l'environnement et la transition énergétique.
\end{enumerate}
\textbf{Barème indicatif (sur 10)} : équation (1,5 pt) ; rôle de l'étincelle (1 pt) ; calcul du volume de \ce{O2} (2 pts) ; gaz restant et volume (2 pts) ; caractère exothermique justifié par observation (1,5 pt) ; argument environnemental comparatif (2 pts).

\textbf{Points de vigilance} : la justification du caractère exothermique doit citer une \emph{observation expérimentale} (flammèche, échauffement, condensation), pas seulement une affirmation ; pour la question environnementale, comparer explicitement avec la combustion du butane qui produit \ce{CO2}.
\end{corrige}

\begin{corrige}[Exercice 11 — Problème ouvert : corriger une erreur sur l'ébullition]
\begin{enumerate}
\item Dans une \cle{transformation physique}, la nature des corps ne change pas : aucun corps nouveau n'apparaît (exemple : changement d'état, dissolution). Dans une \cle{transformation chimique}, des corps disparaissent (réactifs) et des corps nouveaux apparaissent (produits).
\item L'ébullition de l'eau est une transformation \textbf{physique} : c'est un changement d'état (liquide $\to$ gaz). La vapeur formée est encore de l'eau (molécules \ce{H2O} intactes) ; aucun corps nouveau n'apparaît.
\item Expérience : on fait bouillir de l'eau dans une casserole et on tient une assiette froide (ou un verre rempli d'eau glacée) au-dessus de la vapeur. On observe la formation de gouttelettes d'eau liquide sur la surface froide : la vapeur s'est condensée. Cela prouve que la vapeur est de l'eau à l'état gazeux.
\item Texte argumenté : \og L'ébullition de l'eau n'est pas une analyse de l'eau. Quand l'eau bout, elle change simplement d'état : elle passe de l'état liquide à l'état gazeux, c'est une transformation physique. Les bulles que tu vois sont de la vapeur d'eau (\ce{H2O} gazeux), et non du dihydrogène ou du dioxygène : en refroidissant cette vapeur sur une surface froide, on retrouve de l'eau liquide. La molécule d'eau \ce{H2O} n'a pas été détruite, elle est restée identique. Pour décomposer l'eau en dihydrogène et en dioxygène (analyse de l'eau), il faut faire passer un courant électrique : c'est l'électrolyse, qui est une transformation chimique. \fg
\end{enumerate}
\textbf{Point de vigilance} : l'argument décisif est la conservation de la molécule \ce{H2O} lors de l'ébullition (changement d'état) ; citer l'électrolyse comme seul moyen (au programme de 3ème) de décomposer l'eau en ses éléments.
\end{corrige}

\newpage

\coursec{Fiche bilan}

\begin{popfigure}
\begin{tikzpicture}[mindmap, grow cyclic, every node/.style={concept}, text=white,
  root concept/.append style={concept color=popblue, font=\bfseries\small, minimum size=2.6cm},
  level 1/.append style={level distance=3.6cm, sibling angle=60, font=\scriptsize},
  level 1 concept/.append style={minimum size=2.1cm}]
\node[root concept]{Électrolyse et synthèse de l'eau}
  child[concept color=poporange]{ node {Électrolyse de l'eau : courant + eau (salgue)} }
  child[concept color=popgreen]{ node {Test de \ce{H2} : aboiement à la flamme} }
  child[concept color=popgold]{ node {Test de \ce{O2} : ravive une braise} }
  child[concept color=poppurple]{ node {Analyse : \ce{2H2O -> 2H2 + O2}} }
  child[concept color=poppink]{ node {Synthèse : \ce{2H2 + O2 -> 2H2O}} }
  child[concept color=popdark]{ node {Molécule \ce{H2O} : 2 H + 1 O} };
\end{tikzpicture}
\legende{Carte mentale de la leçon : autour de l'eau, l'électrolyse (analyse), les tests des deux gaz, la synthèse et la formule de la molécule.}
\end{popfigure}

\begin{bilan}
\textbf{Définitions clés}
\begin{itemize}
  \item \cle{Électrolyse de l'eau} : décomposition de l'eau en dihydrogène et dioxygène grâce au passage d'un \cle{courant électrique} dans de l'eau rendue conductrice (addition d'un peu de soude ou de sel).
  \item \cle{Analyse de l'eau} : opération qui décompose l'eau en ses éléments constitutifs (hydrogène et oxygène). L'électrolyse est une analyse de l'eau.
  \item \cle{Synthèse de l'eau} : opération inverse : on forme de l'eau en faisant réagir le dihydrogène et le dioxygène (mélange détonant enflammé par une étincelle).
  \item \cle{Électrodes} : conducteurs plongés dans l'eau et reliés au générateur ; la \cle{cathode} (pôle $-$) dégage \ce{H2}, l'\cle{anode} (pôle $+$) dégage \ce{O2}.
\end{itemize}

\textbf{Formules et équations à connaître par cœur}
\begin{center}
\begin{tabularx}{\linewidth}{|l|X|}
\hline
\rowcolor{popblueL} \textbf{Notion} & \textbf{À retenir} \\
\hline
Équation de l'électrolyse (analyse) & \ce{2H2O -> 2H2 + O2} \\
\hline
Équation de la synthèse de l'eau & \ce{2H2 + O2 -> 2H2O} \\
\hline
Test du dihydrogène \ce{H2} & Une flamme approchée produit une petite détonation (\emph{aboiement}) \\
\hline
Test du dioxygène \ce{O2} & Ravive la combustion d'une braise (allumette incandescente) \\
\hline
Volumes recueillis & $V(\ce{H2}) = 2 \times V(\ce{O2})$ : le dihydrogène est deux fois plus abondant \\
\hline
Formule de la molécule d'eau & \ce{H2O} : 2 atomes d'hydrogène et 1 atome d'oxygène \\
\hline
\end{tabularx}
\end{center}

\textbf{Méthodes express}
\begin{itemize}
  \item \emph{Identifier un gaz} : observer le volume (le tube le plus plein contient \ce{H2}), puis confirmer par le test à la flamme (\ce{H2}) ou à la braise (\ce{O2}).
  \item \emph{Équilibrer} : compter les atomes de chaque côté de la flèche ; ajuster les coefficients (2 devant \ce{H2O} et \ce{H2}, 1 devant \ce{O2}).
  \item \emph{Schématiser l'électrolyseur} : générateur, deux électrodes, deux tubes renversés remplis d'eau ; cathode reliée au pôle $-$ (gros volume de \ce{H2}).
\end{itemize}
\end{bilan}

\begin{aretenir}[Checklist avant le BEPC]
\begin{itemize}
  \item $\square$ Je sais définir l'électrolyse de l'eau et dire pourquoi on ajoute un peu de soude ou de sel.
  \item $\square$ Je sais dessiner et légender le schéma de l'électrolyseur (générateur, électrodes, tubes).
  \item $\square$ Je sais que \ce{H2} se dégage à la cathode (pôle $-$) et \ce{O2} à l'anode (pôle $+$).
  \item $\square$ Je sais que le volume de dihydrogène est le double de celui de dioxygène.
  \item $\square$ Je connais le test du dihydrogène : aboiement à la flamme.
  \item $\square$ Je connais le test du dioxygène : il ravive une braise.
  \item $\square$ Je sais écrire et équilibrer \ce{2H2O -> 2H2 + O2} (analyse).
  \item $\square$ Je sais écrire et équilibrer \ce{2H2 + O2 -> 2H2O} (synthèse).
  \item $\square$ Je sais distinguer analyse (décomposition) et synthèse (formation) de l'eau.
  \item $\square$ Je connais la formule \ce{H2O} et la composition de la molécule d'eau.
  \item $\square$ Je respecte les consignes de sécurité : mélange \ce{H2} + \ce{O2} détonant, petites quantités seulement.
  \item $\square$ Je sais citer une application : production de dihydrogène, soudure, piles à combustible.
\end{itemize}
\end{aretenir}

\findecours

\end{document}
